% ================================================================================
% DOCUMENTO PRINCIPAL: Inferencia Ecológica - Elecciones Uruguay 2024
% ================================================================================
% Autor: Proyecto de Análisis Electoral
% Fecha: Febrero 2025
% Descripción: Análisis comprehensivo de transferencias de votos en las
%              elecciones presidenciales uruguayas de 2024 usando inferencia
%              ecológica bayesiana
% ================================================================================

\documentclass[12pt,a4paper]{article}

% ========================================
% PACKAGES
% ========================================
% Codificación y idioma
\usepackage[utf8]{inputenc}
\usepackage[spanish]{babel}
\usepackage[margin=1in]{geometry}

% Fuentes escalables (necesario para microtype)
\usepackage{lmodern}

% Tipografía moderna (ESENCIAL - mejora calidad visual automáticamente)
\usepackage{microtype}

% Espaciado de líneas (1.5× - estándar académico)
\usepackage[onehalfspacing]{setspace}

% Matemáticas y símbolos
\usepackage{amsmath}
\usepackage{amssymb}
\usepackage{amsthm}

% Gráficos y figuras
\usepackage{graphicx}
\usepackage{float}
\usepackage[labelfont=bf]{caption}
\usepackage{subcaption}              % Moderno para subfiguras (si necesario en futuro)

% Tablas
\usepackage{booktabs}
\usepackage{multirow}
\usepackage{array}
\usepackage{siunitx}                 % Alineación decimal y formato profesional
\usepackage{adjustbox}               % Auto-ajuste de tablas anchas
\usepackage{tabularx}                % Tablas con ancho flexible
\renewcommand{\arraystretch}{1.1}    % Más espacio vertical en celdas

% Control de floats (CRÍTICO - evita que figuras/tablas se vayan al final)
\usepackage[section]{placeins}       % FloatBarrier automático en cada \section

% Optimización de parámetros de floats
\renewcommand{\topfraction}{0.85}           % 85% de página puede ser floats arriba
\renewcommand{\bottomfraction}{0.4}         % 40% abajo
\renewcommand{\textfraction}{0.15}          % Mínimo 15% de texto
\renewcommand{\floatpagefraction}{0.7}      % Página de solo floats si >70%

% Auto-centrado de todas las figuras y tablas
\makeatletter
\g@addto@macro\@floatboxreset\centering
\makeatother

% Configuración de párrafos (sangría académica estándar)
\setlength{\parindent}{15pt}
\setlength{\parskip}{0pt plus 1pt}

% Prevención de widows y orphans (líneas solas al inicio/fin de página)
\widowpenalty=10000
\clubpenalty=10000
\displaywidowpenalty=10000
\raggedbottom                        % Permite páginas cortas para evitar widows/orphans

% Citas y referencias
\usepackage{natbib}
\bibliographystyle{apalike}

% Hipervínculos (cargar al final)
\usepackage{hyperref}
\hypersetup{
    colorlinks=true,
    linkcolor=blue,
    citecolor=blue,
    urlcolor=blue,
    pdfauthor={Proyecto de Análisis Electoral},
    pdftitle={Inferencia Ecológica de las Elecciones Uruguayas 2024}
}

% ========================================
% METADATOS DEL DOCUMENTO
% ========================================
\title{%
    \textbf{Inferencia Ecológica de las Elecciones Uruguayas 2024:}\\
    \Large Análisis de Transferencias de Votos a Nivel Circuital
}
\author{Electoral Analysis Project}
\date{Febrero 2025}

% ========================================
% COMANDOS PERSONALIZADOS
% ========================================
\newcommand{\rhat}{$\hat{R}$}
\newcommand{\CI}[2]{[#1, #2]}

% Ruta a figuras y tablas
\graphicspath{{../../outputs/figures/}}
\newcommand{\tablepath}{../../outputs/tables/latex/}

% ========================================
% INICIO DEL DOCUMENTO
% ========================================
\begin{document}

% ========================================
% PRELIMINARES
% ========================================
\maketitle

\begin{abstract}
\noindent
Esta investigación emplea métodos bayesianos de inferencia ecológica para estimar las transferencias de votos entre la primera vuelta y el balotaje de las elecciones presidenciales uruguayas de 2024 y 2019. Utilizando datos de 7,271 circuitos electorales en 2024 y 7,213 en 2019, se revela una paradoja fundamental: la Coalición Republicana perdió en 2024 a pesar de mejorar dramáticamente su cohesión interna. El modelo nacional de inferencia ecológica con PI para 2024 estima que el 46.5\% de los votantes de Cabildo Abierto (CA) defectó al Frente Amplio (FA), con un 49.1\% retenido por el Partido Nacional (PN); el Partido Colorado (PC) transfirió el 10.0\% de sus votos a FA, y el PN solo el 6.3\%. La comparación temporal---realizada mediante promedios departamentales ponderados por votos para garantizar comparabilidad entre años---muestra mejoras generalizadas en cohesión coalicional entre 2019 y 2024: CA redujo su defección a FA en 41.9 puntos porcentuales, PC en 26.6 pp, PN en 32.4 pp, y PI en 92.1 pp. Sin embargo, la coalición perdió igual. El factor determinante fue el colapso del tamaño de CA, que perdió 208,722 votos (-78\%) entre elecciones, dejando a la coalición con 154 mil votos menos para retener en segunda vuelta. El análisis geográfico revela nula estabilidad temporal (correlaciones $r < 0.13$), sugiriendo reconfiguración profunda del electorado uruguayo. En 2024, los circuitos rurales exhibieron tasas de defección 13.3 puntos porcentuales más altas que los urbanos, contradiciendo supuestos convencionales sobre conservadurismo rural. La lección política clave: la cohesión en segunda vuelta no compensa la derrota en primera vuelta---se concluye que las coaliciones deben enfocarse en crecer o mantener su base electoral inicial.
\end{abstract}

\vspace{0.5cm}
\noindent
\textbf{Palabras clave:} Inferencia ecológica, elecciones Uruguay 2024, elecciones Uruguay 2019, transferencias de votos, análisis temporal, heterogeneidad geográfica, análisis bayesiano, comportamiento electoral

\newpage

% ========================================
% TABLA DE CONTENIDOS
% ========================================
\tableofcontents
\newpage

% ========================================
% SECCIONES PRINCIPALES
% ========================================

% Incluir sección de introducción
% ========================================
% INTRODUCTION SECTION
% Ecological Inference Analysis of Uruguay 2024 Elections
% ========================================

\section{Introducción}
\label{sec:introduccion}

\subsection{Contexto político}

El 27 de octubre de 2024, Uruguay celebró la primera vuelta de sus elecciones presidenciales en un escenario político marcado por la experiencia de gobierno de la Coalición Republicana. Desde 2020, esta alianza---integrada por el Partido Nacional (PN), el Partido Colorado (PC), Cabildo Abierto (CA) y el Partido Independiente (PI)---había gobernado el país bajo la presidencia de Luis Lacalle Pou, poniendo fin a quince años de administraciones del Frente Amplio (FA). La primera vuelta de 2024 reveló una coalición fragmentada: mientras el PN obtuvo el 26.8\% de los votos, el PC alcanzó el 16.1\%, CA apenas superó el 2.7\% y el PI registró un 1.8\%. El FA, por su parte, se posicionó como primera fuerza con el 44.1\% de los sufragios, insuficientes para alcanzar la mayoría absoluta requerida para evitar un balotaje.

El 24 de noviembre de 2024, en la segunda vuelta, el candidato del FA, Yamandú Orsi, derrotó al candidato del PN, Álvaro Delgado, por un margen de aproximadamente 95.000 votos. Este resultado planteó una pregunta central para la ciencia política uruguaya: ¿hacia dónde se dirigieron los votos de los partidos de la coalición que no llegaron al balotaje? Más específicamente, ¿en qué medida los votantes de CA, PC, PI y los partidos menores transfirieron su apoyo al candidato de la coalición, y cuántos optaron por respaldar al FA o emitir voto en blanco?

\subsection{El problema de la inferencia ecológica}

Las transferencias de votos entre rondas electorales constituyen un fenómeno de interés fundamental que, sin embargo, no puede observarse directamente. El secreto del voto impide conocer las decisiones individuales de cada elector: solo se dispone de resultados agregados a nivel de circuito electoral. Se sabe cuántos votos obtuvo cada partido en cada circuito en primera vuelta y cuántos sufragios recibió cada candidato en el balotaje, pero la matriz de transición subyacente---es decir, la proporción de votantes de cada partido de origen que se dirigió a cada opción de destino---permanece latente.

Este desafío inferencial se conoce en la literatura estadística como el \textit{problema de la inferencia ecológica} \citep{king1997}. La \textit{falacia ecológica} surge cuando se infiere comportamiento individual directamente a partir de correlaciones observadas a nivel agregado, sin considerar que la heterogeneidad dentro de las unidades puede generar patrones agregados engañosos. \citet{king1997} propuso una solución bayesiana al problema, desarrollando un modelo que estima probabilidades de transición individuales a partir de datos agregados respetando las restricciones naturales del problema (probabilidades acotadas entre 0 y 1 que suman 1 para cada partido de origen). La extensión al caso multidimensional $R \times C$ fue desarrollada por \citet{rosen2001bayesian}, permitiendo estimar matrices de transición completas en sistemas multipartidistas como el uruguayo.

\subsection{Relevancia del estudio}

El análisis de las transferencias de votos en el balotaje uruguayo de 2024 reviste particular relevancia por múltiples razones. En primer lugar, la elección se decidió por un margen relativamente estrecho en el contexto de una segunda vuelta, lo que hace que cada flujo de transferencia resulte potencialmente decisivo para explicar el resultado. En segundo lugar, el estudio permite evaluar la cohesión electoral de la Coalición Republicana, una alianza heterogénea que agrupa partidos con tradiciones ideológicas diversas, desde el centro-derecha liberal del PC hasta el populismo de derecha de CA. La capacidad de esta coalición para retener sus votos en segunda vuelta constituye una prueba empírica de su viabilidad como proyecto político.

En tercer lugar, el caso de Cabildo Abierto merece atención especial. Fundado en 2019 por el general retirado Guido Manini Ríos, CA irrumpió en el sistema de partidos uruguayo obteniendo el 11\% de los votos en su primera elección, para luego colapsar al 2.7\% en 2024---una caída del 78\% en su caudal electoral. Comprender hacia dónde se dirigieron los votantes que abandonaron CA entre 2019 y 2024, y cómo se comportaron los votantes remanentes de CA en el balotaje, resulta esencial para evaluar la estabilidad del sistema de partidos uruguayo, uno de los más institucionalizados de América Latina \citep{mainwaring1995party}.

Finalmente, el análisis permite testear hipótesis sobre la geografía del comportamiento electoral. La literatura sobre divisiones urbano-rurales en democracias contemporáneas \citep{rodden2019s} sugiere que los votantes rurales tienden a ser más conservadores y leales a partidos de derecha. El caso uruguayo ofrece una oportunidad para evaluar si este patrón se reproduce en el contexto de transferencias entre rondas, o si, por el contrario, la dinámica de segunda vuelta genera patrones diferentes a los esperados.

\subsection{Antecedentes en la literatura}

El estudio del comportamiento electoral uruguayo cuenta con una sólida tradición académica. \citet{gonzalez1991electoral} y \citet{gonzalez1993structuring} establecieron el marco analítico fundamental para comprender las estructuras políticas y la transformación democrática en Uruguay, documentando la estabilidad del sistema bipartidista tradicional y su posterior evolución hacia un sistema tripartidista con la consolidación del FA. \citet{mainwaring1995party} ubicó a Uruguay como uno de los sistemas de partidos más institucionalizados de América Latina, con altos niveles de estabilidad electoral, legitimidad de los procesos electorales y arraigo organizacional de los partidos. Más recientemente, \citet{bottinelli2021elecciones} analizó las elecciones de 2019, documentando las continuidades y cambios en el sistema de partidos tras la irrupción de CA y la formación de la Coalición Republicana.

En el plano metodológico, la inferencia ecológica ha sido aplicada extensamente al análisis de transferencias de votos en sistemas multipartidistas. El trabajo fundacional de \citet{king1997} propuso el modelo jerárquico binomial-beta para el caso $2 \times 2$, mientras que \citet{rosen2001bayesian} extendieron el enfoque al caso general $R \times C$ mediante modelos bayesianos con distribuciones Dirichlet, precisamente la formulación requerida para el sistema multipartidista uruguayo. Las aplicaciones de inferencia ecológica en contextos electorales latinoamericanos han sido menos frecuentes que en la tradición anglosajona, lo que confiere al presente estudio un carácter exploratorio relevante para la disciplina en la región.

\subsection{Contribuciones}

El presente estudio realiza las siguientes contribuciones:

\begin{enumerate}
    \item \textbf{Primera aplicación sistemática de inferencia ecológica bayesiana a balotajes uruguayos a nivel circuital.} Se analizan 7.271 circuitos electorales en 2024 y 7.118 en 2019, alcanzando la máxima resolución geográfica disponible en datos oficiales. Esta granularidad permite capturar heterogeneidad local que queda enmascarada en análisis departamentales o nacionales.

    \item \textbf{Comparación temporal 2019--2024 que revela cambios estructurales en la dinámica de coaliciones.} El análisis de dos elecciones consecutivas permite evaluar la estabilidad de los patrones de transferencia y distinguir entre comportamientos electorales coyunturales y tendencias estructurales del sistema de partidos.

    \item \textbf{Estratificación geográfica que identifica un patrón de ``rebelión rural''.} La desagregación urbano/rural y metropolitano/interior revela que los circuitos rurales exhibieron tasas de defección de la coalición sustancialmente superiores a las de los circuitos urbanos, contradiciendo supuestos convencionales sobre el conservadurismo político rural.

    \item \textbf{Análisis de sensibilidad y robustez que valida la metodología.} Se evalúa la estabilidad de las estimaciones frente a variaciones en distribuciones previas, tratamiento de valores atípicos y remuestreo bootstrap, proporcionando evidencia de la confiabilidad de los resultados.

    \item \textbf{Análisis de los patrones de voto en blanco como hipótesis alternativa.} Se examina si el voto en blanco actuó como refugio para electores descontentos de la coalición, o si constituye un fenómeno independiente de las transferencias partidarias.
\end{enumerate}

\subsection{Estructura del artículo}

El resto de este artículo se organiza como sigue. La Sección~\ref{sec:metodologia} describe la metodología empleada: el modelo de inferencia ecológica de King, la implementación bayesiana con PyMC, los datos electorales utilizados y la estrategia analítica. La Sección~\ref{sec:resultados} presenta los resultados del análisis de 2024, incluyendo estimaciones nacionales, estratificaciones urbano/rural y metropolitano/interior, y variación departamental. La Sección~\ref{sec:resultados_temporal} examina la evolución temporal comparando los patrones de transferencia de 2019 y 2024. La Sección~\ref{sec:sensibilidad} evalúa la sensibilidad de las estimaciones a variaciones metodológicas. La Sección~\ref{sec:blank_votes} analiza los patrones de voto en blanco. Finalmente, la Sección~\ref{sec:discusion} discute las implicaciones de los hallazgos para la comprensión del sistema de partidos uruguayo y presenta las conclusiones del estudio.

\newpage

% Incluir sección de metodología
% ========================================
% METHODOLOGY SECTION
% Ecological Inference Analysis of Uruguay 2024 Elections
% ========================================

\section{Metodología}
\label{sec:metodologia}

\subsection{El Problema de la Inferencia Ecológica}

La \textit{ecological inference} se refiere al desafío estadístico fundamental de inferir comportamiento a nivel individual a partir de datos agregados \cite{king1997}. En estudios electorales, frecuentemente se observan totales de votos a nivel de unidades geográficas (circuitos, distritos, departamentos) pero se carece de información directa sobre las transiciones de voto individuales. La pregunta central se convierte en: ¿cómo distribuyeron sus votos en la segunda elección los votantes que apoyaron al partido $i$ en la primera elección?

La \textit{ecological fallacy} surge cuando se hacen inferencias directas sobre comportamiento individual a partir de patrones agregados sin considerar la heterogeneidad dentro de las unidades \cite{king1997}. Por ejemplo, observar que un área geográfica con alto apoyo al partido $i$ en octubre también muestra alto apoyo al partido $j$ en noviembre no implica necesariamente que los votantes se trasladaron de $i$ a $j$—el patrón agregado podría resultar de dinámicas individuales completamente diferentes.

En el contexto de las elecciones presidenciales de Uruguay 2024, se busca estimar las transferencias de votos entre la primera vuelta (27 de octubre de 2024) y el balotaje (24 de noviembre de 2024). Se observan conteos de votos agregados a nivel de circuito pero no se puede observar directamente cómo los votantes individuales cambiaron sus preferencias. Este desafío inferencial requiere métodos estadísticos sofisticados que puedan extraer información a nivel individual mientras respetan los límites impuestos por el proceso generador de datos.

\subsection{El Método de Inferencia Ecológica de King}

El método de Inferencia Ecológica de King (King's Ecological Inference, EI) \cite{king1997} proporciona una solución fundamentada a este problema inferencial. El método estima una \textit{matriz de transición} $\mathbf{T}$ de dimensión $K \times J$, donde $K$ representa el número de partidos de origen (primera vuelta) y $J$ representa el número de partidos de destino (balotaje). Cada elemento $T_{kj}$ de esta matriz representa la probabilidad de que un votante que apoyó al partido $k$ en la primera elección votara por el partido $j$ en la segunda elección.

Formalmente, para cada circuito electoral $i$, sea $\mathbf{X}_i = (X_{i1}, \ldots, X_{iK})$ el vector de votos recibidos por cada uno de los $K$ partidos de origen, y sea $\mathbf{Y}_i = (Y_{i1}, \ldots, Y_{iJ})$ los votos recibidos por cada uno de los $J$ partidos de destino. Las transferencias de votos esperadas se modelan como:

\begin{equation}
\label{eq:transferencias_esperadas}
Y_{ij} = \sum_{k=1}^{K} T_{kj} \cdot X_{ik}
\end{equation}

En notación matricial, esto puede expresarse como $\mathbf{Y}_i = \mathbf{X}_i \mathbf{T}$, indicando que los votos de destino observados son una combinación lineal de los votos de origen ponderados por las probabilidades de transición.

El método de King ofrece ventajas críticas sobre enfoques clásicos como la regresión ecológica de Goodman \cite{king1997}. Más importante aún, respeta los límites naturales sobre las probabilidades de transición: $0 \leq T_{kj} \leq 1$ para todo $k,j$, y $\sum_{j=1}^{J} T_{kj} = 1$ para cada partido de origen $k$. Estas restricciones aseguran que las probabilidades de transición estimadas sean interpretables como probabilidades propias y que todos los votantes de cada partido de origen sean contabilizados en la elección de destino. La regresión de Goodman, en contraste, frecuentemente produce estimaciones fuera del intervalo $[0,1]$, generando resultados sustancialmente carentes de sentido.

\subsection{Implementación Bayesiana}

Se implementa el método EI de King usando un modelo jerárquico bayesiano con el framework de programación probabilística PyMC \cite{abril2023pymc}. El enfoque bayesiano proporciona varias ventajas para la inferencia ecológica: (1) incorpora naturalmente conocimiento previo a través de distribuciones de probabilidad, (2) cuantifica la incertidumbre mediante distribuciones posteriores en lugar de estimaciones puntuales, (3) respeta restricciones mediante elecciones apropiadas de distribuciones previas, y (4) permite la propagación coherente de la incertidumbre desde los parámetros hasta las cantidades de interés.

El núcleo del modelo especifica una distribución previa sobre la matriz de transición que impone las restricciones requeridas. Específicamente, para cada partido de origen $k$, la fila $\mathbf{T}_k = (T_{k1}, \ldots, T_{kJ})$ sigue una distribución Dirichlet:

\begin{equation}
\label{eq:prior_dirichlet}
\mathbf{T}_k \sim \text{Dirichlet}(\boldsymbol{\alpha})
\end{equation}

donde $\boldsymbol{\alpha} = (\alpha_1, \ldots, \alpha_J)$ es un vector de parámetros de concentración. La distribución Dirichlet es la elección natural para modelar vectores de probabilidad: garantiza que todos los elementos se encuentren en $[0,1]$ y sumen 1, precisamente las restricciones necesarias para las probabilidades de transición \cite{gelman2013bayesian}. Se utiliza una distribución previa uniforme con $\alpha_j = 1$ para todo $j$, lo que expresa información previa mínima y permite que los datos dominen la distribución posterior.

Dada la matriz de transición, se modelan los votos de destino observados usando una aproximación Normal a la distribución Multinomial. Para conteos de votos grandes, la aproximación Normal proporciona eficiencia computacional mientras mantiene precisión estadística \cite{martin2021bayesian}. Para el circuito $i$ y el partido de destino $j$, se especifica:

\begin{equation}
\label{eq:verosimilitud}
Y_{ij} \sim \mathcal{N}\left(\mu_{ij}, \sigma_{ij}^2\right)
\end{equation}

donde la media $\mu_{ij} = \sum_{k=1}^{K} T_{kj} \cdot X_{ik}$ captura los votos esperados basados en la ecuación \ref{eq:transferencias_esperadas}, y la desviación estándar $\sigma_{ij} = \sqrt{\mu_{ij} + 1}$ da cuenta de la variabilidad de muestreo con una pequeña constante para evitar inestabilidad numérica.

La distribución posterior combina el conocimiento previo con la verosimilitud de los datos observados a través del teorema de Bayes. Esta distribución posterior sobre la matriz de transición $\mathbf{T}$ representa la creencia actualizada sobre los patrones de transferencia de votos después de observar los datos electorales. Se resume esta distribución usando medias posteriores (estimaciones puntuales) e intervalos de credibilidad (cuantificación de incertidumbre).

\subsection{Muestreo MCMC y Diagnósticos de Convergencia}

La inferencia posterior se realiza usando muestreo de Monte Carlo mediante Cadenas de Markov (Markov Chain Monte Carlo, MCMC), específicamente el No-U-Turn Sampler (NUTS) \cite{betancourt2017conceptual}, una variante de Monte Carlo Hamiltoniano (Hamiltonian Monte Carlo, HMC). NUTS es particularmente adecuado para modelos bayesianos complejos porque explora eficientemente espacios de parámetros de alta dimensión y se adapta automáticamente a la geometría de la distribución posterior.

Se emplean diferentes configuraciones de muestreo dependiendo del alcance del análisis:

\begin{itemize}
    \item \textbf{Análisis nacional}: 4 cadenas, 4000 muestras post-calentamiento por cadena, 2000 muestras de calentamiento, tasa de aceptación objetivo 0.95
    \item \textbf{Análisis estratificados} (urbano/rural, regional): 4 cadenas, 2000 muestras post-calentamiento por cadena, 2000 muestras de calentamiento, tasa de aceptación objetivo 0.90
\end{itemize}

Estas elecciones de parámetros reflejan un equilibrio cuidadoso entre eficiencia computacional y rigurosidad estadística. El análisis nacional usa más muestras y una tasa de aceptación objetivo más alta porque representa el objetivo inferencial primario y requiere la mayor precisión. Los análisis estratificados usan parámetros ligeramente relajados mientras mantienen estándares científicos, dado que estos análisis involucran tamaños de muestra más pequeños y sirven como exploraciones secundarias de heterogeneidad.

La tasa de aceptación objetivo más alta en el análisis nacional (0.95) reduce el tamaño de paso en HMC, llevando a una exploración más cuidadosa de la distribución posterior a costa de un tiempo de cómputo incrementado. Este enfoque conservador minimiza el riesgo de muestras sesgadas y asegura una mezcla exhaustiva a través de la distribución posterior. El período de calentamiento de 2000 muestras permite al muestreador adaptar sus parámetros (tamaño de paso y matriz de masa) a la geometría específica de cada distribución posterior.

Ejecutar 4 cadenas en paralelo sirve dos propósitos: (1) acelera el cómputo mediante paralelización, y (2) permite verificaciones diagnósticas de convergencia comparando la variabilidad dentro de las cadenas y entre cadenas. Múltiples cadenas iniciadas desde valores iniciales dispersos deberían converger a la misma distribución si el muestreador está funcionando correctamente.

Se evalúa la convergencia usando dos diagnósticos primarios:

\begin{enumerate}
    \item \textbf{Estadístico de Gelman-Rubin ($\hat{R}$)} \cite{gelman1992inference}: Este diagnóstico compara la varianza dentro de las cadenas con la varianza entre cadenas. Valores de $\hat{R} < 1.01$ indican que las cadenas han convergido a una distribución común. Se requiere que todos los parámetros satisfagan este criterio.

    \item \textbf{Tamaño de Muestra Efectivo (Effective Sample Size, ESS)}: Debido a la autocorrelación en las muestras MCMC, el número efectivo de extracciones independientes es típicamente menor que el número total de muestras. Se requiere ESS $> 1000$ para todos los parámetros para asegurar suficiente precisión en las estimaciones posteriores e intervalos de credibilidad \cite{gelman2013bayesian}.
\end{enumerate}

Solo los análisis que satisfacen ambos criterios de convergencia ($\hat{R} < 1.01$ y ESS $> 1000$ para todos los parámetros) se incluyen en los resultados reportados. Esta adherencia estricta a los diagnósticos de convergencia asegura la validez científica de las inferencias resultantes.

\subsection{Fuentes de Datos y Procesamiento}

\subsubsection{Datos Electorales 2024}

Los datos electorales oficiales fueron obtenidos de la Corte Electoral de Uruguay, el órgano constitucional responsable de organizar y supervisar las elecciones en Uruguay. Para 2024, se analizan dos elecciones:

\begin{itemize}
    \item \textbf{Primera vuelta}: 27 de octubre de 2024
    \item \textbf{Ballotage} (segunda vuelta): 24 de noviembre de 2024
\end{itemize}

Los datos electorales están organizados a nivel de circuito (\textit{circuito}), la unidad geográfica más pequeña para la cual se reportan resultados oficiales. El conjunto de datos 2024 abarca 7,271 circuitos electorales a través de los 19 departamentos de Uruguay. Cada circuito típicamente atiende entre 100 y 1,000 votantes, proporcionando resolución geográfica de grano fino para el análisis de inferencia ecológica.

Para la primera vuelta 2024, se clasifican los votos en seis categorías (separando el Partido Independiente de la categoría "OTROS" dada su relevancia política):
\begin{itemize}
    \item \textbf{FA}: Frente Amplio (coalición de izquierda)
    \item \textbf{PN}: Partido Nacional (centro-derecha)
    \item \textbf{PC}: Partido Colorado (centro-derecha)
    \item \textbf{CA}: Cabildo Abierto (derecha populista)
    \item \textbf{PI}: Partido Independiente (centro-derecha, 1.79\% de votos)
    \item \textbf{OTROS}: Partidos pequeños remanentes agregados
\end{itemize}

\subsubsection{Datos Electorales 2019}

Para el análisis temporal comparativo, se obtuvieron datos electorales de 2019 también de la Corte Electoral:

\begin{itemize}
    \item \textbf{Primera vuelta}: 27 de octubre de 2019
    \item \textbf{Ballotage}: 24 de noviembre de 2019
\end{itemize}

El conjunto de datos 2019 abarca 7,213 circuitos electorales, ligeramente menor al de 2024 debido a cambios en la organización electoral. Se utiliza idéntica categorización de partidos para asegurar comparabilidad temporal. Se define la \textbf{Coalición Republicana} como la suma de PN + PC + CA + PI en ambos años. Notablemente, en 2019 el Partido Independiente obtuvo solo 1.01\% de los votos (23,554 votos), mientras que en 2024 creció a 1.79\% (41,055 votos).

La disponibilidad de datos estructuralmente idénticos en ambas elecciones (primera vuelta y balotaje, misma resolución circuital, mismos partidos principales) permite análisis temporal riguroso de cambios en patrones de transferencia de votos y estabilidad geográfica de comportamientos electorales.

Para el balotaje, solo dos partidos principales compiten, complementados por votos en blanco y anulados:
\begin{itemize}
    \item \textbf{FA}: Frente Amplio
    \item \textbf{PN}: Partido Nacional (representando a la coalición)
    \item \textbf{Blancos/Nulos}: Votos en blanco y anulados
\end{itemize}

\subsubsection{Covariables Sociodemográficas}

Para permitir análisis estratificados, se incorporan datos sociodemográficos de dos fuentes:

\begin{itemize}
    \item \textbf{Censo 2023} (Instituto Nacional de Estadística): Recuentos poblacionales, clasificación urbano/rural
    \item \textbf{Datos administrativos electorales} (Corte Electoral): Clasificaciones geográficas a nivel de circuito
\end{itemize}

La clasificación urbano/rural divide los circuitos en dos estratos basados en la densidad poblacional y los patrones de asentamiento. Los circuitos urbanos (86.8\% del total) están ubicados en ciudades y pueblos, mientras que los circuitos rurales (13.2\%) sirven a poblaciones rurales dispersas. Esta estratificación es crítica porque investigaciones previas sobre la política uruguaya sugieren diferencias sustanciales en el comportamiento político entre áreas urbanas y rurales \cite{manacorda2011transfers}.

Adicionalmente, se estratifica por macrorregión:
\begin{itemize}
    \item \textbf{Área Metropolitana}: Montevideo y Canelones (conteniendo aproximadamente el 60\% de la población nacional)
    \item \textbf{Interior}: Todos los otros 17 departamentos
\end{itemize}

\subsubsection{Calidad y Validación de Datos}

Se implementan varios procedimientos de control de calidad:

\begin{enumerate}
    \item \textbf{Verificaciones de completitud}: Verificar que los totales de votos coincidan con los agregados oficiales a nivel departamental y nacional
    \item \textbf{Verificaciones de consistencia}: Asegurar que la suma de votos por partido sea igual al total de votos válidos en cada circuito
    \item \textbf{Análisis de datos faltantes}: Examinar patrones de datos faltantes y aplicar eliminación por lista cuando sea necesario (menos del 0.1\% de los circuitos tienen valores faltantes)
    \item \textbf{Detección de valores atípicos}: Identificar circuitos con patrones anómalos que puedan indicar errores en los datos (ninguno detectado)
\end{enumerate}

\subsection{Estrategia Analítica}

La estrategia analítica procede en tres etapas, avanzando desde patrones nacionales hacia la heterogeneidad geográfica:

\subsubsection{Análisis Nacional}

El análisis nacional agrupa los 7,271 circuitos para estimar patrones generales de transferencia de votos entre la primera vuelta y el balotaje. Esto proporciona la matriz de transición base que representa el comportamiento promedio en Uruguay. Las estimaciones nacionales sirven como puntos de referencia para evaluar la heterogeneidad en los análisis estratificados.

\subsubsection{Análisis Estratificados}

Se realizan dos conjuntos de análisis estratificados para explorar la heterogeneidad geográfica:

\begin{enumerate}
    \item \textbf{Urbano vs Rural}: Se estiman matrices de transición separadas para circuitos urbanos (n = 6,312) y circuitos rurales (n = 959). Esta estratificación aborda la hipótesis de que los votantes rurales exhiben patrones diferentes de lealtad partidaria y defección comparados con los votantes urbanos.

    \item \textbf{Área Metropolitana vs Interior}: Se compara la región capital (Montevideo y Canelones) con el resto del país. Esta estratificación captura potenciales dinámicas centro-periferia en el comportamiento electoral.
\end{enumerate}

Cada análisis estratificado produce matrices de transición específicas del estrato con intervalos de credibilidad asociados. Se comparan formalmente los estratos examinando si los intervalos de credibilidad se superponen para las probabilidades de transición clave.

\subsubsection{Análisis a Nivel Departamental}

Para proporcionar perspectivas geográficas detalladas, se estiman matrices de transición separadamente para cada uno de los 19 departamentos de Uruguay. Estos análisis a nivel departamental utilizan tamaños de muestra menores (que van desde aproximadamente 100 hasta 1,500 circuitos por departamento) y consecuentemente exhiben mayor incertidumbre. Sin embargo, revelan importantes variaciones locales en los patrones de transferencia de votos que quedan enmascaradas en los agregados nacionales.

Los resultados a nivel departamental son particularmente valiosos para comprender las dinámicas políticas regionales e identificar áreas donde los partidos nacionales experimentaron ganancias o pérdidas inusuales. Se presentan estos resultados principalmente a través de visualización (mapas y gráficos de bosque) en lugar de tablas detalladas debido al gran número de estimaciones.

\subsection{Validación del Modelo}

Para evaluar la validez de las estimaciones de inferencia ecológica se realizaron varios ejercicios de validación:

\subsubsection{Verificaciones Predictivas Posteriores}

Para cada modelo ajustado, se generan distribuciones predictivas posteriores muestreando de la matriz de transición estimada y simulando recuentos de votos para cada circuito. Luego se compara la distribución de votos simulados con los votos observados. Un buen ajuste del modelo se indica por una superposición sustancial entre las distribuciones predichas y observadas \cite{gelman2013bayesian}.

\subsubsection{Consistencia de Agregación}

Una verificación clave de validez para la inferencia ecológica es la consistencia de agregación: cuando agregamos las predicciones a nivel de circuito a niveles superiores (departamento, región, nación), ¿las predicciones agregadas coinciden con los datos agregados observados? Se verificó que las probabilidades de transición a nivel de circuito reproducen los resultados departamentales y nacionales observados cuando se ponderan por los totales de votos de origen.

\subsubsection{Análisis de Sensibilidad}

Para evaluar la robustez, se realizan análisis de sensibilidad variando:
\begin{itemize}
    \item Especificaciones de priors (priors uniformes vs priors Dirichlet débilmente informativos)
    \item Especificaciones de verosimilitud (Normal vs Multinomial)
    \item Parámetros MCMC (número de muestras, cadenas, tasa de aceptación)
\end{itemize}

Los resultados demuestran ser robustos a través de estas variaciones, con medias posteriores que difieren en menos de 0.5 puntos porcentuales e intervalos de credibilidad que muestran superposición sustancial.

\subsubsection{Comparación con Métodos Alternativos}

Aunque la regresión ecológica de Goodman frecuentemente produce estimaciones fuera de los límites (una falla fundamental), se estima con fines de comparación. Como se esperaba, el método de Goodman produce varias probabilidades de transición fuera de $[0,1]$, confirmando la necesidad del enfoque acotado de King. Donde las estimaciones de Goodman están dentro de los límites, se alinean estrechamente con el EI de King, lo que proporciona confianza adicional en la validez de los resultados.

\subsection{Software y Reproducibilidad}

Todos los análisis se realizaron en Python 3.11 utilizando PyMC 5.10 \cite{abril2023pymc} para inferencia bayesiana, NumPy 1.26 y Pandas 2.1 para manipulación de datos, y ArviZ 0.17 para diagnósticos de convergencia y análisis posterior. La visualización utilizó Matplotlib 3.8 y Seaborn 0.13. El pipeline analítico completo, incluyendo procesamiento de datos, estimación de modelos, diagnósticos y visualización, está disponible en un repositorio público para asegurar la reproducibilidad completa de los resultados.

Todos los generadores de números aleatorios se inicializaron con la semilla 42 para permitir la replicación exacta del muestreo MCMC. El muestreador NUTS utilizó la adaptación de matriz de masa predeterminada de PyMC y el ajuste de tamaño de paso durante la fase de warmup.

\subsection{Consideraciones Éticas}

Este análisis utiliza únicamente datos electorales agregados y disponibles públicamente publicados por la Corte Electoral de Uruguay. No se recolectaron ni analizaron datos a nivel individual. Todos los resultados se presentan a nivel agregado (circuito o superior), preservando el anonimato de los votantes. La metodología de inferencia ecológica reconoce explícitamente que las estimaciones a nivel individual son inferencias probabilísticas de datos agregados, no observaciones determinísticas del comportamiento individual.

\newpage

% Incluir sección de resultados
% ========================================
% RESULTS SECTION
% Ecological Inference Analysis of Uruguay 2024 Elections
% ========================================

\section{Resultados}
\label{sec:resultados}

Esta sección presenta los hallazgos del análisis de inferencia ecológica sobre las transferencias de votos entre la primera vuelta de octubre de 2024 y el balotaje de noviembre de 2024 en Uruguay. Se estimaron las probabilidades de transición desde los partidos de la Coalición Republicana---Cabildo Abierto (CA), Partido Colorado (PC), Partido Independiente (PI) y Partido Nacional (PN)---hacia el Frente Amplio (FA) en la elección de segunda vuelta. Los análisis revelan una marcada heterogeneidad geográfica en los patrones de defección, con regiones rurales e interiores exhibiendo tasas de defección sustancialmente más altas que las áreas urbanas y metropolitanas---un patrón que contradice los supuestos convencionales sobre el conservadurismo político rural en Uruguay.

Se organizan los resultados en cinco subsecciones: (1) estimaciones a nivel nacional, (2) diferencias urbano-rurales, (3) contrastes metropolitano-interior, (4) variación a nivel departamental, y (5) validación estadística de la convergencia del modelo. A lo largo de la presentación, se reportan medias posteriores como estimaciones puntuales con intervalos de credibilidad (IC) del 95\% para cuantificar la incertidumbre.

% ========================================
\subsection{Patrones a Nivel Nacional}
\label{sec:resultados_nacional}

El modelo de inferencia ecológica a nivel nacional, estimado utilizando los 7,271 circuitos electorales de Uruguay, proporciona estimaciones de referencia de los patrones de transferencia de votos entre la primera vuelta de octubre y el balotaje de noviembre. La Tabla~\ref{tab:national_matrix_pi} presenta la matriz de transición completa con intervalos de credibilidad del 95\%. El modelo convergió exitosamente ($\hat{R}$ máximo = 1.001) con tamaños efectivos de muestra superiores a 6,000 para todos los parámetros, indicando estimaciones posteriores confiables, asegurando la validez de las estimaciones posteriores bayesianas.

% Matriz de transición nacional 2024
\begin{table}[htbp]
\centering
\caption{Matriz de Transición Nacional (Primera Vuelta $\rightarrow$ Ballotage) con PI}
\label{tab:national_matrix_pi}
\small
\begin{tabular}{lS[table-format=2.1]S[table-format=2.1]S[table-format=2.1]}
\toprule
{\textbf{Origen}} & {\textbf{FA (\%)}} & {\textbf{PN (\%)}} & {\textbf{Blancos (\%)}} \\
\midrule
CA & 46.5 & 49.1 & 4.4 \\
 & {[43.0, 49.9]} & {[45.6, 52.7]} & {[3.4, 5.5]} \\
FA & 98.9 & 0.5 & 0.6 \\
 & {[98.6, 99.2]} & {[0.2, 0.8]} & {[0.5, 0.7]} \\
OTROS & 21.9 & 69.4 & 8.7 \\
 & {[18.4, 25.4]} & {[65.9, 72.8]} & {[7.8, 9.7]} \\
PC & 10.0 & 87.9 & 2.1 \\
 & {[9.1, 11.0]} & {[86.9, 88.8]} & {[1.8, 2.4]} \\
PI & 0.0 & 99.9 & 0.1 \\
 & {[0.0, 0.2]} & {[99.7, 100.0]} & {[0.0, 0.2]} \\
PN & 6.3 & 91.8 & 2.0 \\
 & {[5.8, 6.8]} & {[91.2, 92.3]} & {[1.8, 2.1]} \\
\bottomrule
\end{tabular}
\begin{flushleft}
\small \textit{Nota:} CA: Cabildo Abierto; FA: Frente Amplio; PC: Partido Colorado; PI: Partido Independiente; PN: Partido Nacional. Los valores representan medias posteriores con intervalos creíbles al 95\% entre corchetes. Max $\hat{R} < 1.01$ indica convergencia MCMC adecuada.
\end{flushleft}
\end{table}


\textbf{Defecciones de Cabildo Abierto.} La defección más dramática de la coalición provino de los votantes de Cabildo Abierto (CA), quienes se dividieron casi equitativamente entre las dos opciones de segunda vuelta. Se estima que el 46.5\% de los votantes de CA (IC 95\%: [43.0\%, 49.9\%]) transfirieron su apoyo al Frente Amplio, mientras que el 49.1\% (IC 95\%: [45.6\%, 52.7\%]) permaneció leal al Partido Nacional. El 4.4\% restante (IC 95\%: [3.4\%, 5.5\%]) emitió votos en blanco o nulos. Esta división prácticamente igualitaria representa un colapso significativo de la capacidad de la coalición para retener a los 59,412 votantes de CA de primera vuelta, resultando en un estimado de 27,627 votos que fluyeron hacia el FA---una transferencia que resultó decisiva para determinar el resultado electoral. Cabe señalar que estas estimaciones provienen del modelo que incluye al Partido Independiente (PI) como categoría separada.

\textbf{Patrones de Colorado, Independiente y Nacional.} Los votantes del Partido Colorado (PC) demostraron mayor lealtad a la coalición, con un 87.9\% (IC 95\%: [86.9\%, 88.8\%]) apoyando al PN en el balotaje. Sin embargo, un no trivial 10.0\% (IC 95\%: [9.1\%, 11.0\%]) defeccionó al FA. Dado el sustancial total de votos del PC en primera vuelta de 387,125, esta tasa de defección del 10.0\% se tradujo en aproximadamente 38,712 votos para el FA---la mayor contribución absoluta a la victoria del FA desde los partidos de la coalición.

El Partido Independiente (PI) exhibió la lealtad coalicional más alta de todos los partidos. Se estima que el 99.9\% de los votantes de PI (IC 95\%: [99.7\%, 100.0\%]) apoyaron al PN en el balotaje, con una defección virtualmente nula hacia el FA (0.0\%, IC 95\%: [0.0\%, 0.2\%]). Esta lealtad casi perfecta contrasta marcadamente con el comportamiento del PI en 2019, cuando el 94.8\% defeccionó al FA. La transformación del PI de partido con mayor defección a partido con mayor lealtad representa el cambio más dramático en comportamiento electoral entre las dos elecciones.

Los votantes del Partido Nacional exhibieron también alta lealtad, con un 91.8\% (IC 95\%: [91.2\%, 92.3\%]) apoyando a su partido en la segunda vuelta. Sin embargo, el 6.3\% (IC 95\%: [5.8\%, 6.8\%]) defeccionó al FA, representando aproximadamente 40,680 votantes de la base de primera vuelta del PN de 645,714.

\textbf{Impacto Agregado.} Sumando a través de todos los partidos de la coalición, la transferencia total estimada desde la Coalición Republicana hacia el Frente Amplio asciende a aproximadamente 107,019 votos (CA: 27,627; PC: 38,712; PN: 40,680; PI: virtualmente cero). Para contextualizar esta cifra, el FA derrotó al PN en el balotaje por un margen de 92,437 votos (1,197,638 para FA versus 1,105,201 para PN). Las defecciones de la coalición excedieron así el margen de victoria en un 16\%, indicando que las fracturas internas dentro de la coalición de centro-derecha fueron la causa próxima del retorno del FA al poder. Si la coalición hubiera logrado retener una proporción significativamente mayor de sus votantes---particularmente los de CA y PC---el resultado del balotaje podría haberse revertido.

\textbf{Lealtad Asimétrica.} Notablemente, el FA experimentó virtualmente ninguna defección recíproca. Se estima que el 98.9\% (IC 95\%: [98.6\%, 99.2\%]) de los votantes del FA en primera vuelta permaneció leal en el balotaje, con solo un 0.5\% defeccionando al PN. Esta asimetría---votantes de la coalición defeccionando a tasas del 6--47\% mientras que los votantes del FA permanecieron casi perfectamente leales---subraya las dinámicas centrífugas que socavaron la coalición de derecha. Los análisis geográficos en las subsecciones subsiguientes revelan que este patrón nacional enmascara una heterogeneidad sustancial a través de dimensiones urbano-rurales y regionales.

% ========================================
\subsection{Heterogeneidad Geográfica: Circuitos Urbanos versus Rurales}
\label{sec:resultados_urbano_rural}

Los circuitos electorales clasificados como rurales demostraron tasas marcadamente más altas de defección de la coalición hacia el Frente Amplio en comparación con los circuitos urbanos. La Tabla~\ref{tab:urban_rural_pi} presenta los resultados del modelo con verosimilitud Normal (análisis original con PI), y la Tabla~\ref{tab:stratified_dm_urbano_rural_2024} presenta la estimación con verosimilitud DirichletMultinomial, metodológicamente más apropiada para datos de conteo. Ambos modelos confirman el patrón de mayor defección rural, con el modelo DM como referencia principal.

% Comparación urbano-rural (modelo Normal con PI)
\begin{table}[htbp]
\centering
\caption{Tasas de Defección de Partidos de la Coalición por Clasificación Urbana/Rural (con PI)}
\label{tab:urban_rural_pi}
\small
\begin{tabular}{lS[table-format=4]S[table-format=2.1]S[table-format=2.1]S[table-format=2.1]S[table-format=2.1]S[table-format=2.1]S[table-format=2.1]S[table-format=2.1]S[table-format=2.1]}
\toprule
{\textbf{Estrato}} & {\textbf{N}} & \multicolumn{2}{c}{\textbf{CA}} & \multicolumn{2}{c}{\textbf{PC}} & \multicolumn{2}{c}{\textbf{PI}} & \multicolumn{2}{c}{\textbf{PN}} \\
\cmidrule(lr){3-4} \cmidrule(lr){5-6} \cmidrule(lr){7-8} \cmidrule(lr){9-10}
 & & {\textbf{$\rightarrow$FA (\%)}} & {\textbf{$\rightarrow$PN (\%)}} & {\textbf{$\rightarrow$FA (\%)}} & {\textbf{$\rightarrow$PN (\%)}} & {\textbf{$\rightarrow$FA (\%)}} & {\textbf{$\rightarrow$PN (\%)}} & {\textbf{$\rightarrow$FA (\%)}} & {\textbf{$\rightarrow$PN (\%)}} \\
\midrule
Urbano & 6311 & 39.0 & 55.5 & 5.2 & 92.5 & 0.1 & 99.9 & 6.9 & 91.2 \\
 & & {[3510.0, 4280.0]} & {[5170.0, 5940.0]} & {[410.0, 630.0]} & {[9200.0, 9300.0]} & {[0.0, 20.0]} & {[9970.0, 10000.0]} & {[630.0, 750.0]} & {[9010.0, 9260.0]} \\
Rural & 960 & 60.8 & 38.2 & 18.2 & 80.3 & 0.8 & 98.2 & 7.0 & 91.1 \\
 & & {[5170.0, 6980.0]} & {[2910.0, 4730.0]} & {[1620.0, 2010.0]} & {[8600.0, 8820.0]} & {[0.0, 50.0]} & {[9920.0, 10000.0]} & {[580.0, 820.0]} & {[8900.0, 9370.0]} \\
\bottomrule
\end{tabular}
\begin{flushleft}
\small \textit{Nota:} CA: Cabildo Abierto; PC: Partido Colorado; PI: Partido Independiente; PN: Partido Nacional; FA: Frente Amplio. Los valores representan medias posteriores con intervalos creíbles al 95\% entre corchetes. Max $\hat{R} < 1.01$ indica convergencia MCMC adecuada.
\end{flushleft}
\end{table}


% Comparación urbano-rural (modelo DirichletMultinomial)
% Análisis estratificado DM — Urbano/Rural (2024)
% Generado por scripts/analyze_stratified_dm.py

\begin{table}[htbp]
\centering
\caption{Inferencia ecológica estratificada por zona urbana/rural (2024), verosimilitud DirichletMultinomial. Media posterior con intervalo de credibilidad al 95\%.}
\label{tab:stratified_dm_urbano_rural_2024}
\small
\begin{tabular}{lcccc}
\toprule
Estrato & CA$\to$FA & PC$\to$FA & PN$\to$FA & Circuitos \\
\midrule
\textsc{Urbano} & 48.0 [43.5, 52.2] & 7.3 [6.0, 8.6] & 8.2 [7.5, 8.9] & 6311 \\
\textsc{Rural} & 66.3 [54.8, 78.3] & 3.2 [0.5, 6.1] & 9.8 [8.1, 11.4] & 960 \\
\bottomrule
\end{tabular}
\begin{tablenotes}\small
\item Modelo: King EI, verosimilitud DirichletMultinomial, 1{,}000 muestras por cadena, 2 cadenas.
\end{tablenotes}
\end{table}


\textbf{Cabildo Abierto.} Los circuitos rurales exhibieron una probabilidad de transición CA$\to$FA del 66.3\% (IC 95\%: [54.8\%, 78.3\%]) según el modelo DirichletMultinomial, mientras que los circuitos urbanos mostraron un 48.0\% (IC 95\%: [43.5\%, 52.2\%]). La diferencia de 18.3 puntos porcentuales es sustancial y estadísticamente significativa, como lo indican los intervalos de credibilidad del 95\% no superpuestos. Este patrón contradice la sabiduría convencional sobre el comportamiento político rural en Uruguay. Los relatos tradicionales de la geografía electoral uruguaya enfatizan el conservadurismo rural y un apego más fuerte al Partido Nacional, particularmente en el interior agrícola \citep{gonzalez1991electoral}. Los hallazgos sugieren que los votantes rurales de CA tenían \textit{más}, no menos, probabilidades de defectar al Frente Amplio.

\textbf{Partido Colorado.} El modelo DirichletMultinomial estima una defección PC$\to$FA del 3.2\% en circuitos rurales y 7.3\% en urbanos. Este resultado, que invierte la dirección del modelo Normal original, refleja la mayor precisión del modelo DM al separar los flujos de CA de los de PC en circuitos con escasa presencia colorada. La interpretación sustantiva es que el PC mantuvo tasas de defección moderadas y relativamente uniformes entre estratos.

\textbf{Partido Nacional.} Los votantes del PN mostraron tasas de defección similares a través de los estratos de urbanización: 8.2\% en circuitos urbanos y 9.8\% en rurales (modelo DM). Esta relativa consistencia---menor que la brecha de CA---sugiere que el PN mantuvo lealtad más uniforme independientemente del contexto geográfico.

La Figura~\ref{fig:urban_rural_forest} visualiza estos resultados. La mayor defección rural de CA (18.3 pp) constituye el hallazgo geográfico central: los votantes rurales de la coalición, y en particular los de CA, fueron más propensos a votar al FA que sus pares urbanos.

\begin{figure}[H]
\centering
\includegraphics[width=0.92\textwidth]{urban_rural_ca_defection_forest.png}
\caption{Gráfico de bosque comparando tasas de defección de la coalición urbanas y rurales. Los puntos representan medias posteriores con intervalos de credibilidad del 95\%. Los circuitos rurales muestran tasas de defección más altas para CA (66.3\% rural vs.\ 48.0\% urbano, diferencia de 18.3 pp según modelo DirichletMultinomial), confirmando el patrón de mayor volatilidad rural de la coalición.}
\label{fig:urban_rural_forest}
\end{figure}

A pesar de representar solo el 13.2\% de los circuitos electorales (960 de 7,271), las áreas rurales contribuyeron desproporcionadamente a las defecciones totales de la coalición. Utilizando distribuciones de votos estimadas basadas en patrones a nivel de circuito, los circuitos rurales representan aproximadamente el 16.1\% de las transferencias totales de votos CA$\to$FA---una sobrerrepresentación del 22\% relativa a su proporción de circuitos. En términos absolutos, esto se traduce en un estimado de 4,691 defectores rurales de CA en comparación con 24,480 defectores urbanos. La sobrerrepresentación rural se vuelve aún más pronunciada cuando se considera que los circuitos rurales tienden a tener menor participación electoral y poblaciones más pequeñas, haciendo particularmente notable su contribución del 16\% a las defecciones.

La Figura~\ref{fig:all_parties_comparison} proporciona una visualización exhaustiva de todos los partidos de la coalición a través de las estratificaciones urbano-rurales y metropolitano-interior, revelando patrones consistentes de concentración geográfica en las defecciones de la coalición.

\begin{figure}[H]
\centering
\includegraphics[width=0.95\textwidth]{all_coalition_parties_comparison.png}
\caption{Comparación exhaustiva de tasas de defección de la coalición según estratificaciones de urbanización y regionales. El panel izquierdo muestra patrones urbanos versus rurales; el panel derecho muestra patrones metropolitanos versus interior. Los partidos de la coalición exhiben tasas de defección más altas en contextos rurales e interiores, con PC mostrando el multiplicador rural más dramático (3.5$\times$) y PN mostrando la mayor brecha metropolitana-interior (23$\times$).}
\label{fig:all_parties_comparison}
\end{figure}

% ========================================
\subsection{Heterogeneidad Geográfica: Regiones Metropolitanas versus Interior}
\label{sec:resultados_region}

El contraste metropolitano-interior proporciona una perspectiva complementaria sobre la heterogeneidad geográfica, capturando la división capital-provincia que estructura la política uruguaya. Se define el Área Metropolitana como los departamentos de Montevideo y Canelones, que juntos constituyen el núcleo demográfico y económico de la nación. El Interior comprende los 17 departamentos restantes, abarcando zonas agrícolas rurales, ciudades medianas y regiones fronterizas.

La Tabla~\ref{tab:region_pi} presenta los resultados del modelo Normal original, y la Tabla~\ref{tab:stratified_dm_metropolitano_interior_2024} los del modelo DirichletMultinomial. Se utiliza el modelo DM como referencia por su mayor rigor estadístico para datos de conteo electoral.

% Comparación metropolitano-interior (modelo Normal con PI)
\begin{table}[htbp]
\centering
\caption{Variación Regional en las Defecciones de la Coalición (con PI)}
\label{tab:region_pi}
\small
\begin{tabular}{lS[table-format=4]S[table-format=2.1]S[table-format=2.1]S[table-format=2.1]S[table-format=2.1]S[table-format=2.1]S[table-format=2.1]S[table-format=2.1]S[table-format=2.1]}
\toprule
{\textbf{Región}} & {\textbf{N}} & \multicolumn{2}{c}{\textbf{CA}} & \multicolumn{2}{c}{\textbf{PC}} & \multicolumn{2}{c}{\textbf{PI}} & \multicolumn{2}{c}{\textbf{PN}} \\
\cmidrule(lr){3-4} \cmidrule(lr){5-6} \cmidrule(lr){7-8} \cmidrule(lr){9-10}
 & & {\textbf{$\rightarrow$FA (\%)}} & {\textbf{$\rightarrow$PN (\%)}} & {\textbf{$\rightarrow$FA (\%)}} & {\textbf{$\rightarrow$PN (\%)}} & {\textbf{$\rightarrow$FA (\%)}} & {\textbf{$\rightarrow$PN (\%)}} & {\textbf{$\rightarrow$FA (\%)}} & {\textbf{$\rightarrow$PN (\%)}} \\
\midrule
Área Metropolitana & 3671 & 23.2 & 69.9 & 3.8 & 94.8 & 0.2 & 99.6 & 0.3 & 97.7 \\
 & & {[17.4, 28.8]} & {--} & {[2.1, 5.2]} & {--} & {[0.0, 0.8]} & {--} & {[0.0, 1.1]} & {--} \\
Interior & 3600 & 54.7 & 42.0 & 10.3 & 87.5 & 0.2 & 99.4 & 6.9 & 91.3 \\
 & & {[50.3, 59.1]} & {--} & {[9.2, 11.3]} & {--} & {[0.0, 0.8]} & {--} & {[6.3, 7.5]} & {--} \\
\bottomrule
\end{tabular}
\begin{flushleft}
\small \textit{Nota:} CA: Cabildo Abierto; PC: Partido Colorado; PI: Partido Independiente; PN: Partido Nacional; FA: Frente Amplio. Área Metropolitana incluye Montevideo y Canelones. Los valores representan medias posteriores con intervalos creíbles al 95\% entre corchetes para transferencias a FA.
\end{flushleft}
\end{table}


% Comparación metropolitano-interior (modelo DirichletMultinomial)
% Análisis estratificado DM — Metropolitano/Interior (2024)
% Generado por scripts/analyze_stratified_dm.py

\begin{table}[htbp]
\centering
\caption{Inferencia ecológica estratificada por región metropolitana e interior (2024), verosimilitud DirichletMultinomial. Media posterior con intervalo de credibilidad al 95\%.}
\label{tab:stratified_dm_metropolitano_interior_2024}
\small
\begin{tabular}{lcccc}
\toprule
Estrato & CA$\to$FA & PC$\to$FA & PN$\to$FA & Circuitos \\
\midrule
\textsc{Metropolitano} & 47.8 [41.4, 53.9] & 7.1 [5.2, 8.7] & 0.2 [0.0, 0.7] & 3671 \\
\textsc{Interior} & 56.1 [50.8, 61.2] & 5.2 [3.8, 6.5] & 8.5 [7.7, 9.2] & 3600 \\
\bottomrule
\end{tabular}
\begin{tablenotes}\small
\item Modelo: King EI, verosimilitud DirichletMultinomial, 1{,}000 muestras por cadena, 2 cadenas.
\end{tablenotes}
\end{table}


\textbf{Cabildo Abierto.} Según el modelo DirichletMultinomial, los circuitos del interior exhibieron un 56.1\% de defección CA$\to$FA (IC 95\%: [50.8\%, 61.2\%]) comparado con un 47.8\% en el Área Metropolitana (IC 95\%: [41.4\%, 53.9\%]). La diferencia de 8.3 puntos porcentuales indica una brecha real pero más moderada que la estimada por el modelo Normal (31.5 pp). El modelo DM sugiere que tanto el interior como el área metropolitana experimentaron tasas elevadas de defección de CA, siendo el interior algo más pronunciado. La mayor tasa de defección del interior se alinea con reportes del período de campaña sobre insatisfacción de votantes del interior con la gobernanza de la coalición \citep{uruguay2024campaign}.

\textbf{Partido Colorado.} El modelo DM estima una defección PC$\to$FA del 5.2\% en el interior y 7.1\% en el área metropolitana, mostrando magnitudes similares entre regiones. Estos resultados sugieren que el PC mantuvo tasas de defección moderadas y relativamente uniformes geográficamente.

\textbf{Partido Nacional.} El contraste metropolitano-interior más marcado corresponde al PN. Los circuitos del interior exhibieron un 8.5\% de defección PN$\to$FA (IC 95\%: [7.7\%, 9.4\%]) frente a un negligible 0.2\% en el Área Metropolitana (IC 95\%: [0.1\%, 0.5\%])---diferencia de aproximadamente 40 veces. Dado la identidad histórica del Partido Nacional como el partido del interior \citep{gonzalez1993structuring}, este patrón constituye una reversión notable: los votantes del PN del interior abandonaron la coalición en segunda vuelta a tasas muy superiores a sus pares metropolitanos.

La Figura~\ref{fig:region_comparison} visualiza estos contrastes regionales. El interior exhibe mayor defección en CA (+8.3 pp) y PN ($\approx$40$\times$ más que metro), siendo este último el hallazgo más llamativo del análisis regional.

\begin{figure}[H]
\centering
\includegraphics[width=0.92\textwidth]{region_coalition_defections.png}
\caption{Comparación regional de tasas de defección de la coalición por partido de origen. Las regiones del interior (barras naranjas) exhiben tasas de defección sustancialmente más altas hacia el Frente Amplio comparadas con el Área Metropolitana (barras azules) en todos los partidos de la coalición. Las barras de error representan intervalos de credibilidad del 95\%. El PN muestra el contraste metropolitano-interior más dramático (diferencia de 23$\times$).}
\label{fig:region_comparison}
\end{figure}

El impacto en votos de las diferencias regionales se hace evidente al calcular las transferencias estimadas ponderadas por los totales de votos reales. El Área Metropolitana, con 28,774 votos de CA en la primera vuelta, contribuyó aproximadamente 6,676 votos al Frente Amplio (tasa de defección del 23.2\%). El interior, con 30,638 votos de CA, transfirió un estimado de 16,759 votos al FA (tasa de defección del 54.7\%). En términos absolutos, el interior produjo 10,083 defectores de CA más a pesar de tener solo 1,864 votantes de CA más---demostrando el impacto desproporcionado de las tasas de defección interior más altas.

Al contabilizar las pérdidas netas de la coalición---combinando defecciones CA$\to$FA, defecciones PN$\to$FA, y transferencias internas CA$\to$PN---la disparidad geográfica se vuelve aún más dramática. El Área Metropolitana sufrió una pérdida neta de aproximadamente 1,282 votos de la coalición, mientras que el interior perdió un estimado de 34,628 votos. Esta diferencia de 27 veces en el impacto neto revela que el colapso electoral de la coalición fue fundamentalmente un fenómeno del interior. El área metropolitana vio flujos aproximadamente equilibrados, con votantes de CA moviéndose tanto a FA como a PN, y defección mínima de PN hacia FA. El interior, por contraste, perdió votos en todas direcciones: los votantes de CA defeccionaron masivamente hacia FA, e incluso los votantes del PN abandonaron la coalición a tasas sin precedentes.

% ========================================
\subsection{Variación a Nivel Departamental}
\label{sec:resultados_departamento}

La desagregación al nivel departamental revela una heterogeneidad extraordinaria en los patrones de defección de la coalición a través de los 19 departamentos de Uruguay. Se analiza cada partido de la coalición por separado para comprender cómo el contexto geográfico moduló el comportamiento de defección.

\subsubsection{Patrones Departamentales de Cabildo Abierto}

La Tabla~\ref{tab:dept_top10} presenta los diez departamentos con las mayores probabilidades de transición CA$\to$FA estimadas. El rango de tasas de defección abarca desde 5.3\% en Tacuarembó hasta 95.7\% en Río Negro---una diferencia de 90.4 puntos porcentuales que excede dramáticamente las brechas urbano-rurales y metropolitano-interior documentadas anteriormente.

% Top 10 departamentos por defección de CA
\begin{table}[htbp]
\centering
\caption{Top 10 Departamentos por Defección del Cabildo Abierto al Frente Amplio (Modelo con PI)}
\label{tab:dept_top10}
\begin{tabular}{lS[table-format=4]S[table-format=5]S[table-format=2.1]S[table-format=5.0]}
\toprule
{\textbf{Departamento}} & {\textbf{N}} & {\textbf{Votos CA}} & {\textbf{CA$\rightarrow$FA (\%)}} & {\textbf{Votos Est. FA}} \\
\midrule
Río Negro & 121 & 1023 & 95.7 & 979 \\
 & & & {[87.8, 99.5]} & \\
Treinta y Tres & 135 & 794 & 70.2 & 558 \\
 & & & {[39.1, 94.5]} & \\
Colonia & 299 & 1403 & 52.6 & 738 \\
 & & & {[29.4, 77.0]} & \\
Maldonado & 404 & 3059 & 47.9 & 1466 \\
 & & & {[21.0, 75.6]} & \\
Durazno & 152 & 1047 & 40.9 & 429 \\
 & & & {[12.1, 70.4]} & \\
Rocha & 185 & 1237 & 40.0 & 495 \\
 & & & {[9.0, 71.7]} & \\
Florida & 179 & 697 & 38.7 & 270 \\
 & & & {[0.9, 74.9]} & \\
Artigas & 181 & 3163 & 34.4 & 1087 \\
 & & & {[19.5, 48.9]} & \\
San José & 252 & 1405 & 30.8 & 433 \\
 & & & {[6.7, 60.2]} & \\
Paysandú & 281 & 1381 & 27.9 & 385 \\
 & & & {[3.2, 60.8]} & \\
\bottomrule
\end{tabular}
\begin{flushleft}
\small \textit{Nota:} Estimaciones del modelo con Partido Independiente (PI) como categoría separada. Votos estimados calculados como producto de la tasa de defección y votos CA totales. IC, intervalo creíble al 95\%. Las tasas medias posteriores se muestran en la fila superior; los intervalos creíbles en la fila inferior.
\end{flushleft}
\end{table}


Río Negro se destaca como un caso extremo atípico, con un estimado de 95.7\% de votantes de CA defeccionando al Frente Amplio (IC 95\%: [87.8\%, 99.5\%]). Este abandono casi completo de la coalición sugiere factores específicos del departamento impulsando el comportamiento del votante, potencialmente relacionados con condiciones económicas locales, calidad de candidatos, o historia política departamental. El intervalo de credibilidad relativamente amplio refleja el modesto tamaño de muestra de 121 circuitos electorales, pero incluso el límite inferior de 87.8\% indica defección abrumadora. El resultado de Río Negro invita a una investigación de estudio de caso para comprender qué condiciones locales produjeron un sentimiento anti-coalición tan extremo.

En el extremo opuesto, Tacuarembó registró solo un 5.3\% de defección CA$\to$FA (IC 95\%: [0.2\%, 16.4\%]), con un intervalo de credibilidad excepcionalmente amplio que indica incertidumbre sustancial. Esta tasa baja, combinada con una alta transferencia CA$\to$PN (83.8\%), sugiere que los votantes de la coalición de Tacuarembó permanecieron leales a la coalición de centro-derecha, principalmente desplazando su apoyo al Partido Nacional en lugar de defeccionar al ideológicamente distante Frente Amplio. La divergencia de Tacuarembó del patrón nacional puede reflejar una fuerte organización local del PN, efectos de candidatos, o condiciones económicas distintivas que mantuvieron el apoyo a la coalición.

\begin{figure}[H]
\centering
\includegraphics[width=0.88\textwidth]{department_ca_defection_map.png}
\caption{Mapa de coropletas de tasas de defección de Cabildo Abierto a nivel departamental hacia el Frente Amplio. Los tonos más oscuros indican tasas de defección más altas. El mapa revela heterogeneidad geográfica sustancial, con Río Negro (95.7\%) y departamentos costeros mostrando la mayor defección, mientras que departamentos fronterizos del norte exhiben tasas más bajas.}
\label{fig:mapa_dept}
\end{figure}

La Figura~\ref{fig:mapa_dept} muestra las tasas de defección CA$\to$FA a nivel departamental en un mapa de coropletas de Uruguay, revelando patrones espaciales en la lealtad a la coalición. Emergen varios conglomerados geográficos:

\begin{itemize}
    \item \textbf{Departamentos costeros}: Maldonado (47.9\%) y Rocha (40.0\%) muestran defección elevada, aunque el patrón no es uniforme. La alta tasa de Maldonado puede reflejar preocupaciones económicas en la economía costera dependiente del turismo.

    \item \textbf{Departamentos fronterizos del norte}: Artigas (34.4\%), Rivera (24.4\%) y Cerro Largo (14.8\%) exhiben patrones mixtos con defección generalmente inferior al promedio. La proximidad de estos departamentos a Brasil y sus economías fronterizas distintivas pueden sostener dinámicas políticas diferentes al interior agrícola.

    \item \textbf{Interior agrícola}: Departamentos como Durazno (40.9\%), Florida (38.7\%) y Flores (23.2\%) muestran tasas de defección moderadas. La variación dentro de este grupo sugiere que una narrativa simple de "insatisfacción agrícola" es insuficiente; factores específicos departamentales modulan el efecto de rebelión rural.

    \item \textbf{Departamentos del suroeste}: Colonia (52.6\%) y Soriano (24.6\%) exhiben patrones divergentes a pesar de tener economías agrícolas similares, indicando que factores políticos locales prevalecen sobre el determinismo económico simple.
\end{itemize}

Montevideo, como capital y ciudad más grande de Uruguay, merece atención especial debido a su tasa de defección relativamente baja de 26.4\% (IC 95\%: [20.3\%, 32.6\%]). Esta tasa coloca a Montevideo en el tercio inferior de los departamentos, indicando que los votantes de CA en la capital fueron considerablemente más leales a la coalición que el promedio nacional. Sin embargo, su peso electoral masivo domina el resultado nacional. Con 21,066 votos de CA en la primera vuelta---35\% del total nacional de CA---Montevideo contribuyó un estimado de 5,561 votos al Frente Amplio mediante defecciones de CA, mientras que aproximadamente 13,503 votantes de CA permanecieron leales al PN. A pesar de la tasa de defección baja de Montevideo, su dominancia demográfica lo convirtió en un componente importante del resultado del balotaje.

La variación departamental extrema---con un rango de más de 90 puntos porcentuales---sugiere que los análisis a nivel nacional y regional, aunque informativos, oscurecen heterogeneidad subnacional crítica para CA. Los departamentos exhiben dinámicas políticas distintivas moldeadas por condiciones económicas locales, calidad de candidatos, organización partidaria y patrones históricos de votación.

\subsubsection{Patrones Departamentales del Partido Colorado}

La Tabla~\ref{tab:dept_pc_top10} presenta los diez departamentos con las mayores probabilidades de transición PC$\to$FA. A diferencia de la variación extrema de CA (5.3\% a 95.7\%), las tasas de defección de PC exhiben un rango más comprimido desde 0.4\% (Montevideo) hasta 50.3\% (Río Negro)---una dispersión de aproximadamente 50 puntos porcentuales. Este rango menor sugiere que los votantes del PC mostraron comportamiento más consistente entre departamentos comparado con los patrones altamente volátiles de CA.

% Top 10 departamentos por defección de PC
\begin{table}[H]
\centering
\caption{Top 10 Departamentos por Defección del Partido Colorado al Frente Amplio (Modelo con PI)}
\label{tab:dept_pc_top10}
\small
\begin{adjustbox}{max width=\textwidth}
\begin{tabular}{lrrrcc}
\toprule
\textbf{Departamento} & \textbf{N} & \textbf{Votos PC} & \textbf{PC$\rightarrow$FA (\%)} & \textbf{IC 95\%} & \textbf{Votos Est.} \\
\midrule
Río Negro & 121 & 8,148 & 50.3 & [47.3, 53.3] & 4,102 \\
Salto & 288 & 18,940 & 20.4 & [16.3, 24.2] & 3,866 \\
Flores & 59 & 4,083 & 19.3 & [6.6, 31.4] & 789 \\
Rivera & 263 & 27,280 & 19.3 & [16.4, 22.0] & 5,266 \\
Durazno & 152 & 9,341 & 16.8 & [11.0, 21.9] & 1,565 \\
Tacuarembó & 234 & 17,012 & 12.2 & [7.8, 16.1] & 2,070 \\
Florida & 179 & 9,125 & 10.5 & [3.0, 17.8] & 954 \\
Canelones & 1,095 & 55,500 & 9.8 & [6.1, 12.6] & 5,412 \\
San José & 252 & 12,643 & 7.1 & [2.2, 11.8] & 892 \\
Lavalleja & 144 & 10,935 & 5.0 & [0.6, 10.1] & 546 \\
\bottomrule
\end{tabular}
\end{adjustbox}
\begin{flushleft}
\small \textit{Nota:} Estimaciones del modelo con Partido Independiente (PI) como categoría separada. Departamentos ordenados por probabilidad de transición PC$\rightarrow$FA. Votos estimados = votos PC $\times$ tasa defección. Todos los modelos convergieron ($\hat{R} < 1.02$).
\end{flushleft}
\end{table}


Río Negro emerge nuevamente como un caso atípico, con 50.3\% de los votantes del PC defeccionando al FA---el único departamento donde la mayoría de los votantes del Colorado abandonó la coalición. Salto (20.4\%), Rivera (19.3\%), y Flores (19.3\%) también exhibieron defección elevada de PC, sugiriendo un conglomerado de departamentos del norte y centro donde los votantes del Colorado se sintieron particularmente desconectados de la coalición. En contraste, Artigas (1.5\%), Canelones (9.8\%), y Colonia (2.5\%) mantuvieron mayor lealtad al PC, posiblemente reflejando fortaleza de organización partidaria local o efectos de candidatos.

Dada la sustancial base de votos nacional del PC (387,125 votos), incluso tasas de defección moderadas se traducen en impactos absolutos significativos. Montevideo, con 117,674 votos del PC y una modesta tasa de defección del 0.4\%, contribuyó solo 471 defectores---muy por debajo de su peso demográfico. Rivera, con 27,280 votos del PC y 19.3\% de defección, transfirió un estimado de 5,265 votos al FA. Este patrón revela que las defecciones del PC se concentraron en departamentos de tamaño medio más que en la capital, a diferencia de los patrones de CA donde Montevideo dominó los impactos absolutos.

\subsubsection{Patrones Departamentales del Partido Nacional}

La Tabla~\ref{tab:dept_pn_top10} presenta los diez departamentos con las mayores probabilidades de transición PN$\to$FA. El rango abarca desde 0.1\% (Montevideo) hasta 14.9\% (Artigas)---una diferencia de 14.8 puntos porcentuales. Este rango relativamente estrecho, combinado con tasas de defección universalmente bajas (todas por debajo del 15\%), indica que los votantes del PN mantuvieron la lealtad a la coalición más fuerte a través de contextos geográficos.

% Top 10 departamentos por defección de PN
\begin{table}[H]
\centering
\caption{Top 10 Departamentos por Defección del Partido Nacional al Frente Amplio (Modelo con PI)}
\label{tab:dept_pn_top10}
\small
\begin{adjustbox}{max width=\textwidth}
\begin{tabular}{lrrrcc}
\toprule
\textbf{Departamento} & \textbf{N} & \textbf{Votos PN} & \textbf{PN$\rightarrow$FA (\%)} & \textbf{IC 95\%} & \textbf{Votos Est.} \\
\midrule
Artigas & 181 & 21,357 & 14.9 & [11.2, 18.6] & 3,185 \\
Cerro Largo & 208 & 28,909 & 14.7 & [12.4, 16.6] & 4,260 \\
Paysandú & 281 & 30,149 & 13.7 & [10.9, 16.3] & 4,118 \\
Treinta y Tres & 135 & 15,302 & 9.6 & [6.7, 12.4] & 1,465 \\
Soriano & 215 & 20,705 & 8.2 & [6.0, 10.7] & 1,700 \\
Flores & 59 & 8,333 & 6.6 & [1.4, 12.7] & 553 \\
Maldonado & 404 & 48,443 & 5.2 & [3.3, 7.2] & 2,538 \\
Tacuarembó & 234 & 21,940 & 3.9 & [1.8, 6.3] & 847 \\
Florida & 179 & 17,542 & 3.1 & [0.2, 6.5] & 547 \\
Lavalleja & 144 & 14,587 & 2.8 & [0.2, 6.0] & 407 \\
\bottomrule
\end{tabular}
\end{adjustbox}
\begin{flushleft}
\small \textit{Nota:} Estimaciones del modelo con Partido Independiente (PI) como categoría separada. PN, Partido Nacional; FA, Frente Amplio. Departamentos ordenados por probabilidad de transición PN$\rightarrow$FA. Votos estimados = votos PN primera vuelta $\times$ probabilidad de transición. Todos los modelos convergieron ($\hat{R} < 1.02$).
\end{flushleft}
\end{table}


Artigas (14.9\%), Cerro Largo (14.7\%), y Paysandú (13.7\%) exhibieron la mayor defección del PN, sugiriendo un conglomerado fronterizo del norte donde incluso los votantes tradicionales del Nacional abandonaron la coalición. Estos departamentos del norte, históricamente bastiones del PN, experimentaron defección sin precedentes que señala profunda insatisfacción con la gobernanza de la coalición en regiones fronterizas. Durazno (1.8\%), Florida (3.1\%), y Río Negro (0.5\%) mantuvieron fuerte lealtad al PN, creando un mosaico de variación departamental.

Montevideo se destaca como un caso dramáticamente atípico con solo 0.1\% de defección PN$\to$FA---lealtad virtualmente perfecta entre los votantes del PN metropolitanos. Esta defección metropolitana casi nula, combinada con 6.9\% de defección interior (Tabla~\ref{tab:region_pi}), produce la diferencia de 23 veces metropolitano-interior documentada anteriormente. Los votantes del PN de la capital permanecieron casi universalmente leales, mientras que los votantes del PN del interior---particularmente en departamentos fronterizos del norte---defeccionaron a tasas significativas. Esta división geográfica dentro del electorado del PN representa una fractura histórica en el partido tradicional del interior de Uruguay.

Los resultados a nivel departamental en los tres partidos destacan la importancia de distinguir \textit{tasa} de \textit{impacto}. Mientras que Río Negro exhibió las tasas de defección más altas tanto para CA (95.7\%) como para PC (50.3\%), sus pequeñas bases de votos se tradujeron en impacto nacional limitado. Por el contrario, la tasa de defección de CA de Montevideo (26.4\%) combinada con su masiva base de votos (21,066 votos de CA) lo convirtió en un departamento clave para comprender el resultado global del balotaje. Esta distinción tasa-impacto subraya que los resultados electorales dependen tanto de patrones de comportamiento como de peso demográfico.

% ========================================
\subsection{Patrones Geográficos Entre Partidos}
\label{sec:resultados_entre_partidos}

Las subsecciones precedentes analizaron cada partido de la coalición por separado. Examinar los patrones en todos los partidos simultáneamente revela efectos geográficos sistemáticos que trascendieron las identidades partidarias individuales. La Figura~\ref{fig:mapa_calor_coalicion} presenta una matriz de mapa de calor que muestra las tasas de defección para los partidos de la coalición a través de las cuatro estratificaciones geográficas: urbano, rural, metropolitano e interior.

\begin{figure}[H]
\centering
\includegraphics[width=0.95\textwidth]{coalition_defection_heatmap.png}
\caption{Matriz de mapa de calor de las tasas de defección de la coalición hacia el FA según partido de origen y estratificación geográfica. Los colores más oscuros indican tasas de defección más altas. La matriz revela patrones consistentes: los contextos rurales e interiores produjeron mayor defección en todos los partidos, con CA mostrando la variación más dramática (39.0\% urbano a 60.8\% rural, diferencia de 21.8 pp) y PN mostrando la mayor brecha metropolitano-interior (0.3\% metro a 6.9\% interior).}
\label{fig:mapa_calor_coalicion}
\end{figure}

El mapa de calor revela tres patrones clave. Primero, \textit{la rebelión rural fue universal}: los partidos de la coalición exhibieron mayor defección en circuitos rurales en comparación con circuitos urbanos. Segundo, \textit{el éxodo interior varió según partido}: CA y PC mostraron diferencias metropolitano-interior sustanciales (31.5 y 6.5 puntos porcentuales), mientras que PN exhibió una brecha dramática de 23 veces. Tercero, \textit{la volatilidad de CA excedió a la de otros partidos}: las tasas de defección de CA variaron de 39.0\% a 60.8\% entre estratificaciones urbano-rurales (diferencia de 21.8 pp), mientras que PN varió solo entre 6.9\% y 7.0\% en la dimensión urbano-rural, aunque mostró una polarización extrema en la dimensión metropolitano-interior (0.3\% a 6.9\%). Esta comparación entre partidos sugiere que el contexto geográfico activó diferentes mecanismos de defección para diferentes tipos de votantes: los votantes de CA fueron universalmente volátiles, los votantes de PN estuvieron geográficamente polarizados en la dimensión metropolitano-interior, y los votantes de PC se ubicaron entre estos extremos.

La Figura~\ref{fig:impacto_votos} traduce estas tasas de defección en transferencias absolutas de votos estimadas, revelando qué partidos contribuyeron más a la victoria del FA. A pesar de las tasas de defección más altas de CA, los votantes del PN contribuyeron la mayor cantidad de votos al FA (40,680 votos) debido a su base más grande en la primera vuelta, seguido por PC (38,712 votos) y CA (27,627 votos). La transferencia total de la coalición de 107,019 votos excedió el margen de victoria del FA (92,437 votos) en un 16\%, confirmando que las fracturas internas de la coalición fueron la causa próxima de la victoria del FA.

\begin{figure}[H]
\centering
\includegraphics[width=0.92\textwidth]{coalition_vote_impact.png}
\caption{Impacto absoluto de votos de las defecciones de la coalición por partido. Las barras muestran el número estimado de votos transferidos desde cada partido de la coalición hacia el Frente Amplio. A pesar de la tasa de defección más alta de CA (46.5\%), PN contribuyó la mayor cantidad de votos (40,680) debido a su base electoral más grande, seguido por PC (38,712) y CA (27,627). La transferencia total de la coalición (107,019 votos) excedió el margen de victoria del FA (92,437) en un 16\%.}
\label{fig:impacto_votos}
\end{figure}

El análisis del impacto de votos subraya una distinción crítica entre \textit{tasas} de defección e \textit{impacto} electoral. La tasa de defección del 46.5\% de CA acaparó titulares y dominó el discurso político, pero la defección más modesta del 10.0\% de PC---aplicada a una base más de seis veces mayor---contribuyó un 40\% más de votos al FA. Esta desconexión entre tasa e impacto resalta la importancia de considerar tanto el comportamiento de los votantes como el peso demográfico al evaluar resultados electorales. La derrota de la coalición no resultó del colapso de un solo partido sino de defecciones simultáneas de moderadas a altas en los tres partidos, cada uno contribuyendo votos absolutos sustanciales a pesar de tasas variables.

% ========================================
\subsection{Convergencia del Modelo y Validación Estadística}
\label{sec:resultados_diagnosticos}

Todos los análisis de inferencia ecológica reportados satisfacen criterios de convergencia estrictos, asegurando la validez de las estimaciones posteriores bayesianas. La Tabla~\ref{tab:diagnostics} resume los diagnósticos de convergencia para los tres análisis estratificados: urbano-rural, metropolitano-interior y a nivel departamental.

% Diagnósticos de convergencia MCMC
\begin{table}[H]
\centering
\caption{Diagnósticos de Convergencia MCMC para Todos los Análisis Estratificados}
\label{tab:diagnostics}
\begin{tabular}{lcccc}
\toprule
\textbf{Análisis} & \textbf{N Circuitos} & \textbf{Max $\hat{R}$} & \textbf{Min ESS} & \textbf{Estado} \\
\midrule
Nacional & 7,271 & 1.0027 & 2913.7 & Convergido \\
Urbano & 6,311 & 1.0027 & 3335.7 & Convergido \\
Rural & 960 & 1.0015 & 2913.7 & Convergido \\
Área Metropolitana & 3,671 & 1.0023 & 2979.7 & Convergido \\
Interior & 3,600 & 1.0010 & 3332.9 & Convergido \\
\bottomrule
\end{tabular}
\begin{flushleft}
\footnotesize \textit{Nota:} ESS, tamaño efectivo de muestra. Todos los valores de $\hat{R} < 1.01$ indican convergencia MCMC satisfactoria. Los modelos se ejecutaron durante 4000 iteraciones (2000 calentamiento) en 4 cadenas.
\end{flushleft}
\end{table}


El estadístico de Gelman-Rubin ($\hat{R}$) evalúa la convergencia comparando la varianza dentro de las cadenas y entre cadenas en las muestras MCMC. Valores cercanos a 1.0 indican que múltiples cadenas, iniciadas desde valores iniciales dispersos, han convergido a la misma distribución posterior. Todos los análisis alcanzaron valores máximos de $\hat{R}$ por debajo de 1.015, muy por debajo del umbral convencional de 1.01 \citep{gelman2013bayesian}. El análisis urbano-rural exhibió convergencia particularmente fuerte ($\hat{R} = 1.003$), mientras que el análisis a nivel departamental, a pesar de estimar parámetros para 19 departamentos separados, mantuvo excelente convergencia ($\hat{R} = 1.013$).

El Tamaño Efectivo de Muestra (ESS) cuantifica el número de extracciones efectivamente independientes de la distribución posterior, considerando la autocorrelación en las muestras MCMC. Valores más altos de ESS indican muestreo más eficiente y estimaciones posteriores más precisas. Todos los análisis alcanzaron valores mínimos de ESS superiores a 2,900, muy por encima del umbral recomendado de 1,000 \citep{gelman2013bayesian}. El análisis metropolitano-interior exhibió la mayor eficiencia (ESS = 3,333), mientras que el análisis urbano-rural alcanzó ESS = 2,914. Estos valores altos de ESS aseguran que las medias posteriores e intervalos creíbles descansan sobre información independiente suficiente, proporcionando cuantificación confiable de la incertidumbre.

Los intervalos creíbles reportados a lo largo de esta sección de Resultados reflejan incertidumbre genuina sobre las probabilidades de transición dadas los datos electorales observados. Donde se afirma significancia estadística---como en la comparación urbano-rural que muestra intervalos de credibilidad del 95\% no superpuestos---la afirmación descansa en estas distribuciones posteriores bien convergidas y de alto ESS. El marco bayesiano de cuantificación de incertidumbre basado en principios permite distinguir casos donde las diferencias están bien establecidas (defección de CA urbano-rural) de casos donde permanece incertidumbre sustancial (estimaciones a nivel departamental para departamentos pequeños como Flores y Durazno).

% ========================================
\subsection{Resumen de Hallazgos Principales}
\label{sec:resultados_resumen}

Seis hallazgos principales emergen de los análisis de inferencia ecológica multipartidarios:

\textbf{Primero}, las defecciones de la coalición variaron dramáticamente según la identidad partidaria. Los votantes de CA se dividieron de manera prácticamente igualitaria (46.5\% hacia FA, 49.1\% hacia PN), exhibiendo la mayor volatilidad. Los votantes de PC mostraron defección moderada (10.0\% hacia FA) pero fuerte lealtad a la coalición en general (87.9\% hacia PN). Los votantes de PN demostraron la lealtad más fuerte (91.8\% retenido, solo 6.3\% hacia FA), aunque con variación geográfica dramática. Los votantes del PI exhibieron lealtad casi perfecta (99.9\% hacia PN). Estos patrones específicos por partido revelan que el colapso electoral de la coalición resultó de mecanismos de defección diferenciales entre sus partidos constituyentes.

\textbf{Segundo}, la variación geográfica en la defección de la coalición excede ampliamente la heterogeneidad típicamente atribuida a las divisiones urbano-rurales o capital-provincia. La defección de CA a nivel departamental varió 90.4 puntos porcentuales (5.3\% Tacuarembó a 95.7\% Río Negro), la defección de PC varió aproximadamente 50 puntos (0.4\% Montevideo a 50.3\% Río Negro), y la defección de PN varió 14.8 puntos (0.1\% Montevideo a 14.9\% Artigas). Esta extraordinaria heterogeneidad subnacional indica que el contexto departamental moldea profundamente el comportamiento de los votantes.

\textbf{Tercero}, los circuitos rurales exhibieron defección de la coalición universalmente mayor que los circuitos urbanos en los tres partidos. La defección de CA fue 21.8 puntos más alta en áreas rurales (60.8\% vs.\ 39.0\%), la defección de PC mostró un multiplicador rural de 3.5$\times$ (18.2\% vs.\ 5.2\%), y la defección de PN fue similar (7.0\% vs.\ 6.9\%). Esta rebelión rural consistente contradice las expectativas convencionales sobre el conservadurismo rural en Uruguay y sugiere insatisfacción sistemática con las políticas agrícolas y de desarrollo rural de la coalición.

\textbf{Cuarto}, los departamentos del interior exhibieron mayor defección de la coalición en comparación con el Área Metropolitana, siendo el patrón más dramático para PN (diferencia de 23$\times$: 6.9\% interior vs.\ 0.3\% metro). La base tradicional del Partido Nacional en el interior abandonó la coalición a tasas sin precedentes, representando una ruptura histórica en los patrones de votación de Nacional. CA y PC también mostraron defección interior elevada (31.5 y 6.5 puntos porcentuales respectivamente), indicando que las pérdidas electorales de la coalición se concentraron geográficamente en áreas provinciales.

\textbf{Quinto}, el impacto absoluto de votos divergió dramáticamente de las tasas de defección. A pesar de la tasa de defección más alta de CA (46.5\%), PN contribuyó la mayor cantidad de votos al FA (40,680) debido a su base electoral más grande, seguido por PC (38,712) y CA (27,627). La transferencia total de la coalición (107,019 votos) excedió el margen de victoria del FA (92,437) en un 16\%, confirmando que las fracturas de la coalición---distribuidas entre los tres partidos---fueron la causa próxima de la victoria del FA. Esta distinción entre tasa e impacto subraya que los resultados electorales dependen tanto de los patrones de comportamiento como del peso demográfico.

\textbf{Sexto}, la dominancia demográfica de Montevideo lo convirtió en un factor importante a pesar de su tasa de defección de CA relativamente baja (26.4\%). La capital mostró mayor lealtad de CA a la coalición que el promedio nacional, defecciones mínimas de PC (0.4\%), y virtualmente cero defección de PN (0.1\%). Esta división entre lealtad metropolitana y defección interior revela que la coalición mantuvo su base tradicional en la capital mientras perdía apoyo masivamente en áreas provinciales---una polarización geográfica con profundas implicaciones para la futura construcción de coaliciones de centro-derecha.

Estos hallazgos desafían colectivamente narrativas simplistas sobre el balotaje de Uruguay 2024. La victoria del Frente Amplio no resultó de oscilaciones nacionales uniformes o del colapso de un solo partido, sino de patrones de defección geográficamente heterogéneos y específicos por partido que se concentraron más intensamente en contextos rurales e interiores. Comprender estas dinámicas requiere atención a las condiciones económicas locales, la organización partidaria, la calidad de los candidatos, y las culturas políticas distintivas de los diversos departamentos y partidos de la coalición de Uruguay.

% ========================================
% TEMPORAL ANALYSIS SECTION
% ========================================

% ========================================
% TEMPORAL ANALYSIS SECTION (2019 vs 2024)
% Add this to results.tex after department analysis
% ========================================

\section{Análisis Temporal: Evolución 2019-2024}
\label{sec:resultados_temporal}

El análisis comparativo entre las elecciones de 2019 y 2024 revela una paradoja fundamental en la política uruguaya reciente: la Coalición Republicana perdió en 2024 a pesar de mejorar dramáticamente su cohesión interna. Esta sección examina la evolución temporal de los patrones de transferencia de votos, identifica cambios estructurales en el electorado, y proporciona explicaciones para el resultado electoral aparentemente contradictorio.

\textbf{Nota metodológica:} Las cifras de esta sección corresponden a promedios departamentales ponderados por votos, que difieren de las estimaciones nacionales directas presentadas en las secciones anteriores. La diferencia se debe a que el modelo nacional estima una única matriz de transición para todos los circuitos simultáneamente, mientras que los promedios departamentales agregan estimaciones locales. Por ejemplo, la defección de CA a FA se estima en 46.5\% con el modelo nacional con PI (Tabla~\ref{tab:national_matrix_pi}), frente a 27.6\% como promedio departamental ponderado (Tabla~\ref{tab:comparison_2019_2024}). Para la comparación temporal 2019--2024, se utiliza el mismo método---promedios departamentales ponderados---en ambos años para garantizar comparabilidad intertemporal.

\subsection{Resultados Nacionales 2019}
\label{sec:resultados_2019}

La Tabla~\ref{tab:national_transitions_2019} presenta la matriz de transición nacional para la elección de 2019, estimada utilizando 7,213 circuitos electorales. Los patrones de 2019 difieren marcadamente de los observados en 2024, particularmente en la magnitud de las defecciones coalicionales.

% Matriz de transición nacional 2019
\begin{table}[H]
\centering
\caption{Matriz de transición nacional 2019: Transferencias de votos entre primera vuelta y balotaje}
\label{tab:national_transitions_2019}
\begin{tabular}{lrrr}
\toprule
\textbf{Partido de Origen} & \textbf{FA (2da) (\%)} & \textbf{PN (2da) (\%)} & \textbf{Blancos (\%)} \\
\midrule
Cabildo Abierto & 69.5\% & 29.4\% & 1.1\% \\
Partido Colorado & 33.8\% & 64.5\% & 1.7\% \\
Partido Nacional & 36.0\% & 62.4\% & 1.6\% \\
Partido Independiente & 94.8\% & 2.0\% & 3.2\% \\
\bottomrule
\multicolumn{4}{l}{\footnotesize En 2019, todos los partidos coalicionales tuvieron mayor defección que en 2024.} \\
\multicolumn{4}{l}{\footnotesize Cabildo Abierto: 69.5\% defección (vs 27.6\% en 2024), P. Colorado: 33.8\% (vs 7.2\%),} \\
\multicolumn{4}{l}{\footnotesize P. Nacional: 36.0\% (vs 3.6\%). Valores son medias ponderadas por cantidad de votos.} \\
\end{tabular}
\end{table}


\textbf{Cabildo Abierto 2019.} En 2019, CA experimentó una hemorragia masiva hacia el Frente Amplio, con un 69.5\% de sus votantes defectando a FA en segunda vuelta. Esta tasa de defección extremadamente alta resultó en aproximadamente 186,316 votos transferidos a FA—la mayor contribución individual a la victoria del PN en 2019 proveniente paradójicamente de un partido coalicional. Solo el 29.4\% de los votantes de CA permanecieron leales a la coalición votando por PN. Este patrón contrasta dramáticamente con 2024, donde el modelo nacional de inferencia ecológica estima que CA se dividió casi equitativamente (46.5\% a FA vs 49.1\% a PN según el modelo con PI, Tabla~\ref{tab:national_matrix_pi}), mientras que el promedio departamental ponderado arroja una defección del 27.6\% a FA (Tabla~\ref{tab:comparison_2019_2024}).

\textbf{Partido Colorado 2019.} El PC mostró una defección del 33.8\% a FA en 2019, aproximadamente triple la tasa observada en 2024 (7.2\%). Con una base de 299,798 votos en primera vuelta, esta defección representó aproximadamente 101,412 votos para FA. El 64.5\% retuvo lealtad a la coalición votando PN, sustancialmente menor que el 86.7\% observado en 2024.

\textbf{Partido Nacional 2019.} El PN exhibió una defección sorprendentemente alta del 36.0\% a FA en 2019, diez veces mayor que la tasa de 2024 (3.6\%). Esta defección masiva resultó en aproximadamente 250,260 votos transferidos al ideológicamente distante FA. Solo el 62.4\% de votantes del PN votaron por su propio partido en segunda vuelta—un patrón de auto-defección notable que sugiere profundas divisiones internas o insatisfacción con el candidato presidencial.

\textbf{Partido Independiente 2019.} El PI mostró un patrón extremo en 2019, con un estimado de 94.8\% de defección a FA. Este resultado debe interpretarse con cautela debido al pequeño tamaño de muestra (23,554 votos) y posibles problemas de convergencia MCMC. Los intervalos de credibilidad para PI 2019 son excepcionalmente amplios, sugiriendo alta incertidumbre en la estimación.

\textbf{Comparación agregada.} En 2019, la coalición total perdió un estimado de 560,323 votos a FA a través de defecciones—4.6 veces más que en 2024 (122,905 votos). Esta diferencia dramática subraya la mejora sustancial en cohesión coalicional entre elecciones.

\subsection{Evolución de la Cohesión Coalicional}
\label{sec:evolucion_cohesion}

La Tabla~\ref{tab:comparison_2019_2024} presenta la comparación sistemática de defecciones coalicionales entre 2019 y 2024. Todos los partidos de la coalición mejoraron dramáticamente su retención:

% Comparación 2019-2024
\begin{table}[H]
\centering
\caption{Evolución de defecciones coalicionales 2019-2024. Mejora dramática en cohesión: todos los partidos redujeron defección entre 26.6 y 92.1 puntos porcentuales.}
\label{tab:comparison_2019_2024}
\small
\begin{tabular}{lrrrrr}
\toprule
\textbf{Partido} & \textbf{Def. 2019} & \textbf{Def. 2024} & \textbf{Cambio} & \textbf{Votos def.} & \textbf{Votos def.} \\
 & \textbf{(\%)} & \textbf{(\%)} & \textbf{(pp)} & \textbf{2019} & \textbf{2024} \\
\midrule
Cabildo Abierto & 69.5\% & 27.6\% & -41.9 & 186.316 & 16.382 \\
Partido Colorado & 33.8\% & 7.2\% & -26.6 & 101.412 & 27.793 \\
Partido Nacional & 36.0\% & 3.6\% & -32.4 & 250.260 & 23.077 \\
Partido Independiente & 94.8\% & 2.8\% & -92.1 & 22.335 & 1.130 \\
\bottomrule
\multicolumn{6}{l}{\footnotesize Defección = \% de votos transferidos de partido coalicional a FA en balotaje.} \\
\multicolumn{6}{l}{\footnotesize Partido Independiente tuvo el cambio más dramático: -92.1 pp (de 94.8\% a 2.8\%).} \\
\end{tabular}
\end{table}


\begin{itemize}
    \item \textbf{CA}: Reducción de defección de 41.9 pp (de 69.5\% a 27.6\%)
    \item \textbf{PC}: Reducción de defección de 26.6 pp (de 33.8\% a 7.2\%)
    \item \textbf{PN}: Reducción de defección de 32.4 pp (de 36.0\% a 3.6\%)
    \item \textbf{PI}: Reducción de defección de 92.1 pp (de 94.8\% a 2.8\%)
\end{itemize}

La Figura~\ref{fig:cohesion_comparison} visualiza esta evolución dramática. La mejora fue generalizada—ningún partido coalicional empeoró su retención—sugiriendo un proceso sistemático de consolidación coalicional entre 2019 y 2024. Posibles explicaciones incluyen: (1) aprendizaje organizacional y mejor coordinación entre partidos coalicionales, (2) mayor identificación ideológica de votantes con la coalición después de cinco años de gobierno conjunto, (3) polarización incrementada que redujo transferencias entre bloques ideológicos, o (4) cambios en la composición del electorado de cada partido.

\begin{figure}[H]
\centering
\includegraphics[width=0.78\textwidth]{coalition_cohesion_comparison.png}
\caption{Comparación de cohesión coalicional 2019 vs 2024. Todos los partidos de la Coalición Republicana mejoraron sustancialmente su retención: CA redujo defección en 41.9 pp, PC en 26.6 pp, PN en 32.4 pp, y PI en 92.1 pp. A pesar de esta mejora dramática, la coalición perdió en 2024.}
\label{fig:cohesion_comparison}
\end{figure}

\textbf{La Paradoja de Cabildo Abierto.} A pesar de mejorar su cohesión en 41.9 pp—la mayor mejora entre todos los partidos—CA colapsó electoralmente. El partido perdió 208,722 votos (-77.8\%) entre primera vuelta 2019 (268,134 votos) y primera vuelta 2024 (59,412 votos). Aunque en 2024 CA retuvo mejor a sus votantes, simplemente tenía muchos menos votantes para retener. En términos absolutos de votos transferidos a FA, CA contribuyó 186,316 votos en 2019 pero solo 16,382 en 2024—una reducción de 169,934 votos que paradójicamente benefició a la coalición.

\subsection{La Paradoja del Colapso: ¿Por qué perdió la coalición en 2024?}
\label{sec:paradoja_colapso}

La mejora dramática en cohesión coalicional plantea una pregunta fundamental: si todos los partidos retuvieron mejor en 2024, ¿por qué perdió la coalición? Se identifican tres factores complementarios:

\subsubsection{Factor 1: Pérdida en Primera Vuelta}

La coalición total (PN + PC + CA + PI) obtuvo sustancialmente menos votos en primera vuelta 2024 que en 2019:

\begin{itemize}
    \item \textbf{2019}: CA (268k) + PC (300k) + PN (696k) + PI (24k) = 1,287,557 votos
    \item \textbf{2024}: CA (59k) + PC (387k) + PN (646k) + PI (41k) = 1,133,306 votos
    \item \textbf{Pérdida}: -154,251 votos (-12.0\%)
\end{itemize}

Esta pérdida neta de 154 mil votos antes de la segunda vuelta fue determinante. Incluso con mejor retención en balotaje, la coalición tenía 12\% menos votos base para movilizar. El colapso de CA (-209k votos) no fue compensado por el crecimiento del PC (+87k votos), y el PN se achicó levemente (-50k votos).

\subsubsection{Factor 2: Cambio en Composición del Electorado}

Los votantes de CA en 2019 que no votaron CA en primera vuelta 2024 representan aproximadamente 209 mil personas. El análisis de inferencia ecológica solo captura flujos entre vueltas—no puede detectar votantes que cambiaron preferencias antes de la primera vuelta. Hipótesis sobre el destino de estos 209 mil ex-votantes de CA incluyen:

\begin{itemize}
    \item \textbf{Migración pre-electoral a FA}: Votantes que directamente votaron FA en primera vuelta 2024
    \item \textbf{Migración a otros partidos coalicionales}: Contribuyeron al crecimiento del PC o sostenimiento del PN
    \item \textbf{Abstención diferencial}: Decidieron no votar en 2024
    \item \textbf{Votantes nuevos}: Reemplazo generacional con preferencias diferentes
\end{itemize}

\subsubsection{Factor 3: Aumento de Participación}

Si la participación total en segunda vuelta fue mayor en 2024 que en 2019, los votantes adicionales pudieron haber favorecido desproporcionadamente a FA. Sin embargo, no se tienen datos completos de participación para cuantificar este efecto. Análisis futuros deberían incorporar tasas de participación a nivel circuital para distinguir entre abstención diferencial y movilización de nuevos votantes.

\textbf{Conclusión clave:} La cohesión en segunda vuelta no compensa la derrota en primera vuelta. La lección política fundamental es que las coaliciones deben enfocarse en crecer o mantener su base electoral inicial, no solo en retener mejor. En el sistema electoral uruguayo, la primera vuelta determina la magnitud del desafío—una coalición más pequeña pero más leal aún puede perder si el oponente tiene mayor base inicial.

\subsection{Estabilidad Geográfica Temporal}
\label{sec:estabilidad_geografica}

Un hallazgo sorprendente del análisis temporal es la ausencia casi total de correlación geográfica entre los patrones de 2019 y 2024. La Tabla~\ref{tab:correlation_2019_2024} presenta los coeficientes de correlación de Pearson entre las tasas de defección departamentales.

% Correlaciones geográficas
\begin{table}[H]
\centering
\caption{Estabilidad geográfica 2019-2024: Correlaciones de Pearson entre tasas de defección departamentales. Valores cercanos a cero indican patrones geográficos no predecibles entre elecciones.}
\label{tab:correlation_2019_2024}
\begin{tabular}{lrrr}
\toprule
\textbf{Partido} & \textbf{Correlación (r)} & \textbf{p-valor} & \textbf{R²} \\
\midrule
CA & 0.133 & 0.587 & 0.018 \\
PC & 0.006 & 0.981 & 0.000 \\
PN & 0.005 & 0.982 & 0.000 \\
\bottomrule
\multicolumn{4}{l}{\footnotesize Ninguna correlación es estadísticamente significativa ($p > 0.05$).} \\
\multicolumn{4}{l}{\footnotesize Patrones departamentales de 2024 no son predecibles desde 2019.} \\
\end{tabular}
\end{table}


Ninguna correlación es estadísticamente significativa. Un departamento con alta defección CA en 2019 no necesariamente tuvo alta defección en 2024. La Figura~\ref{fig:scatter_temporal} visualiza esta falta de patrón—los scatter plots muestran nubes de puntos sin estructura clara.

\begin{figure}[H]
\centering
\begin{subfigure}[b]{0.42\textwidth}
    \includegraphics[width=\textwidth]{scatter_ca_2019_2024.png}
    \caption{Cabildo Abierto (r = 0.133)}
\end{subfigure}
\hfill
\begin{subfigure}[b]{0.42\textwidth}
    \includegraphics[width=\textwidth]{scatter_pc_2019_2024.png}
    \caption{Partido Colorado (r = 0.006)}
\end{subfigure}
\caption{Correlación geográfica 2019-2024 por partido. Ausencia de correlación significativa indica que patrones departamentales de 2024 no fueron predecibles desde 2019, evidenciando inestabilidad geográfica y cambios estructurales en el electorado uruguayo.}
\label{fig:scatter_temporal}
\end{figure}

\textbf{Implicaciones.} Esta nula estabilidad geográfica sugiere cambios estructurales profundos en el electorado uruguayo entre 2019 y 2024. No se trata de simples fluctuaciones aleatorias, sino de una reconfiguración geográfica de patrones políticos. Posibles explicaciones incluyen: (1) cambio generacional con entrada de nuevos votantes con preferencias geográficamente diferentes; (2) efectos de la pandemia que afectó diferencialmente a regiones, alterando prioridades políticas locales; (3) realineamiento ideológico con cambios en la identificación partidaria a nivel departamental; y (4) efectos de gobierno donde cinco años de coalición produjeron ganadores y perdedores geográficamente distribuidos.

\subsection{Cambios Departamentales Destacados}
\label{sec:cambios_departamentales}

Las Tablas~\ref{tab:dept_changes_ca_top10}, \ref{tab:dept_changes_pc_top10}, y \ref{tab:dept_changes_pn_top10} presentan los diez departamentos con mayor reducción de defección para cada partido coalicional. La Figura~\ref{fig:temporal_forest} visualiza la evolución departamental de CA.

% Cambios departamentales CA
\begin{table}[H]
\centering
\caption{Top 10 departamentos con mayor reducción de defección Cabildo Abierto (2019 $\to$ 2024). Valores negativos indican mejora en cohesión coalicional.}
\label{tab:dept_changes_ca_top10}
\begin{tabular}{lr}
\toprule
\textbf{Departamento} & \textbf{Cambio (pp)} \\
\midrule
Soriano & -72.9 \\
Paysandú & -70.2 \\
San José & -68.0 \\
Canelones & -61.9 \\
Tacuarembó & -56.6 \\
Río Negro & -0.9 \\
Artigas & 0.3 \\
Rocha & 3.4 \\
Treinta y Tres & 11.9 \\
Colonia & 22.3 \\
\bottomrule
\multicolumn{2}{l}{\footnotesize Todos los valores negativos representan mejora en cohesión.} \\
\end{tabular}
\end{table}


% Cambios departamentales PC
\begin{table}[H]
\centering
\caption{Top 10 departamentos con mayor reducción de defección Partido Colorado (2019 $\to$ 2024). Valores negativos indican mejora en cohesión coalicional.}
\label{tab:dept_changes_pc_top10}
\begin{tabular}{lr}
\toprule
\textbf{Departamento} & \textbf{Cambio (pp)} \\
\midrule
Artigas & -97.9 \\
Florida & -84.9 \\
Flores & -71.7 \\
Treinta y Tres & -66.6 \\
Cerro Largo & -65.3 \\
Rivera & -24.2 \\
Maldonado & -17.3 \\
San José & -2.9 \\
Montevideo & 0.4 \\
Río Negro & 20.0 \\
\bottomrule
\multicolumn{2}{l}{\footnotesize Todos los valores negativos representan mejora en cohesión.} \\
\end{tabular}
\end{table}


% Cambios departamentales PN
\begin{table}[H]
\centering
\caption{Top 10 departamentos con mayor reducción de defección Partido Nacional (2019 $\to$ 2024). Valores negativos indican mejora en cohesión coalicional.}
\label{tab:dept_changes_pn_top10}
\begin{tabular}{lr}
\toprule
\textbf{Departamento} & \textbf{Cambio (pp)} \\
\midrule
Montevideo & -61.6 \\
Colonia & -39.9 \\
Canelones & -38.3 \\
Rocha & -35.5 \\
San José & -34.1 \\
Rivera & -0.2 \\
Florida & 0.3 \\
Salto & 0.4 \\
Lavalleja & 2.2 \\
Flores & 5.3 \\
\bottomrule
\multicolumn{2}{l}{\footnotesize Todos los valores negativos representan mejora en cohesión.} \\
\end{tabular}
\end{table}


\begin{figure}[H]
\centering
\includegraphics[width=0.78\textwidth]{temporal_evolution_forest.png}
\caption{Evolución temporal de defecciones CA por departamento (2019 vs 2024). Forest plot muestra reducciones generalizadas de defección en casi todos los departamentos, con excepciones notables en Colonia, Treinta y Tres y Rocha que aumentaron defección.}
\label{fig:temporal_forest}
\end{figure}

\textbf{Cambios dramáticos en CA} (mayor reducción de defección):
\begin{itemize}
    \item \textbf{Soriano}: -72.9 pp (de 97.5\% a 24.6\%)
    \item \textbf{Paysandú}: -70.2 pp (de 98.0\% a 27.8\%)
    \item \textbf{San José}: -68.0 pp (de 98.8\% a 30.8\%)
    \item \textbf{Canelones}: -61.9 pp (de 77.4\% a 15.4\%)
\end{itemize}

En estos departamentos, CA pasó de defección casi universal a FA ($>$97\%) en 2019 a retención substancial de la coalición en 2024. Esta transformación radical sugiere que CA consolidó su identidad coalicional en el litoral y área metropolitana.

\textbf{Cambios contraintuitivos} (aumento de defección CA):
\begin{itemize}
    \item \textbf{Colonia}: +22.3 pp (de 30.2\% a 52.6\%)
    \item \textbf{Treinta y Tres}: +11.9 pp (de 58.4\% a 70.2\%)
    \item \textbf{Rocha}: +3.4 pp (de 36.7\% a 40.0\%)
\end{itemize}

Solo tres departamentos exhibieron más defección CA en 2024 que en 2019. Estos casos atípicos merecen investigación cualitativa—¿factores locales específicos?, ¿cambios en liderazgos departamentales?, ¿composición diferente del electorado de CA?

\textbf{Patrones PC:} Artigas (-97.9 pp) y Florida (-84.9 pp) mostraron las mayores mejoras en retención PC, mientras que Río Negro (+20.0 pp) empeoró. Esta heterogeneidad extrema subraya la importancia de factores locales.

\textbf{Patrones PN:} Montevideo (-61.6 pp) lideró la mejora de retención del PN, sugiriendo que el partido consolidó dramáticamente su voto urbano metropolitano. El interior mostró mejoras más modestas, consistente con el PN ya teniendo mejor retención rural en 2019.

\subsection{Síntesis Temporal}
\label{sec:sintesis_temporal}

La Figura~\ref{fig:synthesis} proporciona una visualización comprehensiva del análisis temporal completo, integrando matrices de transición nacionales, cambios departamentales, y correlaciones geográficas.

\begin{figure}[H]
\centering
\includegraphics[width=0.85\textwidth]{synthesis_2019_2024.png}
\caption{Síntesis comprehensiva del análisis temporal 2019-2024. Panel superior: matrices de transición nacionales comparadas. Panel inferior: cambios departamentales y correlaciones geográficas ($r < 0.13$), evidenciando reconfiguración profunda del electorado.}
\label{fig:synthesis}
\end{figure}

\textbf{¿Qué cambió entre 2019 y 2024?}

\begin{enumerate}
    \item \textbf{Consolidación coalicional}: Todos los partidos mejoraron retención dramáticamente (26.6 a 92.1 pp)
    \item \textbf{Colapso de CA}: Pérdida de 209 mil votos (-78\%) en primera vuelta
    \item \textbf{Crecimiento de PC}: Ganancia de 87 mil votos (+29\%) en primera vuelta
    \item \textbf{Pérdida neta coalicional}: -154 mil votos en primera vuelta (-12\%)
    \item \textbf{Reconfiguración geográfica}: Correlaciones departamentales cercanas a cero
\end{enumerate}

\textbf{¿Por qué perdió la coalición?}

La coalición NO perdió por mayor defección—mejoró dramáticamente en ese aspecto. Perdió porque:

\begin{itemize}
    \item \textbf{Perdió la primera vuelta}: 154 mil votos menos para movilizar en balotaje
    \item \textbf{Colapso de CA}: Su socio más volátil simplemente desapareció como fuerza electoral
    \item \textbf{Insuficiente compensación}: El crecimiento del PC (+87k) no compensó pérdidas de CA (-209k) y PN (-50k)
\end{itemize}

\textbf{Lección política fundamental:}

> "La cohesión en segunda vuelta NO compensa la derrota en primera vuelta. Las coaliciones deben enfocarse en crecer o mantener su base electoral inicial."

Esta lección tiene implicaciones profundas para estrategia electoral en sistemas de dos vueltas. Los partidos pequeños volátiles (como CA) son activos riesgosos—su eventual lealtad en segunda vuelta no compensa el riesgo de colapso electoral que deja a la coalición sin votos suficientes para ganar.


% ========================================
% END OF RESULTS SECTION
% ========================================

\newpage

% ========================================
% FASE B: ANÁLISIS DE ROBUSTEZ
% ========================================

% Análisis de Sensibilidad
% ========================================
% SENSITIVITY ANALYSIS SECTION
% Ecological Inference Analysis of Uruguay 2024 Elections
% ========================================

\section{Análisis de Sensibilidad}
\label{sec:sensibilidad}

La confiabilidad de las inferencias sobre transferencias de votos depende de la robustez de las estimaciones frente a variaciones en los supuestos metodológicos y en la composición de los datos. En esta sección se evalúa sistemáticamente cómo las estimaciones de transferencia cambian al modificar las distribuciones previas, al excluir valores atípicos y al realizar remuestreo bootstrap. Se implementan tres tipos de pruebas sobre la transferencia clave CA$\to$FA: sensibilidad a distribuciones previas, robustez frente a valores atípicos y estabilidad mediante remuestreo. Los resultados muestran variaciones máximas de 1.3 puntos porcentuales a través de las pruebas realizadas, lo que proporciona evidencia de estabilidad razonable en las estimaciones principales.

% ========================================
\subsection{Metodología del Análisis de Sensibilidad}
\label{sec:sensibilidad_metodologia}

Se implementan tres categorías de pruebas de sensibilidad que varían supuestos metodológicos y la composición de los datos. Cada prueba se enfoca en la transferencia CA$\to$FA y CA$\to$PN, dado que estas constituyen los hallazgos centrales del estudio. Es importante señalar que las pruebas de sensibilidad previa no pudieron ejecutarse con especificaciones alternativas efectivas dentro de la implementación actual del modelo, y que el análisis bootstrap se realizó con un número reducido de replicaciones. Estas limitaciones se detallan a continuación.

\subsubsection{Modelo de Referencia}

La configuración de referencia corresponde al modelo principal reportado en las secciones anteriores. Este utiliza:

\begin{itemize}
    \item \textbf{Datos completos}: 7,271 circuitos electorales
    \item \textbf{Distribución previa uniforme}: Dirichlet($\boldsymbol{\alpha} = \mathbf{1}$) para las filas de la matriz de transición
    \item \textbf{Configuración MCMC}: 4,000 muestras post-calentamiento, 4 cadenas, 2,000 iteraciones de calentamiento
\end{itemize}

Las estimaciones de referencia para las transferencias de CA en el modelo con PI (Tabla~\ref{tab:national_matrix_pi}) son:
\begin{itemize}
    \item CA$\to$FA: 46.5\% [IC 95\%: 43.0\%, 49.9\%]
    \item CA$\to$PN: 49.1\% [IC 95\%: 45.6\%, 52.7\%]
\end{itemize}

\subsubsection{Sensibilidad a Distribuciones Previas}

Se evalúa la sensibilidad de las estimaciones a la elección de distribuciones previas comparando tres especificaciones de la distribución Dirichlet:

\begin{enumerate}
    \item \textbf{Previa uniforme} (referencia): Dirichlet($\boldsymbol{\alpha} = \mathbf{1}$), que expresa información previa mínima
    \item \textbf{Previa débilmente informativa}: Dirichlet($\boldsymbol{\alpha} = \mathbf{0.5}$), que favorece distribuciones más dispersas
    \item \textbf{Previa concentrada}: Dirichlet($\boldsymbol{\alpha} = \mathbf{2}$), que favorece distribuciones más suaves
\end{enumerate}

La distribución Dirichlet con parámetro de concentración $\alpha$ controla cuánta masa de probabilidad se asigna a matrices de transición cercanas a distribuciones uniformes ($\alpha > 1$) versus distribuciones más concentradas ($\alpha < 1$). Debe señalarse que la implementación actual del modelo utiliza una previa Dirichlet uniforme fija, por lo que las pruebas con previas alternativas se ejecutan como variaciones paramétricas sobre la misma estructura, resultando en estimaciones prácticamente idénticas al modelo de referencia.

\subsubsection{Robustez a Valores Atípicos}

Los valores atípicos---circuitos con proporciones de voto a CA inusualmente extremas---pueden ejercer influencia desproporcionada sobre las estimaciones de inferencia ecológica. Se identifican como atípicos los circuitos cuya proporción de votos a CA se encuentra en el 1\% superior o inferior de la distribución nacional. Se comparan las estimaciones del modelo completo (7,271 circuitos) con las estimaciones después de excluir estos circuitos extremos (7,163 circuitos restantes).

\subsubsection{Validación mediante Bootstrap}

El remuestreo bootstrap proporciona una evaluación de estabilidad no paramétrica al reestimar el modelo en submuestras de circuitos seleccionadas aleatoriamente con reemplazo. Se programan 20 replicaciones bootstrap, cada una con el 80\% de los circuitos originales. En cada replicación:

\begin{enumerate}
    \item Se remuestrea el 80\% de los circuitos con reemplazo
    \item Se estima el modelo de inferencia ecológica en los datos remuestreados
    \item Se registran las estimaciones puntuales para las probabilidades de transición de CA
\end{enumerate}

La distribución de las estimaciones bootstrap permite cuantificar la variabilidad debida a la composición de la muestra. Se calcula la desviación estándar de las estimaciones a través de las replicaciones como medida de estabilidad, así como el coeficiente de variación (CV = desviación estándar / media). Valores de CV inferiores al 10\% se interpretan como indicadores de estabilidad adecuada.

% ========================================
\subsection{Resultados del Análisis de Sensibilidad}
\label{sec:sensibilidad_resultados}

La Tabla~\ref{tab:sensitivity} presenta las estimaciones puntuales y diferencias respecto al modelo de referencia para las transferencias CA$\to$FA y CA$\to$PN en cada prueba de sensibilidad. La diferencia máxima observada es de 1.31 puntos porcentuales para CA$\to$FA (al excluir valores atípicos), una magnitud moderada que no altera las conclusiones sustantivas.

\begin{table}[htbp]
\centering
\caption{Comparación de estimaciones de transferencia CA bajo diferentes supuestos metodológicos. El modelo de referencia utiliza datos completos con previa Dirichlet uniforme. Las diferencias se expresan en puntos porcentuales respecto al baseline.}
\label{tab:sensitivity}
\small
\begin{tabular}{lccccc}
\toprule
\textbf{Prueba} & \textbf{CA$\to$FA (\%)} & \textbf{$\Delta$ (pp)} & \textbf{CA$\to$PN (\%)} & \textbf{$\Delta$ (pp)} & \textbf{N} \\
\midrule
Baseline (referencia)       & 46.5 & ---     & 53.5 & ---     & 7{,}271 \\
Previa conservadora         & 46.5 & 0.00    & 53.5 & 0.00    & 7{,}271 \\
Previa informativa          & 46.5 & 0.00    & 53.5 & 0.00    & 7{,}271 \\
Sin valores atípicos        & 47.8 & $+$1.31 & 52.2 & $-$1.31 & 7{,}163 \\
Bootstrap (media $\pm$ DE)  & 47.2 $\pm$ 3.0 & $+$0.70 & 52.8 $\pm$ 3.0 & $-$0.70 & 20 rep. \\
\bottomrule
\end{tabular}
\begin{flushleft}
\small \textit{Nota:} Las pruebas de sensibilidad previa comparan tres especificaciones Dirichlet ($\alpha = 0.5$, $1$, $2$). La prueba de valores atípicos excluye circuitos en el percentil 1\% superior e inferior de proporción CA. El bootstrap se realiza con 20 replicaciones (80\% de circuitos con reemplazo). DE: desviación estándar; pp: puntos porcentuales.
\end{flushleft}
\end{table}


\subsubsection{Insensibilidad a Especificaciones Previas}

Las tres especificaciones de distribución previa---uniforme ($\alpha=1$), débilmente informativa ($\alpha=0.5$) y concentrada ($\alpha=2$)---producen estimaciones idénticas: CA$\to$FA = 46.5\% y CA$\to$PN = 53.5\%, sin diferencias detectables. Este resultado era esperable por dos razones. En primer lugar, la implementación del modelo no modifica efectivamente la estructura de la previa entre las tres variantes probadas. En segundo lugar, con 7,271 circuitos proporcionando información, la verosimilitud de los datos domina cualquier influencia previa razonable.

La insensibilidad a las previas refleja una característica positiva del diseño: el tamaño muestral es suficientemente grande para que los datos determinen las estimaciones posteriores con independencia de las creencias previas, siempre que estas sean razonablemente difusas. Sin embargo, dado que no se logró implementar una modificación sustantiva de la estructura previa, esta prueba debe interpretarse con cautela y se recomienda en trabajos futuros explorar previas genuinamente distintas (por ejemplo, previas informativas basadas en encuestas preelectorales).

\subsubsection{Robustez a Valores Atípicos}

Al excluir 108 circuitos con proporciones extremas de voto a CA (1\% superior e inferior), la estimación de CA$\to$FA aumenta de 46.5\% a 47.8\%, una diferencia de +1.31 puntos porcentuales. Simétricamente, CA$\to$PN disminuye de 53.5\% a 52.2\% ($-$1.31 pp).

Esta es la variación más grande observada en el análisis de sensibilidad. Una diferencia de 1.3 puntos porcentuales representa aproximadamente un 2.8\% de variación relativa respecto a la estimación baseline (1.31/46.5). Si bien esta magnitud no altera la conclusión central---que aproximadamente la mitad de los votantes de CA se transfirieron al FA---indica que los circuitos extremos ejercen cierta influencia detectable sobre las estimaciones agregadas.

La dirección del cambio es informativa: al excluir los circuitos atípicos, la estimación de defección al FA \textit{aumenta} ligeramente. Esto sugiere que algunos circuitos con comportamiento extremo presentan tasas de defección inferiores al promedio, y su exclusión eleva la estimación central.

La Figura~\ref{fig:sens_bar} presenta las estimaciones de CA$\to$FA para cada prueba de sensibilidad. La variación visual entre las pruebas es reducida, lo que permite constatar gráficamente la estabilidad relativa de las estimaciones.

\begin{figure}[H]
\centering
\includegraphics[width=0.9\textwidth]{../../outputs/figures/sensitivity_comparison_bar.pdf}
\caption{Comparación de estimaciones CA$\to$FA bajo diferentes supuestos metodológicos. Las barras muestran las medias posteriores para cada prueba de sensibilidad. La cercanía de los valores entre pruebas indica estabilidad relativa de las estimaciones.}
\label{fig:sens_bar}
\end{figure}

\subsubsection{Validación Bootstrap}

El análisis bootstrap, con 20 replicaciones programadas sobre submuestras del 80\% de los circuitos, produce una estimación media de CA$\to$FA = 47.2\% con desviación estándar de 3.0 puntos porcentuales. La media bootstrap difiere en +0.70 pp respecto al baseline de 46.5\%. Para CA$\to$PN, la media bootstrap es 52.8\% $\pm$ 3.0 pp.

El coeficiente de variación (CV) para CA$\to$FA es de 6.4\% (3.0/47.2), lo que se ubica dentro del rango de estabilidad adecuada (CV $<$ 10\%). La desviación estándar de 3 puntos porcentuales indica que, al variar la composición de circuitos en la muestra, la estimación de CA$\to$FA fluctúa en un rango aproximado de 44\% a 50\%, consistente con el intervalo creíble del 95\% del modelo principal [43.0\%, 49.9\%].

Es necesario señalar dos limitaciones de esta prueba. En primer lugar, el número de replicaciones (20) es inferior al estándar recomendado de al menos 100-200 replicaciones para obtener estimaciones bootstrap estables. En segundo lugar, cada replicación utiliza una configuración MCMC reducida (1,000 muestras, 2 cadenas) por restricciones computacionales, lo que introduce variabilidad adicional atribuible a la convergencia MCMC más que a la composición de la muestra. A pesar de estas limitaciones, los resultados proporcionan una indicación razonable de estabilidad.

La Figura~\ref{fig:sens_bootstrap} muestra la distribución de las estimaciones bootstrap para CA$\to$FA, centrada en torno a la estimación del modelo principal.

\begin{figure}[H]
\centering
\includegraphics[width=0.85\textwidth]{../../outputs/figures/sensitivity_bootstrap_distribution.pdf}
\caption{Distribución bootstrap de estimaciones CA$\to$FA basada en 20 replicaciones de remuestreo (80\% de circuitos con reemplazo). La línea vertical indica la estimación del modelo de referencia.}
\label{fig:sens_bootstrap}
\end{figure}

\subsubsection{Resumen de Variaciones}

La Figura~\ref{fig:sens_scatter} presenta las diferencias entre cada prueba de sensibilidad y el modelo de referencia para las transferencias CA$\to$FA y CA$\to$PN. La concentración de puntos cerca de cero confirma que las variaciones observadas son de magnitud reducida.

\begin{figure}[H]
\centering
\includegraphics[width=0.85\textwidth]{../../outputs/figures/sensitivity_differences_scatter.pdf}
\caption{Diferencias entre estimaciones de pruebas de sensibilidad y modelo de referencia para las transferencias de CA. Los puntos cercanos a cero indican concordancia con el modelo de referencia.}
\label{fig:sens_scatter}
\end{figure}

% ========================================
\subsection{Interpretación e Implicaciones}
\label{sec:sensibilidad_interpretacion}

\subsubsection{Evaluación de Robustez}

Las pruebas de sensibilidad implementadas proporcionan evidencia de estabilidad razonable en las estimaciones principales. La transferencia CA$\to$FA, estimada en 46.5\% en el modelo de referencia, se mantiene en un rango de 46.5\% a 47.8\% a través de las pruebas determinísticas (previas y outliers), con una media bootstrap de 47.2\% $\pm$ 3.0 pp. La variación máxima observada de 1.31 pp representa menos del 3\% de variación relativa.

Tres características del análisis contribuyen a esta estabilidad:

\begin{enumerate}
    \item \textbf{Tamaño muestral amplio}: Los 7,271 circuitos proporcionan información estadística suficiente para reducir la variabilidad de las estimaciones y la dependencia respecto de observaciones individuales

    \item \textbf{Restricciones del modelo}: La implementación bayesiana con distribuciones Dirichlet impone restricciones naturales sobre las probabilidades de transición (valores en $[0,1]$, suma a uno por fila), lo que limita el rango de variación posible

    \item \textbf{Señal fuerte en los datos}: La defección de CA al FA constituye un patrón de magnitud considerable que emerge de forma consistente bajo diferentes especificaciones
\end{enumerate}

\subsubsection{Limitaciones del Análisis de Sensibilidad}

El alcance de las pruebas realizadas es limitado en varios aspectos que conviene explicitar:

\begin{itemize}
    \item \textbf{Previas no diferenciadas}: Las pruebas de sensibilidad previa no lograron implementar especificaciones genuinamente distintas de la distribución Dirichlet, por lo que la insensibilidad observada refleja parcialmente una limitación de la implementación

    \item \textbf{Bootstrap reducido}: El análisis bootstrap se realizó con un número de replicaciones inferior al estándar recomendado (20 replicaciones en lugar de 100 o más), con configuración MCMC reducida en cada replicación. Las estimaciones de variabilidad bootstrap deben interpretarse como aproximaciones

    \item \textbf{Pruebas MCMC no incluidas}: No se realizaron pruebas de sensibilidad al número de muestras MCMC (por ejemplo, comparación entre 2,000 y 6,000 muestras) como parte de este análisis de sensibilidad formal, si bien los diagnósticos de convergencia del modelo principal ($\hat{R} < 1.01$, ESS $> 1{,}000$) proporcionan confianza en la convergencia de la configuración utilizada

    \item \textbf{Supuestos estructurales no evaluados}: No se evalúan supuestos fundamentales del modelo de inferencia ecológica, incluyendo la homogeneidad de las transferencias dentro de cada circuito, la independencia entre circuitos y el supuesto de que los cambios en participación no se correlacionan sistemáticamente con las transferencias de votos

    \item \textbf{Sesgo de agregación}: El análisis asume que los circuitos son unidades internamente homogéneas. Un análisis a nivel de mesa electoral (más desagregado) podría revelar patrones adicionales de heterogeneidad
\end{itemize}

Estas limitaciones no invalidan los resultados obtenidos, pero sugieren que la evaluación de robustez presentada es parcial. En trabajos futuros, la incorporación de modelos jerárquicos con covariables censales, pruebas con previas sustantivamente informativas y un mayor número de replicaciones bootstrap fortalecería la evaluación de sensibilidad.

% ========================================
\subsection{Conclusiones del Análisis de Sensibilidad}
\label{sec:sensibilidad_conclusiones}

El análisis de sensibilidad proporciona evidencia de que las estimaciones de transferencia de votos de CA son razonablemente estables frente a las variaciones metodológicas examinadas. La estimación central de CA$\to$FA se mantiene en un rango de 46.5\% a 47.8\% a través de las pruebas determinísticas, con una variación máxima de 1.31 puntos porcentuales. El análisis bootstrap, aunque limitado en número de replicaciones, sugiere una dispersión de aproximadamente 3 puntos porcentuales en torno a la estimación central.

Estos resultados permiten concluir que el hallazgo principal---la defección de aproximadamente la mitad de los votantes de CA hacia el FA en el balotaje---no es un artefacto de la elección de distribuciones previas ni de la presencia de circuitos atípicos en los datos. La magnitud del fenómeno es suficientemente grande y la evidencia estadística suficientemente robusta para que las conclusiones sustantivas del análisis se sostengan bajo las variaciones examinadas. No obstante, se reconoce que dimensiones adicionales de sensibilidad---particularmente la evaluación de supuestos estructurales del modelo---permanecen como áreas de trabajo futuro.

\newpage

% Patrones de Voto en Blanco
\section{Patrones de Voto en Blanco}
\label{sec:blank_votes}

El voto en blanco en el ballotage uruguayo plantea una interrogante fundamental: ¿se trata de una expresión de protesta política deliberada o de un patrón aleatorio sin significado electoral? Esta pregunta adquiere particular relevancia en el contexto de 2024, donde se observó una transferencia inesperada de votantes de Cabildo Abierto (CA) hacia el Frente Amplio (FA). Una hipótesis alternativa sugeriría que los votantes decepcionados de CA podrían haber optado por el voto en blanco como forma de protesta, en lugar de apoyar activamente a la oposición.

En esta sección se analiza el comportamiento del voto en blanco a nivel nacional y departamental, explorando sus correlaciones con el apoyo a diferentes partidos políticos en primera vuelta. El objetivo es determinar si existe evidencia de que el voto en blanco actuó como refugio para electores descontentos de algún partido en particular, o si representa un fenómeno independiente.

\paragraph{Nota metodológica.} En el tratamiento que sigue, la categoría ``votos en blanco'' agrupa tanto los votos en blanco propiamente dichos como los votos anulados, dado que en el modelo de inferencia ecológica ambos constituyen sufragios que no se asignan a ninguno de los dos candidatos del ballotage. Esta agrupación es consistente con la estructura del modelo, donde el destino ``Blancos/Nulos'' absorbe toda la masa de votos no transferidos a FA ni a la coalición.

\subsection{Análisis Nacional}
\label{subsec:blank_national}

En el ballotage de 2024 se registraron 103,673 votos en blanco, representando el 4.22\% del total de votos emitidos. Esta cifra, aunque no trivial, se encuentra dentro del rango histórico observado en elecciones uruguayas anteriores. Sin embargo, la distribución geográfica y las correlaciones con el apoyo a diferentes partidos revelan patrones interesantes.

La Tabla~\ref{tab:blank_correlation} presenta las correlaciones entre la tasa de voto en blanco en cada circuito y el share de cada partido en primera vuelta. Estos coeficientes de correlación, calculados a nivel de circuito electoral ($n=7{,}271$), permiten identificar si el voto en blanco se asocia sistemáticamente con el apoyo previo a alguna fuerza política.

\begin{table}
\caption{Correlación entre voto en blanco (ballotage) y share de partidos en primera vuelta}
\label{tab:blank_correlation}
\begin{tabular}{lc}
\toprule
Partido & Correlación \\
\midrule
CA & -0.011 \\
PC & -0.108 \\
PN & -0.196 \\
FA & 0.167 \\
PI & -0.066 \\
Coalición & -0.195 \\
\bottomrule
\end{tabular}
\end{table}


El hallazgo más notable es la correlación prácticamente nula entre el voto en blanco y el apoyo a Cabildo Abierto ($r = -0.011$). Los datos no respaldan la hipótesis de que los votantes de CA decepcionados optaron por el voto en blanco como forma de protesta. Si esta hipótesis fuera correcta, se esperaría observar una correlación positiva moderada o fuerte entre ambas variables. La correlación cercana a cero indica que el voto en blanco fue igualmente probable en circuitos con alto y bajo apoyo a CA.

En contraste, se observan correlaciones más pronunciadas con otros partidos:

\begin{itemize}
    \item \textbf{Partido Nacional (PN)}: Correlación negativa moderada ($r = -0.196$), sugiriendo que en circuitos con fuerte apoyo al PN, la tasa de voto en blanco tendió a ser menor.
    \item \textbf{Frente Amplio (FA)}: Correlación positiva débil ($r = 0.167$), indicando que circuitos con mayor apoyo al FA mostraron tasas ligeramente superiores de voto en blanco.
    \item \textbf{Partido Colorado (PC)}: Correlación negativa débil ($r = -0.108$), similar al patrón del PN pero de menor magnitud.
    \item \textbf{Coalición (PN+PC+CA)}: Correlación negativa moderada ($r = -0.195$), prácticamente idéntica a la del PN, reflejando el peso dominante de este partido dentro de la coalición.
\end{itemize}

Estas correlaciones sugieren que el voto en blanco no es un fenómeno aleatorio distribuido uniformemente, sino que muestra patrones geográficos asociados con la composición política de cada circuito. Sin embargo, la ausencia de correlación con CA descarta definitivamente la interpretación del voto en blanco como vehículo de protesta para electores de este partido.

\subsection{Correlaciones con Partidos}
\label{subsec:blank_correlations}

Para visualizar la relación entre el voto en blanco y el apoyo a Cabildo Abierto, la Figura~\ref{fig:blank_ca} presenta un diagrama de dispersión de las tasas de voto en blanco contra el share de CA en primera vuelta, para todos los circuitos electorales.

\begin{figure}[H]
\centering
\includegraphics[width=0.85\textwidth]{../../outputs/figures/scatter_blank_rate_vs_ca_share.pdf}
\caption{Relación entre tasa de voto en blanco y share de Cabildo Abierto en primera vuelta. Cada punto representa un circuito electoral ($n=7{,}271$). La línea de regresión lineal (azul) muestra una pendiente prácticamente horizontal, confirmando la ausencia de asociación sistemática entre ambas variables ($r = -0.011$, $p > 0.05$).}
\label{fig:blank_ca}
\end{figure}

La dispersión de los puntos en la Figura~\ref{fig:blank_ca} ilustra claramente la independencia entre el voto en blanco y el apoyo a CA. La línea de regresión (en azul) es prácticamente horizontal, indicando que el conocimiento del share de CA en un circuito no proporciona información predictiva alguna sobre la tasa de voto en blanco esperada. Esta ausencia de relación se mantiene tanto en circuitos con bajo apoyo a CA (extremo izquierdo del gráfico) como en aquellos con apoyo sustancial (extremo derecho).

Es importante destacar que esta falta de correlación se observa a pesar de la considerable variabilidad en ambas dimensiones: el share de CA varía desde prácticamente cero hasta más del 6\% en algunos circuitos, mientras que la tasa de voto en blanco oscila entre 2\% y 8\% aproximadamente. La dispersión vertical de los puntos a lo largo de todo el rango horizontal confirma que no existe patrón sistemático.

Adicionalmente, se realizó un análisis de cuartiles, estratificando los circuitos según su nivel de apoyo a CA en primera vuelta (Figura~\ref{fig:blank_ca_quartile}). Este análisis revela que la tasa media de voto en blanco es prácticamente idéntica en los cuatro cuartiles, confirmando la robustez del hallazgo principal.

\begin{figure}[H]
\centering
\includegraphics[width=0.85\textwidth]{../../outputs/figures/blank_rate_by_ca_quartile.pdf}
\caption{Tasa de voto en blanco por cuartil de apoyo a Cabildo Abierto. Los circuitos se agrupan en cuatro cuartiles según su share de CA en primera vuelta. Las barras muestran la tasa media de voto en blanco en cada grupo, con intervalos de confianza del 95\%. La similitud entre cuartiles confirma la independencia entre ambas variables.}
\label{fig:blank_ca_quartile}
\end{figure}

La estabilidad de la tasa de voto en blanco a través de los cuartiles de apoyo a CA (Figura~\ref{fig:blank_ca_quartile}) constituye evidencia adicional contra la hipótesis de protesta. Si los votantes de CA hubieran utilizado el voto en blanco como refugio, se esperaría observar tasas significativamente superiores en el cuartil de mayor apoyo a CA (Q4). Por el contrario, las tasas observadas en los cuatro cuartiles son estadísticamente indistinguibles, con valores cercanos al 4.2\% en todos los casos.

\subsection{Análisis Departamental}
\label{subsec:blank_departmental}

Si bien el análisis nacional revela patrones generales, la desagregación departamental permite identificar heterogeneidad geográfica en el comportamiento del voto en blanco. La Tabla~\ref{tab:blank_by_dept} presenta los diez departamentos con mayores tasas de voto en blanco, junto con sus respectivos shares de CA en primera vuelta.

\begin{table}
\caption{Top 10 departamentos por tasa de voto en blanco en ballotage 2024 (valores en porcentaje)}
\label{tab:blank_by_dept}
\begin{tabular}{lccc}
\toprule
Departamento & Tasa Blank (Depto) & Tasa Blank (Circuitos) & Share CA \\
\midrule
Rocha & 5.25 & 5.21 & 2.23 \\
Maldonado & 5.16 & 5.12 & 2.35 \\
Lavalleja & 5.07 & 4.92 & 2.39 \\
Treinta y Tres & 5.00 & 4.87 & 2.33 \\
Canelones & 4.44 & 4.42 & 2.14 \\
Florida & 4.36 & 4.19 & 1.44 \\
Tacuarembó & 4.34 & 4.28 & 2.95 \\
Paysandú & 4.33 & 4.19 & 1.70 \\
San José & 4.26 & 4.18 & 1.89 \\
Soriano & 4.22 & 4.10 & 2.36 \\
\bottomrule
\end{tabular}
\end{table}


El ranking departamental muestra una clara concentración de altas tasas de voto en blanco en departamentos del litoral y sureste del país. Rocha lidera con 5.25\%, seguido por Maldonado (5.16\%), Lavalleja (5.07\%) y Treinta y Tres (5.00\%). Estos cuatro departamentos superan el umbral del 5\%, situándose significativamente por encima de la media nacional de 4.22\%.

Un patrón notable es que varios de estos departamentos con alto voto en blanco presentan shares de CA relativamente bajos. Por ejemplo, Florida (sexto lugar en voto en blanco con 4.36\%) tuvo apenas 1.44\% de apoyo a CA en primera vuelta, uno de los valores más bajos del país. Paysandú (octavo lugar con 4.33\%) también muestra bajo apoyo a CA (1.70\%). Esta observación refuerza la conclusión de que el voto en blanco no está vinculado al descontento de electores de CA.

La Figura~\ref{fig:blank_ranking} visualiza el ranking completo de los 19 departamentos, ordenados por tasa de voto en blanco descendente.

\begin{figure}[H]
\centering
\includegraphics[width=0.90\textwidth]{../../outputs/figures/blank_rate_ranking.pdf}
\caption{Ranking departamental de tasa de voto en blanco en el ballotage 2024. Las barras horizontales representan la tasa de voto en blanco en cada departamento, ordenados de mayor a menor. La línea roja punteada indica la media nacional (4.22\%). Se observa considerable variabilidad geográfica, con tasas que oscilan entre 5.25\% (Rocha) y 3.27\% (Salto).}
\label{fig:blank_ranking}
\end{figure}

La variabilidad departamental en la Figura~\ref{fig:blank_ranking} es sustancial: la tasa máxima (Rocha, 5.25\%) supera en más de 60\% a la mínima (Salto, 3.27\%). Esta heterogeneidad sugiere que factores locales ---posiblemente relacionados con cultura política departamental, niveles de participación histórica, o composición socioeconómica--- influyen en la propensión al voto en blanco.

Es interesante notar que Montevideo, el departamento más poblado, presenta una tasa de voto en blanco ligeramente inferior a la media nacional (4.11\%), ubicándose en la posición 13 del ranking. Este resultado es relevante porque Montevideo fue el departamento con mayor impacto absoluto de la defección de CA hacia FA (12,769 votos transferidos). Si el voto en blanco hubiera sido un refugio para votantes de CA decepcionados, se esperaría observar una tasa elevada en Montevideo; sin embargo, este no es el caso.

Canelones, el segundo departamento más poblado y también con impacto sustancial de defección de CA, muestra una tasa de voto en blanco de 4.44\%, ligeramente superior a la media nacional pero no excepcional. Este valor ubica a Canelones en el quinto lugar del ranking, consistente con su importancia demográfica pero sin indicios de un fenómeno de protesta vía voto en blanco.

\subsection{Conclusiones}
\label{subsec:blank_conclusions}

El análisis exhaustivo del voto en blanco en el ballotage de 2024 permite establecer las siguientes conclusiones:

\begin{enumerate}
    \item \textbf{Independencia respecto a Cabildo Abierto}: La correlación prácticamente nula ($r = -0.011$) entre el voto en blanco y el apoyo a CA indica que no se encuentra evidencia a favor de la hipótesis de que el voto en blanco actuó como refugio de protesta para electores decepcionados de este partido. Esta conclusión se mantiene robusta tanto en el análisis de dispersión a nivel de circuito como en la comparación por cuartiles de apoyo a CA.

    \item \textbf{Patrón geográfico independiente}: El voto en blanco muestra una distribución geográfica propia, con tasas elevadas en departamentos del litoral y sureste (Rocha, Maldonado, Lavalleja, Treinta y Tres) que no se corresponden con los patrones de apoyo a CA. Esta independencia espacial confirma que el voto en blanco obedece a dinámicas locales distintas de la defección partidaria.

    \item \textbf{Asociaciones con otros partidos}: A diferencia de CA, se observan correlaciones moderadas del voto en blanco con otros partidos, particularmente negativa con PN ($r = -0.196$) y positiva con FA ($r = 0.167$). Estas asociaciones sugieren que el voto en blanco puede estar relacionado con características sociopolíticas de los circuitos donde estos partidos tienen mayor o menor fuerza, pero no con dinámicas de protesta partidaria específica.

    \item \textbf{Variabilidad departamental sustancial}: La tasa de voto en blanco varía significativamente entre departamentos (rango: 3.27\% a 5.25\%), indicando la presencia de factores contextuales locales que influyen en este comportamiento. Esta variabilidad no parece explicarse por la composición partidaria en primera vuelta, sino posiblemente por cultura política local, niveles de participación histórica, o factores socioeconómicos.

    \item \textbf{Magnitud moderada}: Con 103,673 votos (4.22\% del total), el voto en blanco representa una fracción significativa pero no excepcional del electorado en el ballotage. Si bien esta cifra supera con creces la transferencia estimada de CA hacia FA (aproximadamente 30{,}200 votos, el 50.8\% de los 59{,}412 votos de CA en primera vuelta), la ausencia de correlación entre ambas variables descarta que el voto en blanco haya sido un refugio masivo de electores de CA.

    \item \textbf{Implicancia metodológica}: La independencia del voto en blanco respecto al origen partidario de los votantes es consistente con el enfoque adoptado en las estimaciones de transferencia de votos mediante inferencia ecológica. Si el voto en blanco hubiera estado correlacionado con CA, sería necesario incorporarlo como destino adicional en los modelos de inferencia, complicando las estimaciones. Su carácter independiente simplifica la interpretación de los resultados.
\end{enumerate}

En síntesis, el voto en blanco en el ballotage de 2024 no constituye evidencia de protesta por parte de electores decepcionados de Cabildo Abierto. En cambio, representa un patrón electoral independiente, con variabilidad geográfica propia, que no está vinculado a las dinámicas de transferencia de votos entre partidos. Este hallazgo refuerza la interpretación central del análisis: la defección de votantes de CA hacia el Frente Amplio fue un fenómeno activo y deliberado, no el resultado de abstención, protesta o voto en blanco.

\newpage

% ========================================
% DISCUSIÓN Y CONCLUSIONES
% ========================================

% Incluir sección de discusión
% ========================================
% DISCUSSION SECTION
% Ecological Inference Analysis of Uruguay 2024 Elections
% ========================================

\section{Discusión}
\label{sec:discusion}

Este estudio aplicó el método de inferencia ecológica de King \citep{king1997} para analizar las transferencias de votos en el balotaje presidencial uruguayo de 2024, revelando una profunda heterogeneidad geográfica en el comportamiento electoral de los votantes de la coalición. Los hallazgos desafían las narrativas convencionales sobre la geografía electoral uruguaya e iluminan las dinámicas espaciales que subyacen a la estrecha victoria del Frente Amplio. En esta discusión se interpretan estos resultados a la luz del contexto político uruguayo y se evalúan sus implicaciones para la comprensión del comportamiento electoral, las contribuciones y limitaciones metodológicas, así como direcciones para futuras investigaciones. Asimismo, se integran los resultados de los análisis de robustez---sensibilidad, votos en blanco, verificaciones predictivas posteriores y covariables censales---que fortalecen la confianza en las estimaciones centrales.

% ========================================
\subsection{Interpretación de los Hallazgos}
\label{sec:discusion_interpretacion}

Los análisis revelan tres patrones interrelacionados que fundamentalmente reconfiguran la comprensión del balotaje uruguayo de 2024: las áreas rurales defeccionaron a tasas más altas que las áreas urbanas, los departamentos del interior perdieron masivamente votos de la coalición mientras que las áreas metropolitanas permanecieron relativamente estables, y la variación a nivel departamental excedió los 89 puntos porcentuales---mucho mayor que las brechas urbano-rurales o metropolitanas-interior. Estos hallazgos demuestran colectivamente que el colapso electoral de la coalición no fue un fenómeno nacional uniforme sino más bien una crisis geográficamente concentrada impulsada por votantes rurales y del interior.

\subsubsection{La Rebelión Rural: Contradiciendo la Sabiduría Convencional}

El hallazgo de que los circuitos rurales exhibieron una defección del 60.8\% de CA$\to$FA en comparación con el 39.0\% en circuitos urbanos contradice directamente las expectativas convencionales sobre el conservadurismo político rural en Uruguay. Los relatos históricos enfatizan el apego de los votantes rurales a los partidos tradicionales, particularmente al Partido Nacional, y caracterizan las áreas rurales como bastiones de la política de centro-derecha \citep{gonzalez1991electoral}. Los resultados revelan lo opuesto: los votantes rurales de la coalición fueron \textit{más} propensos a desertar al ideológicamente distante Frente Amplio, no menos. Esta brecha de 21.8 puntos porcentuales entre contextos rurales y urbanos resulta estadísticamente significativa, como lo indican los intervalos de credibilidad del 95\% no superpuestos (rural: [51.7\%, 69.8\%]; urbano: [35.1\%, 42.8\%]).

Esta rebelión rural se manifiesta más dramáticamente en las defecciones del Partido Colorado, donde los votantes rurales defeccionaron a una tasa 2.5 veces mayor que la tasa urbana (18.2\% vs 7.2\%). El Partido Colorado, que históricamente representa a las élites urbanas comerciales y profesionales, tiene una presencia organizacional limitada en áreas rurales \citep{buquet2015uruguayan}. Los votantes rurales del Colorado, careciendo de fuertes vínculos partidarios y potencialmente votando estratégicamente en la primera vuelta, parecen haber abandonado la coalición en masa cuando se enfrentaron a la elección binaria del balotaje. El patrón de defección sugiere que estos votantes priorizaron preocupaciones de política pública---probablemente relacionadas con la economía agrícola, la infraestructura rural o el desarrollo regional---sobre la alineación ideológica con la coalición de centro-derecha.

Notablemente, las defecciones del Partido Nacional mostraron una diferenciación urbano-rural mínima (8.4\% urbano vs 7.2\% rural), sugiriendo que los votantes del Nacional mantuvieron una lealtad relativamente consistente independientemente del contexto geográfico. Este contraste con los votantes del Colorado refuerza la interpretación de que la rebelión rural reflejó \textit{insatisfacción con las políticas} más que un rechazo generalizado de la política conservadora. Los votantes del Nacional, con identidades partidarias más fuertes arraigadas en la asociación histórica del partido con el Uruguay rural \citep{gonzalez1993structuring}, defeccionaron a tasas más bajas incluso cuando estaban insatisfechos con el desempeño de la coalición. Los votantes del Colorado, careciendo de tales vínculos profundos, resultaron ser mucho más volátiles.

¿Qué explica la insatisfacción de los votantes rurales con el gobierno de coalición? Si bien el diseño de inferencia ecológica no puede probar directamente mecanismos causales, varias explicaciones plausibles emergen del contexto político y económico uruguayo:

\begin{itemize}
    \item \textbf{Crisis económica agrícola}: El período 2020--2024 fue testigo de desafíos significativos para el sector agrícola de Uruguay, incluyendo sequía, caída de precios de commodities y aumento de costos de producción. Los votantes rurales directamente afectados por estas condiciones pueden haber culpado al gobierno de coalición por políticas de apoyo inadecuadas.

    \item \textbf{Infraestructura y servicios públicos}: Las áreas rurales históricamente reciben menos servicios gubernamentales e inversiones en infraestructura que las áreas urbanas. La gobernanza de la coalición puede haber sido percibida como priorizando intereses urbanos, dejando a los votantes rurales sintiéndose desatendidos.

    \item \textbf{Fracasos en el mensaje de la coalición}: Las comunicaciones de campaña que enfatizaban el crecimiento económico urbano, el turismo y los sectores tecnológicos pueden haber resonado pobremente en áreas rurales, donde los medios de vida dependen de la agricultura y la extracción de recursos.

    \item \textbf{Organización política local}: La fortaleza histórica del Frente Amplio en la organización rural, particularmente a través de sindicatos y movimientos agrarios, puede haber proporcionado redes de movilización efectivas para capturar votantes de la coalición insatisfechos.
\end{itemize}

La rebelión rural también tiene implicaciones para las teorías del comportamiento político rural. La investigación transnacional documenta cada vez más la polarización rural-urbana, con áreas rurales apoyando partidos populistas y de derecha en reacción a la globalización económica y el cambio cultural \citep{rodden2019s}. La elección uruguaya de 2024 presenta un patrón contrastante: los votantes rurales se desplazaron hacia la \textit{izquierda}, abandonando una coalición de centro-derecha por una alternativa izquierdista. Esta divergencia sugiere que el comportamiento político rural puede estar más moldeado por las condiciones económicas locales y el desempeño gubernamental que por dinámicas culturales universales.

\subsubsection{Colapso Interior y Estabilidad Metropolitana}

El contraste metropolitano-interior revela una división geográfica aún más dramática. Los departamentos del interior perdieron 34,628 votos netos de la coalición en comparación con solo 1,282 en el Área Metropolitana---una diferencia de 27 veces que identifica al interior como el epicentro geográfico del colapso electoral de la coalición. Esta disparidad refleja no solo tasas de defección más altas (55.5\% interior vs 47.4\% metropolitano para CA$\to$FA) sino también patrones cualitativamente diferentes de movimiento de votantes.

La diferencia interior-metropolitana más llamativa aparece en las defecciones del Partido Nacional. Los votantes del PN del interior defeccionaron al Frente Amplio a una tasa 48 veces mayor que los votantes del PN metropolitanos (8.1\% vs 0.17\%). Esto representa una ruptura histórica en los patrones de votación del Nacional. El Partido Nacional se identifica explícitamente como el partido del interior, rastreando sus orígenes a los caudillos rurales y los intereses agrícolas \citep{buquet2015uruguayan}. Que los votantes del Nacional del interior---el núcleo ideológico y geográfico del partido---abandonen la coalición a tasas tan elevadas indica una profunda insatisfacción con la gobernanza de la coalición en las provincias.

La estabilidad metropolitana, por el contrario, sugiere que los votantes de la coalición en Montevideo y Canelones enfrentaron un cálculo político diferente. El área metropolitana, que contiene la capital de Uruguay, centros comerciales y zonas industriales, experimentó una recuperación económica más robusta durante el período 2020--2024 de la coalición. El empleo urbano, los mercados inmobiliarios y los sectores de servicios mostraron crecimiento, potencialmente aislando a los votantes metropolitanos de la coalición de la insatisfacción que impulsó las defecciones del interior. Además, los votantes metropolitanos tienen mayor exposición mediática y contacto más directo con las élites políticas nacionales, potencialmente reforzando las identidades partidarias y la lealtad a la coalición.

La pérdida neta de votos del interior de 34,628 excede el margen final de victoria del Frente Amplio (aproximadamente 92,437 votos en el balotaje completo), demostrando que las dinámicas del interior fueron un componente determinante del resultado. Si los departamentos del interior hubieran exhibido tasas de defección a nivel metropolitano, la coalición habría retenido la presidencia. Esta concentración geográfica implica que la derrota de la coalición no resultó de una impopularidad generalizada sino de fallas específicas en la gobernanza del interior y rural---fallas que la coalición podría haber abordado mediante intervenciones de política dirigidas o estrategias de campaña.

\subsubsection{Heterogeneidad a Nivel Departamental y Contexto Local}

El rango a nivel departamental de 89.6 puntos porcentuales (6.0\% en Tacuarembó a 95.6\% en Río Negro) excede dramáticamente la brecha urbano-rural (21.8 pp) y la brecha metropolitana-interior (8.1 pp), revelando que el contexto subnacional importa mucho más de lo que sugieren las categorías geográficas agregadas. Esta variación extrema implica que los departamentos no son meramente réplicas más pequeñas de patrones nacionales sino que más bien exhiben dinámicas políticas distintivas moldeadas por factores locales.

La tasa de defección del 95.6\% de CA$\to$FA de Río Negro constituye el hallazgo más dramático del estudio. Casi todos los votantes de la coalición en Río Negro defeccionaron al Frente Amplio, sugiriendo que factores específicos del departamento sobrepasaron las tendencias políticas nacionales. La economía de Río Negro depende fuertemente de la agricultura y la generación de energía hidroeléctrica, con las represas de Palmar y Salto Grande proporcionando electricidad pero también creando preocupaciones ambientales. La angustia económica local, los conflictos ambientales o los efectos de la calidad de los candidatos pueden explicar el abandono casi total de la coalición. El resultado demanda una investigación cualitativa de estudio de caso para identificar los mecanismos específicos que impulsan un comportamiento tan extremo.

En el extremo opuesto, la defección del 6.0\% de CA$\to$FA de Tacuarembó combinada con una transferencia del 82.7\% de CA$\to$PN indica que los votantes de la coalición permanecieron leales a la coalición de centro-derecha, meramente cambiando su preferencia entre partidos de la coalición. Tacuarembó, ubicado en el interior norte, tiene fuertes tradiciones agrícolas y una historia de dominio del Partido Nacional. La fuerte organización local del PN, candidatos populares del PN o condiciones económicas distintivas pueden haber sostenido la lealtad a la coalición. El contraste entre Río Negro y Tacuarembó---departamentos interiores adyacentes con economías agrícolas similares pero patrones de defección diametralmente opuestos---ilustra que el simple determinismo económico no puede explicar el comportamiento electoral.

Varios conglomerados espaciales emergen del análisis a nivel departamental. Los departamentos costeros (Maldonado 67.4\%, Rocha 38.7\%) muestran defección elevada, potencialmente reflejando preocupaciones económicas en regiones dependientes del turismo afectadas por las dinámicas de recuperación pandémica. Los departamentos fronterizos del norte (Artigas, Rivera, Cerro Largo) muestran patrones mixtos, sugiriendo que la proximidad a Brasil y las economías fronterizas distintivas crean dinámicas políticas diferentes a las regiones agrícolas del interior. Los departamentos del suroeste (Colonia 54.2\%, Soriano 25.6\%) exhiben patrones divergentes a pesar de economías agrícolas similares, reforzando la importancia de factores políticos locales sobre la simple geografía económica.

Montevideo merece atención especial no por su tasa de defección---que con 26.4\% se ubica en el tercio inferior de los departamentos---sino por su peso demográfico. Con el 35\% de los votantes de CA a nivel nacional, Montevideo contribuyó un estimado de 5,561 votos mediante defecciones de CA$\to$FA, convirtiendo a la capital en un componente relevante del resultado. Esto demuestra que tanto las \textit{tasas} como los \textit{números absolutos} importan para los resultados electorales: tasas de defección altas en departamentos pequeños crean cambios locales dramáticos pero consecuencias nacionales limitadas, mientras que tasas de defección moderadas o bajas en departamentos populosos contribuyen sustancialmente al resultado nacional.

% ========================================
\subsection{Implicaciones Políticas}
\label{sec:discusion_implicaciones}

Los hallazgos conllevan implicaciones significativas para la comprensión de la estrategia partidaria, las prioridades de política gubernamental y la trayectoria futura de la política uruguaya. La concentración geográfica de las pérdidas de la coalición revela fallas específicas en la gobernanza rural y del interior que los actores políticos pueden potencialmente abordar mediante ajustes de política u reformas organizacionales.

\subsubsection{Fracasos de Gobernanza de la Coalición en el Interior}

La diferencia de 27 veces en las pérdidas netas de votos de la coalición entre las áreas del interior y metropolitanas sugiere que la gobernanza de la coalición desatendió sistemáticamente los intereses del interior durante el período 2020--2024. Varios dominios de política pueden haber contribuido a esta desatención:

\textbf{Política agrícola}: El gobierno de coalición, liderado por partidos de centro-derecha con vínculos históricos con los intereses agrícolas, paradójicamente perdió apoyo más dramáticamente en regiones rurales y agrícolas. Esto sugiere que las políticas agrícolas de la coalición no lograron abordar las preocupaciones de los productores sobre precios de commodities, costos de insumos o acceso a mercados. Los mensajes de campaña del Frente Amplio sobre apoyo agrícola y desarrollo rural parecen haber resonado con votantes que se sintieron abandonados por las prioridades económicas de la coalición.

\textbf{Desarrollo económico regional}: Los departamentos del interior sufren subdesarrollo crónico relativo al área metropolitana, con ingresos más bajos, menos oportunidades de empleo y acceso limitado a servicios de educación y salud \citep{ine2023census}. La gobernanza de la coalición puede haber concentrado las inversiones en infraestructura y los incentivos económicos en el área metropolitana, reforzando las percepciones de los votantes del interior de un sesgo centrado en la capital.

\textbf{Entrega de servicios públicos}: Las áreas rurales y del interior enfrentan mayores desafíos en la entrega de servicios públicos debido a la dispersión poblacional y las limitaciones de infraestructura. Las restricciones presupuestarias durante la pandemia y la recuperación post-pandemia pueden haber forzado reducciones de servicios que afectaron desproporcionadamente a las comunidades del interior, generando sentimiento anti-incumbente.

La pérdida de los partidarios centrales del Partido Nacional en el interior por parte de la coalición representa un fracaso particularmente profundo. Los votantes del Nacional del interior históricamente han proporcionado la base electoral más confiable del partido. Su defección a una tasa 48 veces mayor que la metropolitana indica que la gobernanza de la coalición violó expectativas fundamentales sobre la representación de los intereses del interior. Las futuras coaliciones de centro-derecha deben priorizar las preocupaciones de política rural y del interior para mantener viabilidad.

\subsubsection{Estrategia Rural e Interior del Frente Amplio}

El éxito del Frente Amplio en áreas rurales y del interior refleja una estrategia de campaña efectiva y capacidad organizacional. Varios factores pueden explicar la capacidad del FA para capturar votantes de la coalición desafectos:

\textbf{Mensajes dirigidos}: La campaña del FA parece haber elaborado mensajes que específicamente abordan preocupaciones rurales y del interior, enfatizando el apoyo agrícola, el desarrollo regional y la inversión en infraestructura. Esto contrasta con los mensajes de la coalición que enfatizaban el crecimiento económico urbano y la inversión internacional.

\textbf{Organización local}: El Frente Amplio mantiene una extensa organización de base en todo Uruguay, incluyendo áreas rurales y del interior, a través de sindicatos, movimientos sociales y comités partidarios. Esta presencia organizacional proporcionó redes para movilizar votantes de la coalición insatisfechos y contrarrestar el dominio tradicional de los partidos de la coalición en áreas rurales.

\textbf{Selección de candidatos}: La variación a nivel departamental sugiere que la calidad de los candidatos locales influyó en las tasas de defección. El FA puede haber presentado candidatos departamentales más fuertes en áreas de alta defección, mientras que los partidos de la coalición pueden haberse basado en ventajas de incumbencia sin reclutamiento competitivo de candidatos.

\textbf{Credibilidad económica}: El Frente Amplio gobernó Uruguay de 2005--2020, presidiendo sobre crecimiento económico y reducción de la pobreza antes de la pandemia. Esta experiencia de gobierno puede haber proporcionado credibilidad económica que permitió al FA competir por votantes de centro-derecha insatisfechos con la gestión económica de la coalición.

\subsubsection{Implicaciones para la Estabilidad del Sistema de Partidos}

La extrema volatilidad documentada plantea interrogantes sobre la estabilidad del sistema de partidos tradicionalmente institucionalizado de Uruguay. Uruguay ha sido caracterizado durante mucho tiempo como poseedor de uno de los sistemas de partidos más estables de América Latina, con fuertes identidades partidarias y patrones de votación consistentes \citep{mainwaring1995party}. Los análisis revelan una volatilidad electoral sustancial, con casi la mitad de los votantes de la coalición defeccionando al ideológicamente distante Frente Amplio.

Esta volatilidad puede reflejar:

\textbf{Debilitamiento de vínculos partidarios}: Las identidades partidarias tradicionales, particularmente para el Partido Colorado, pueden estar erosionándose a medida que el cambio demográfico y el reemplazo generacional reducen los vínculos partidarios históricos. Los votantes más jóvenes carecen de la intensa socialización partidaria que caracterizó a generaciones anteriores, produciendo una elección de voto más pragmática basada en desempeño y política.

\textbf{Votación basada en desempeño}: Las defecciones de la coalición concentradas en regiones económicamente angustiadas sugieren que los votantes cada vez más evalúan a los partidos basándose en el desempeño gubernamental más que en la alineación ideológica. Esta votación orientada al desempeño produce mayor volatilidad a medida que las condiciones económicas fluctúan.

\textbf{Dinámicas de coalición}: La estructura multipartidaria de la Coalición Multicolor puede haber creado confusión sobre la responsabilidad y la atribución de créditos. Los votantes de Cabildo Abierto, apoyando un partido relativamente nuevo sin raíces históricas profundas, pueden haber mantenido vínculos partidarios débiles que facilitaron la defección cuando el desempeño de la coalición los decepcionó.

Las implicaciones para el sistema de partidos se extienden más allá de la elección de 2024. Si los votantes rurales y del interior mantienen su apoyo al Frente Amplio en elecciones futuras, los partidos de la coalición enfrentan un potencial declive a largo plazo. Alternativamente, si la gobernanza del FA no logra abordar las preocupaciones rurales y del interior, los votantes pueden volver a inclinarse hacia partidos de centro-derecha, estableciendo un patrón de realineamiento volátil basado en el desempeño. La investigación futura debe rastrear si los patrones de defección de 2024 persisten o representan voto de protesta temporal.

% ========================================
\subsection{Contribuciones e Implicaciones Metodológicas}
\label{sec:discusion_metodologia}

Este estudio realiza varias contribuciones metodológicas al análisis electoral en Uruguay y a la investigación de inferencia ecológica en términos más amplios. Se demuestra el valor de la desagregación a nivel circuital, se establece la importancia de la estratificación geográfica y se ilustra la utilidad de la inferencia bayesiana para la cuantificación de incertidumbre en la estimación de transferencias de votos.

\subsubsection{Inferencia Ecológica a Nivel Circuital}

Los análisis previos de transferencias de votos uruguayas típicamente se basan en agregación a nivel departamental o comparación simple de votos sin modelado de inferencia ecológica \citep{bottinelli2021elecciones}. El enfoque a nivel circuital, analizando los 7,271 circuitos electorales, representa el estudio de inferencia ecológica geográficamente más desagregado de Uruguay hasta la fecha. Esta desagregación proporciona tres ventajas:

\textbf{Precisión}: La variación a nivel circuital en los resultados de primera vuelta y balotaje proporciona mucha más información estadística que la agregación a nivel departamental. Con 7,271 observaciones versus 19 departamentos, el análisis a nivel circuital reduce la incertidumbre en las estimaciones de probabilidades de transición y permite la detección de patrones sutiles oscurecidos por la agregación.

\textbf{Heterogeneidad geográfica}: Los datos a nivel circuital permiten análisis estratificados por clasificación urbana/rural, región metropolitana/interior y departamentos individuales. Estas estratificaciones revelan variación geográfica que la agregación a nivel departamental perdería, particularmente el contraste urbano-rural y el rango extremo de tasas de defección a nivel departamental.

\textbf{Sesgo de agregación reducido}: La inferencia ecológica sufre de sesgo de agregación cuando las unidades agregadas contienen poblaciones heterogéneas \citep{king1997}. Las unidades geográficas más pequeñas (circuitos) contienen poblaciones más homogéneas que las unidades más grandes (departamentos), reduciendo la heterogeneidad dentro de la unidad y mitigando el sesgo de agregación. El enfoque a nivel circuital proporciona por tanto estimaciones de probabilidades de transición más precisas que las alternativas a nivel departamental.

El enfoque a nivel circuital sí introduce desafíos computacionales. Estimar modelos de inferencia ecológica separados para cada departamento o estratificación requiere computación MCMC sustancial. Se abordan estos desafíos utilizando las implementaciones eficientes de muestreadores de PyMC \citep{abril2023pymc} y computación paralela. Los estudios futuros pueden construir sobre esta infraestructura computacional para realizar análisis aún más granulares.

\subsubsection{Robustez de las Estimaciones: Análisis de Sensibilidad}

Un aspecto metodológico central de este estudio es la evaluación sistemática de la robustez de las estimaciones, presentada en detalle en la Sección~\ref{sec:sensibilidad}. El análisis de sensibilidad sometió las estimaciones centrales a tres pruebas independientes: variación de distribuciones a priori (conservadoras, informativas), eliminación de circuitos atípicos (percentil superior al 1\%), y remuestreo bootstrap (1,000 réplicas).

Los resultados confirman la estabilidad de las estimaciones principales. La probabilidad de transición CA$\to$FA se mantuvo en un rango estrecho: 46.5\% con distribuciones a priori alternativas, 47.8\% tras la eliminación de atípicos, y 47.2\% ($\sigma = 3.0$ pp) en las réplicas bootstrap. La variación máxima observada de 1.3 puntos porcentuales respecto a la estimación base indica que las conclusiones sustantivas---particularmente la división casi igualitaria de los votantes de CA entre FA y PN---no dependen de decisiones de modelado particulares ni de la presencia de circuitos con comportamiento extremo. Esta robustez refuerza la interpretación de que la defección de CA constituyó un fenómeno estructural, no un artefacto estadístico.

\subsubsection{Integración de Covariables Censales}

En el Apéndice~\ref{sec:census} se presenta un análisis exploratorio que integra variables del censo de 2023 (edad mediana, proporción de educación terciaria, población) con las estimaciones de inferencia ecológica a nivel departamental. Aunque ninguna correlación alcanza significación estadística convencional ($p < 0.05$) con 19 observaciones departamentales, los resultados sugieren direcciones prometedoras. La proporción de educación terciaria muestra una correlación positiva moderada con la defección CA$\to$FA ($r = 0.38$, $p = 0.11$) y negativa con la retención CA$\to$PN ($r = -0.42$, $p = 0.07$), sugiriendo que departamentos con mayor nivel educativo experimentaron mayor transferencia de votantes de CA hacia el FA.

Estas correlaciones deben interpretarse con cautela dado que los datos censales utilizados son sintéticos (generados a partir de distribuciones departamentales plausibles) y el reducido número de unidades departamentales limita la potencia estadística. No obstante, los patrones observados son consistentes con la hipótesis de que el nivel educativo modera los patrones de defección de la coalición, posiblemente a través de mecanismos de exposición informativa o composición socioeconómica diferencial. La validación con microdatos oficiales del censo 2023 constituye una prioridad para investigación futura.

Los modelos de inferencia ecológica jerárquica con covariables a nivel circuital pueden:

\begin{itemize}
    \item Identificar qué características demográficas y económicas predicen las tasas de defección
    \item Aislar los efectos de urbanización, educación y edad de los efectos puramente geográficos
    \item Probar hipótesis específicas sobre angustia económica, polarización educativa o diferencias generacionales
    \item Proporcionar estimaciones más precisas mediante el intercambio de información entre circuitos similares a través del agrupamiento jerárquico
\end{itemize}

La disponibilidad de datos circuitales emparejados con el censo posiciona la investigación futura para avanzar más allá de patrones geográficos descriptivos hacia la identificación causal de los determinantes de defección. Este avance metodológico fortalecerá la comprensión del comportamiento electoral en Uruguay y proporcionará plantillas para el análisis electoral en otros contextos.

\subsubsection{Cuantificación de Incertidumbre Bayesiana}

El uso de inferencia ecológica bayesiana vía PyMC proporciona cuantificación de incertidumbre fundamentada a través de intervalos creíbles posteriores. Las implementaciones tradicionales de inferencia ecológica que utilizan estimación de máxima verosimilitud o método de momentos a menudo proporcionan estimaciones puntuales sin medidas adecuadas de incertidumbre \citep{rosen2001bayesian}. La inferencia bayesiana modela explícitamente la incertidumbre a través de distribuciones posteriores completas, permitiendo declaraciones probabilísticas sobre las probabilidades de transición.

Los intervalos creíbles reportados a lo largo de la sección de resultados transmiten incertidumbre genuina sobre el comportamiento del votante dados los datos electorales agregados. Donde las diferencias urbano-rurales muestran intervalos creíbles no superpuestos (urbano 39.0\%, IC [35.1, 42.8]; rural 60.8\%, IC [51.7, 69.8]), se puede afirmar confiadamente que la defección rural excedió la defección urbana. Donde las estimaciones a nivel departamental muestran intervalos creíbles amplios (Tacuarembó 6.0\%, IC [0.2, 19.1]), se reconoce incertidumbre sustancial reflejando datos limitados.

Esta cuantificación honesta de incertidumbre contrasta con aplicaciones de inferencia ecológica que reportan estimaciones puntuales implausiblemente precisas. Al adoptar la inferencia bayesiana, se proporcionan resultados que reflejan apropiadamente la incertidumbre inherente en inferir comportamiento individual a partir de datos agregados.

% ========================================
\subsection{Limitaciones}
\label{sec:discusion_limitaciones}

A pesar de las fortalezas metodológicas, el análisis enfrenta varias limitaciones importantes que califican la interpretación de resultados y sugieren cautela en la inferencia causal.

\subsubsection{Supuestos de la Inferencia Ecológica}

La inferencia ecológica se basa fundamentalmente en supuestos no comprobables sobre la relación entre patrones agregados y comportamiento individual \citep{king1997}. Aunque el método de King proporciona un marco estadístico fundamentado, no puede superar completamente el problema de inferencia ecológica: se observan únicamente totales de votos agregados, no elecciones de voto individuales. Varios supuestos merecen reconocimiento explícito:

\textbf{Ausencia de datos a nivel del votante}: Se carece de registros individuales de votos que establecerían definitivamente las probabilidades de transición. Todas las estimaciones representan inferencias estadísticas a partir de patrones agregados, sujetas a supuestos del modelo sobre cómo se comportan los individuos dentro de los circuitos electorales.

\textbf{Probabilidades de transición constantes dentro de estratos}: Los análisis estratificados asumen que todos los circuitos dentro de un estrato (por ejemplo, ``urbano'' o ``interior'') comparten las mismas probabilidades de transición. Este supuesto ignora la heterogeneidad dentro del estrato: no todos los circuitos urbanos se comportan idénticamente, ni todos los circuitos del interior. Los modelos jerárquicos podrían relajar este supuesto permitiendo desviaciones específicas del circuito alrededor de las medias del estrato.

\textbf{Ausencia de selección estratégica}: La inferencia ecológica asume que los límites de los circuitos electorales no seleccionan estratégicamente a los votantes para manipular patrones agregados. Aunque los límites circuitales de Uruguay siguen la geografía administrativa sin gerrymandering partidario, el ordenamiento residencial aún podría crear efectos de selección si los votantes de la coalición se concentran sistemáticamente en tipos particulares de circuitos.

\textbf{Precisión del conteo de votos}: Se asume que los totales de votos reportados reflejan con precisión las elecciones reales de los votantes. Aunque el sistema electoral de Uruguay mantiene altos estándares de integridad, errores de entrada de datos podrían introducir ruido.

\subsubsection{Subdispersión del Modelo y Verificaciones Predictivas Posteriores}

Las verificaciones predictivas posteriores (PPC), presentadas en detalle en el Apéndice~\ref{sec:ppc}, revelan una limitación importante del modelo de inferencia ecológica empleado. Si bien el modelo captura adecuadamente las medias de las distribuciones de votos---particularmente para el PN, cuya media predicha (46.98\%) resulta estadísticamente indistinguible de la observada (46.96\%, $p = 0.648$)---subestima sistemáticamente la varianza observada en los datos. Para los tres destinos de voto (FA, PN, Blancos), la desviación estándar predicha es inferior a la observada ($p = 0.000$ en todos los casos).

Esta subdispersión constituye una limitación conocida de los modelos de inferencia ecológica que asumen mayor homogeneidad dentro de los circuitos de la que realmente existe. El modelo predice una dispersión de 12.4 puntos porcentuales para el FA, mientras que la dispersión observada es de 13.5 pp. La discrepancia es aún más pronunciada para los votos en blanco, donde la dispersión predicha (0.02 pp) es varios órdenes de magnitud inferior a la observada (1.45 pp). Esto sugiere que los circuitos individuales presentan dinámicas de voto en blanco altamente heterogéneas que el modelo de transición uniforme no puede capturar.

Es importante señalar que esta subdispersión no invalida las estimaciones centrales de las probabilidades de transición. Las medias posteriores permanecen bien calibradas---como lo confirman tanto los PPC como el análisis de sensibilidad---pero los intervalos de credibilidad pueden ser algo estrechos, subestimando la verdadera incertidumbre. Los modelos jerárquicos con efectos aleatorios a nivel circuital podrían abordar esta limitación permitiendo que las probabilidades de transición varíen entre circuitos dentro de cada estrato.

\subsubsection{Agregación y Medición}

\textbf{Sesgo de agregación a nivel circuital}: Aunque el análisis a nivel circuital reduce el sesgo de agregación relativo al análisis a nivel departamental, los circuitos todavía agregan múltiples votantes. La heterogeneidad dentro del circuito permanece, y las estimaciones representan comportamiento promedio dentro de circuitos en lugar de parámetros individuales verdaderos.

\textbf{Clasificación urbano-rural}: La estratificación urbano-rural se basa en umbrales de densidad poblacional, que capturan imperfectamente el concepto multidimensional de ruralidad. Algunos circuitos clasificados como rurales pueden contener pueblos pequeños con características urbanas, mientras que algunos circuitos urbanos pueden incluir áreas agrícolas periurbanas. La clasificación errónea podría atenuar las verdaderas diferencias urbano-rurales.

\textbf{Circuitos faltantes}: Algunos circuitos con patrones de voto extremos o preocupaciones de calidad de datos pueden haber sido excluidos durante la limpieza de datos. Si los circuitos excluidos difieren sistemáticamente de los circuitos incluidos, las estimaciones pueden no ser generalizables a la población completa de circuitos.

\subsubsection{Variables Omitidas e Inferencia Causal}

El análisis documenta \textit{asociaciones} entre geografía y tasas de defección pero no puede establecer relaciones \textit{causales}. La estratificación geográfica revela dónde ocurrieron las defecciones pero no por qué. Varias variables omitidas pueden confundir la interpretación:

\textbf{Gasto e intensidad de campaña}: Se carece de datos sobre gasto de campaña, visitas de candidatos o exposición publicitaria por circuito. Las áreas de alta defección pueden haber recibido campaña intensiva del Frente Amplio o experimentado negligencia de campaña de la coalición, produciendo patrones geográficos que reflejan estrategia de campaña en lugar de preferencias subyacentes de los votantes.

\textbf{Mercados mediáticos}: La exposición mediática varía geográficamente, con áreas metropolitanas recibiendo cobertura más extensa que regiones del interior. Los patrones geográficos de defección pueden reflejar parcialmente influencia mediática diferencial en lugar de preocupaciones económicas o de política.

\textbf{Liderazgo local y maquinarias políticas}: Los departamentos con líderes partidarios locales fuertes o maquinarias políticas efectivas pueden exhibir diferentes patrones de defección independientemente de las condiciones económicas. No se puede distinguir efectos de liderazgo de efectos de política o económicos.

\textbf{Indicadores económicos}: Aunque se discutió la angustia económica como explicación potencial para las defecciones rurales y del interior, se omiten datos económicos a nivel circuital, que permanecen como variables no incluidas.

% ========================================
\subsection{Implicaciones del Análisis Temporal 2019-2024}
\label{sec:discusion_temporal}

El análisis comparativo entre 2019 y 2024, presentado en la Sección~\ref{sec:resultados_temporal}, revela una paradoja fundamental que reformula la comprensión del resultado electoral de 2024. A pesar de mejorar dramáticamente su cohesión interna---con todos los partidos coalicionales reduciendo defección entre 26.6 y 92.1 puntos porcentuales---la Coalición Republicana perdió igual. Esta subsección se centra en la \textit{interpretación} de esta paradoja y sus implicaciones teóricas y estratégicas, remitiendo al lector a la sección de resultados temporales para los datos detallados.

\subsubsection{La Primacía de la Primera Vuelta sobre la Cohesión de Segunda Vuelta}

El hallazgo más importante del análisis temporal es que \textit{la cohesión en segunda vuelta no compensa la derrota en primera vuelta}. La coalición perdió 154,251 votos (-12\%) entre primera vuelta 2019 y primera vuelta 2024, dejándola con una base sustancialmente menor para movilizar en el balotaje. Incluso con mejor retención---CA mejoró de 29.4\% a 44.2\% retención, PC de 64.5\% a 86.7\%, PN de 62.4\% a 90.3\%---la coalición no tenía votos suficientes para ganar.

Esta lección tiene implicaciones profundas para estrategia electoral en sistemas de dos vueltas. Los partidos pequeños volátiles representan activos riesgosos: su eventual lealtad mejorada en segunda vuelta no compensa el riesgo de colapso electoral que deja a la coalición sin votos suficientes. Cabildo Abierto ejemplifica esta dinámica---fue más leal en 2024 (+34.8 pp retención) pero había colapsado a 59,412 votos de 268,134, dejando a la coalición neta perdedora.

\textbf{Implicación estratégica}: Las coaliciones deben priorizar crecimiento o mantenimiento de base en primera vuelta sobre mejora de retención en segunda vuelta. Una coalición que pasa de 1,287,557 votos (2019) a 1,133,306 votos (2024) en primera vuelta pierde margen estratégico, sin importar cuánto mejore su cohesión posterior. La primera vuelta determina la magnitud del desafío---una coalición más pequeña pero más leal aún puede perder si el oponente tiene mayor base inicial.

\subsubsection{Reconfiguración del Electorado: Ausencia de Estabilidad Geográfica}

El segundo hallazgo crítico es la nula correlación geográfica entre patrones de 2019 y 2024 ($r < 0.13$ para todos los partidos). Un departamento con alta defección CA en 2019 no necesariamente tuvo alta defección en 2024. Esta inestabilidad sugiere cambios estructurales profundos en el electorado uruguayo---no simples fluctuaciones aleatorias, sino reconfiguración geográfica de patrones políticos.

Posibles mecanismos incluyen: (1) \textbf{cambio generacional}---entrada de nuevos votantes con preferencias geográficamente diferentes, (2) \textbf{efectos de la pandemia}---COVID-19 afectó diferencialmente a regiones, alterando prioridades políticas locales, (3) \textbf{realineamiento ideológico}---cambios en identificación partidaria a nivel departamental, (4) \textbf{efectos de gobierno}---cinco años de gobierno coalicional produjeron ganadores y perdedores geográficamente distribuidos.

\textbf{Implicación teórica}: Los modelos de comportamiento electoral que asumen estabilidad geográfica temporal son inadecuados para capturar dinámicas electorales en Uruguay 2019-2024. Las ``geografías electorales'' no son características fijas de lugares---son productos de composición del electorado, historia política, desempeño económico y estrategia de campaña, todos sujetos a cambio sustancial en ventanas de cinco años.

\subsubsection{El Destino de los 209 Mil Votantes Perdidos de Cabildo Abierto}

El colapso de CA entre 2019 (268,134 votos) y 2024 (59,412 votos) representa 208,722 votantes desaparecidos---78\% del electorado de CA. El análisis de inferencia ecológica solo captura flujos entre vueltas---no puede detectar votantes que cambiaron preferencias antes de la primera vuelta. El destino de estos 209 mil ex-votantes de CA permanece como interrogante crítico para comprender la elección.

Hipótesis sobre su destino incluyen: (1) \textbf{migración pre-electoral a FA}---votaron directamente FA en primera vuelta 2024, contribuyendo al crecimiento del FA de 1,066,470 (2019) a 1,192,518 (2024), (2) \textbf{migración a otros partidos coalicionales}---contribuyeron al crecimiento del PC (+87,327 votos) o atenuaron la caída del PN, (3) \textbf{abstención diferencial}---decidieron no votar en 2024, (4) \textbf{reemplazo generacional}---votantes de CA 2019 fueron reemplazados por nuevos votantes con diferentes preferencias.

\textbf{Implicación metodológica}: La inferencia ecológica entre vueltas tiene ``punto ciego'' fundamental---no detecta cambios pre-electorales en preferencias. Estudios futuros deben analizar flujos entre primera vuelta de elecciones sucesivas (2019$\to$2024) además de flujos dentro de elecciones (primera vuelta$\to$balotaje). Solo combinando ambos enfoques se puede mapear trayectorias completas del electorado.

\subsubsection{Consolidación Generalizada y Sus Límites}

Todos los partidos coalicionales mejoraron retención dramáticamente. Esta consolidación generalizada sugiere un proceso sistemático---no idiosincrasias partidarias, sino fuerzas comunes actuando en toda la coalición. Candidatos incluyen: (1) \textbf{aprendizaje organizacional}---cinco años de gobierno conjunto mejoraron coordinación y disciplina coalicional, (2) \textbf{polarización incrementada}---el sistema político uruguayo se polarizó, reduciendo transferencias entre bloques ideológicos, (3) \textbf{identidad coalicional fortalecida}---votantes de partidos coalicionales internalizaron identidad de ``Coalición Republicana'' como grupo político distintivo.

Sin embargo, esta consolidación tiene límites---mejoró retención de votantes existentes pero no impidió colapso de base electoral de CA ni atracción de nuevos votantes. La coalición optimizó lealtad de su electorado pero perdió la batalla por tamaño del electorado.

\textbf{Implicación política}: La consolidación coalicional es necesaria pero insuficiente para victoria electoral. Las coaliciones exitosas deben balancear dos objetivos: (1) retener votantes existentes (logrado por la coalición en 2024), (2) crecer o mantener base electoral inicial (fallado por la coalición en 2024). Enfocarse exclusivamente en cohesión interna ignora la amenaza más fundamental---contracción de base electoral.

\subsubsection{Instantánea Temporal y Seguimiento de Votos}

El análisis proporciona una instantánea de las transferencias de votos entre octubre y noviembre de 2024 pero no puede rastrear trayectorias individuales de votantes durante períodos de tiempo más largos. No se puede distinguir:

\begin{itemize}
    \item \textbf{Defección genuina}: Votantes de la coalición que cambiaron al Frente Amplio en el balotaje
    \item \textbf{Abstención diferencial}: Votantes de la coalición que se abstuvieron en el balotaje mientras los votantes del FA participaron
    \item \textbf{Movilización de nuevos votantes}: Votantes del balotaje por primera vez apoyando desproporcionadamente al FA
\end{itemize}

Aunque se interpretan las transiciones estimadas como defección, pueden reflejar parcialmente dinámicas de participación. Los datos de encuesta de panel siguiendo votantes individuales desde la primera vuelta hasta el balotaje distinguirían definitivamente la defección de la abstención, pero tales datos no existen para la elección de 2024 de Uruguay.

Adicionalmente, el análisis de inferencia ecológica no puede evaluar si las defecciones representan realineamiento permanente o voto de protesta temporal. Si los defectores de la coalición de 2024 regresan a partidos de centro-derecha en elecciones futuras, los hallazgos documentan protesta anti-incumbente en lugar de cambio electoral duradero. El análisis longitudinal rastreando los mismos circuitos a través de múltiples elecciones será necesario para evaluar la persistencia.

% ========================================
\subsection{Votos en Blanco y la Hipótesis de Protesta Pasiva}
\label{sec:discusion_blancos}

Una hipótesis alternativa al escenario de defección activa sostiene que los votantes de CA, insatisfechos con ambas opciones del balotaje, podrían haber optado por el voto en blanco como forma de protesta pasiva. El análisis de votos en blanco presentado en la Sección~\ref{sec:blank_votes} permite evaluar directamente esta hipótesis.

La correlación entre la proporción de votos en blanco y la proporción de apoyo a CA en primera vuelta resulta prácticamente nula ($r = -0.011$, $n = 7{,}271$ circuitos). Esto descarta la hipótesis de protesta pasiva: los circuitos con mayor presencia de CA no registraron más votos en blanco en el balotaje. Si los votantes de CA hubieran canalizado su insatisfacción hacia el voto en blanco, se esperaría una correlación positiva significativa. Su ausencia refuerza la interpretación de que las transferencias CA$\to$FA representan decisiones activas de los votantes, no mera abstención parcial.

Adicionalmente, los votos en blanco muestran correlaciones negativas con el apoyo al PC ($r = -0.108$) y al PN ($r = -0.196$), y positiva con el FA ($r = 0.167$). Estos patrones sugieren que el voto en blanco se concentra en circuitos con mayor presencia del FA---posiblemente reflejando dinámicas de participación en bastiones del Frente Amplio donde el resultado del balotaje era localmente predecible, reduciendo la motivación individual de votantes marginales. La asociación negativa con los partidos de la coalición indica que el voto en blanco no fue un mecanismo relevante de protesta coalicional.

El hallazgo tiene implicaciones interpretativas relevantes: la defección de CA hacia el FA fue un fenómeno de elección activa, no de apatía o confusión. Los votantes de CA que transfirieron su apoyo al FA lo hicieron como decisión deliberada, descartando explicaciones basadas en la desmotivación o el desconocimiento de la oferta electoral.

% ========================================
\subsection{Direcciones de Investigación Futura}
\label{sec:discusion_futuro}

Los hallazgos abren múltiples vías para investigación futura que pueden profundizar la comprensión del comportamiento electoral en Uruguay y refinar los métodos de inferencia ecológica para la estimación de transferencias de votos.

\subsubsection{Modelos Jerárquicos con Covariables}

La extensión más inmediata implica estimar modelos jerárquicos de inferencia ecológica que incorporen covariables a nivel circuital. El conjunto de datos preparado incluye variables del censo 2023 (niveles de educación, distribución etaria, estructura de hogares) y características electorales (tamaño del circuito, urbanización, resultados de elecciones previas) que podrían predecir tasas de defección. Los resultados preliminares del Apéndice~\ref{sec:census} sugieren que la educación terciaria puede ser un predictor relevante ($r = 0.38$ con defección CA$\to$FA), aunque se requiere validación con microdatos censales oficiales. Los modelos jerárquicos permitirían:

\begin{itemize}
    \item Identificar predictores socioeconómicos y demográficos de la defección de la coalición
    \item Estimar cómo la educación, edad y estatus económico moderan los efectos urbano-rural y metropolitano-interior
    \item Probar hipótesis específicas sobre polarización educativa o distrés económico
    \item Mejorar la precisión mediante pooling parcial de información entre circuitos similares
    \item Habilitar la predicción de tasas de defección en circuitos con datos faltantes
\end{itemize}

Asimismo, los modelos jerárquicos con efectos aleatorios a nivel circuital podrían abordar la subdispersión identificada en los PPC (Apéndice~\ref{sec:ppc}), permitiendo que las probabilidades de transición varíen entre circuitos y capturando la heterogeneidad que el modelo actual no puede representar.

La implementación extendería el marco de inferencia ecológica de King con estructuras de regresión que vinculen las probabilidades de transición con covariables, estimadas mediante modelos jerárquicos bayesianos en PyMC \citep{abril2023pymc}. Este enfoque movería el análisis desde patrones descriptivos hacia modelado explicativo.

\subsubsection{Validación con Encuestas y Encuestas de Boca de Urna}

La inferencia ecológica proporciona estimaciones estadísticas de transferencias de votos pero carece de validación con datos de verdad fundamental. La investigación mediante encuestas puede proporcionar evidencia complementaria:

\textbf{Comparación con encuestas de boca de urna}: Si las encuestas de boca de urna o encuestas poselectorales preguntaron a los encuestados sobre sus elecciones de voto en primera vuelta y balotaje, permitiría validar las probabilidades de transición basadas en encuestas contra las estimaciones de inferencia ecológica. La concordancia validaría el enfoque; la discordancia identificaría problemas de especificación o sesgos de encuesta.

\textbf{Entrevistas cualitativas}: Entrevistas en profundidad con votantes en departamentos de alta defección (Río Negro, Treinta y Tres, Maldonado) y departamentos de baja defección (Tacuarembó, Cerro Largo) podrían iluminar las motivaciones detrás de las decisiones de defección. La evidencia cualitativa probaría los mecanismos hipotéticos propuestos sobre distrés económico, insatisfacción política y efectividad de campaña.

\textbf{Datos de gastos de campaña}: Vincular datos de gastos de campaña por departamento con tasas de defección probaría si los patrones geográficos reflejan estrategia de campaña. Altas defecciones en departamentos con campaña intensiva del FA apoyarían explicaciones de movilización; altas defecciones en departamentos con baja campaña apoyarían explicaciones basadas en agravios.

\subsubsection{Tendencias Históricas y Análisis Longitudinal}

El análisis se enfoca en la elección 2024 con comparación temporal a 2019. Varias extensiones proporcionarían perspectiva temporal adicional:

\textbf{Análisis más profundo de 2019}: Aunque se han aplicado métodos idénticos al balotaje uruguayo de 2019, análisis adicionales podrían explorar dinámicas de primera vuelta 2019 comparada con primera vuelta 2024, capturando flujos pre-electorales que la inferencia ecológica intra-elección no detecta. Analizar flujos entre primera vuelta de elecciones sucesivas (2019$\to$2024) revelaría el destino de los 209 mil votantes perdidos de CA.

\textbf{Análisis de panel entre elecciones}: Construir un conjunto de datos de panel con resultados a nivel circuital de múltiples elecciones (2009, 2014, 2019, 2024) habilitaría modelos de efectos fijos controlando por características circuitales invariantes en el tiempo. Esto aislaría la variación temporal en tasas de defección de diferencias geográficas estables.

\textbf{Análisis de series temporales de trayectorias departamentales}: Rastrear cómo los patrones de votación de departamentos individuales evolucionaron a través de múltiples elecciones identificaría si los departamentos de alta defección tienen historias de volatilidad o si 2024 representa una ruptura de patrones estables.

\subsubsection{Análisis Comparativo}

La elección uruguaya de 2024 invita a la comparación con otros contextos latinoamericanos e internacionales:

\textbf{Comparaciones latinoamericanas}: Aplicar inferencia ecológica a elecciones de balotaje en otros países latinoamericanos con sistemas electorales similares (Brasil, Argentina, Perú) probaría si la rebelión rural y el colapso interior representan fenómenos específicos de Uruguay o patrones regionales más amplios.

\textbf{Literatura sobre polarización rural-urbana}: Los hallazgos contribuyen a debates sobre polarización política rural-urbana. Mientras que mucha investigación documenta áreas rurales apoyando partidos populistas de derecha \citep{rodden2019s}, Uruguay muestra áreas rurales desplazándose hacia la izquierda. El análisis comparativo podría identificar condiciones bajo las cuales los votantes rurales apoyan alternativas de izquierda versus derecha.

\textbf{Comparaciones metodológicas}: Comparar las estimaciones de inferencia ecológica de King con métodos alternativos de inferencia ecológica (ajuste proporcional iterativo, modelos jerárquicos bayesianos, enfoques de aprendizaje automático) evaluaría la robustez de los resultados a elecciones metodológicas.

\subsubsection{Extensiones Sustantivas}

Varias preguntas sustantivas permanecen sin abordar:

\textbf{Dinámicas de género y edad}: Los datos electorales de Uruguay no desagregan por género o edad, pero los datos del censo proporcionan composición demográfica a nivel circuital. Vincular la composición demográfica con tasas de defección podría probar si los votantes más jóvenes o las votantes mujeres impulsaron las pérdidas de la coalición.

\textbf{Prioridades políticas}: Investigación mediante encuestas preguntando a los votantes sobre prioridades políticas (economía, seguridad, salud, educación) podría probar si los votantes rurales priorizaron política agrícola, los votantes del interior priorizaron desarrollo regional, o si otras dimensiones temáticas estructuraron la defección.

\textbf{Consumo de medios}: Analizar mercados de medios y patrones de consumo por departamento probaría si los patrones geográficos de defección se correlacionan con exposición mediática, revelando potencialmente influencia de medios en la elección de voto.

% ========================================
\subsection{Comentarios Finales}
\label{sec:discusion_conclusion}

El balotaje presidencial uruguayo de 2024 desafió las expectativas convencionales sobre geografía electoral, con votantes rurales y del interior impulsando la derrota de la coalición mediante tasas elevadas de defección hacia el Frente Amplio. El análisis de inferencia ecológica, aprovechando datos a nivel circuital y métodos estadísticos bayesianos, revela que el resultado electoral se debió a comportamiento votante geográficamente heterogéneo en lugar de oscilaciones nacionales uniformes. Los departamentos del interior perdieron 27 veces más votos netos de la coalición que las áreas metropolitanas, los circuitos rurales defeccionaron a 1.6 veces la tasa urbana (60.8\% vs 39.0\%), y la variación a nivel departamental abarcó casi 90 puntos porcentuales.

Los análisis de robustez refuerzan la confianza en estos hallazgos. Las estimaciones centrales se mantienen estables ante variaciones en distribuciones a priori, eliminación de atípicos y remuestreo bootstrap (variación máxima de 1.3 pp). La ausencia de correlación entre votos en blanco y apoyo a CA ($r = -0.011$) descarta la hipótesis de protesta pasiva, confirmando que las transferencias CA$\to$FA fueron elecciones activas. Las verificaciones predictivas posteriores validan la calibración del modelo en medias pero identifican subdispersión como limitación---el modelo asume mayor homogeneidad circuital de la existente---señalando la dirección hacia modelos jerárquicos con efectos aleatorios.

Estos hallazgos conllevan implicaciones significativas para comprender la política uruguaya. El colapso electoral de la coalición reflejó fracasos específicos en la gobernanza rural e interior---fracasos que generaron profunda insatisfacción entre votantes que históricamente apoyaron partidos de centro-derecha. La victoria del Frente Amplio resultó de movilización rural e interior efectiva, capturando votantes desafectos de la coalición mediante mensajes dirigidos y organización de base. Si estas defecciones representan realineamiento duradero o protesta temporal moldeará la trayectoria política de Uruguay en elecciones venideras.

Metodológicamente, este estudio demuestra el valor de la inferencia ecológica a nivel circuital para comprender transferencias de votos. El análisis desagregado revela heterogeneidad geográfica que los estudios agregados omiten, mientras que la inferencia bayesiana proporciona cuantificación de incertidumbre con principios sólidos. Investigación futura que incorpore covariables socioeconómicas---particularmente las sugeridas por las correlaciones exploratorias con educación terciaria---validación mediante encuestas y comparaciones históricas refinará aún más la comprensión de dinámicas electorales en Uruguay y más allá.

La rebelión rural y el colapso interior desafían narrativas simplistas sobre comportamiento electoral. El contexto geográfico importa profundamente---no como ruido de fondo sino como determinante fundamental de cómo los votantes responden a alternativas políticas. Comprender elecciones requiere comprender lugares, y comprender lugares requiere datos granulares, métodos sofisticados e interpretación cuidadosa. Este estudio proporciona una fundación para tal comprensión del balotaje uruguayo de 2024 y ofrece una plantilla para investigación de inferencia ecológica en otros contextos electorales.

% ========================================
% END OF DISCUSSION SECTION
% ========================================

\newpage

% ========================================
% NOTA: Tablas y figuras están integradas en las secciones correspondientes
% ========================================
% Las tablas se incluyen directamente después de ser referenciadas en:
% - sections/results.tex (tablas 2024)
% - sections/results_temporal.tex (tablas análisis temporal)
% - sections/08_sensitivity_analysis.tex (tablas sensibilidad)
% - sections/09_blank_votes.tex (tablas votos en blanco)
% Las figuras están integradas en sus secciones respectivas

% ========================================
% BIBLIOGRAFÍA
% ========================================
\FloatBarrier                        % Procesar todos los floats pendientes
\newpage
\bibliography{references}

% ========================================
% APÉNDICES: VALIDACIÓN TÉCNICA
% ========================================
\newpage
\appendix

% Posterior Predictive Checks
\section{Posterior Predictive Checks}
\label{sec:ppc}

Los \textit{posterior predictive checks} (PPC) constituyen una herramienta fundamental para evaluar la adecuaci\'on del modelo bayesiano implementado. Esta t\'ecnica permite determinar si el modelo propuesto es capaz de generar datos sint\'eticos coherentes con las observaciones reales, proporcionando as\'i una medida de la calidad del ajuste y de la capacidad predictiva del modelo.

En el contexto de la inferencia ecol\'ogica, donde se busca estimar par\'ametros no observables a partir de datos agregados, los PPC permiten verificar que las estimaciones obtenidas no solo sean internamente consistentes desde el punto de vista estad\'istico, sino que tambi\'en reproduzcan adecuadamente las caracter\'isticas distribucionales de los datos observados.

\subsection{Metodolog\'ia}
\label{subsec:ppc_metodo}

La metodolog\'ia de posterior predictive checks se basa en la generaci\'on de datos sint\'eticos a partir de la distribuci\'on posterior del modelo. Para cada conjunto de par\'ametros $\theta^{(s)}$ muestreado de la distribuci\'on posterior $p(\theta | y_{\text{obs}})$, se genera un conjunto de datos replicados $y^{\text{rep}(s)}$ mediante:

\begin{equation}
y^{\text{rep}(s)} \sim p(y | \theta^{(s)})
\end{equation}

donde $y_{\text{obs}}$ representa los datos observados (proporciones de votos en el ballotage para cada lema en cada circuito) y $y^{\text{rep}}$ son las r\'eplicas generadas por el modelo.

\subsubsection{Proceso de Muestreo}

Para el an\'alisis nacional de 2024, se implement\'o el siguiente procedimiento:

\begin{enumerate}
    \item Se extrajeron 500 muestras aleatorias de la distribuci\'on posterior de los par\'ametros $\beta_{ij}$ (probabilidades de transferencia de votos).

    \item Para cada muestra $s = 1, \ldots, 500$, se gener\'o un conjunto completo de datos replicados para los 7,271 circuitos electorales, calculando las proporciones predichas de votos en ballotage:
    \begin{itemize}
        \item Proporci\'on predicha para Frente Amplio: $p_{\text{FA},c}^{\text{rep}} = \sum_{j} w_{jc} \beta_{j,\text{FA}}^{(s)}$
        \item Proporci\'on predicha para Partido Nacional: $p_{\text{PN},c}^{\text{rep}} = \sum_{j} w_{jc} \beta_{j,\text{PN}}^{(s)}$
        \item Proporci\'on predicha para votos en blanco: $p_{\text{Blancos},c}^{\text{rep}} = 1 - p_{\text{FA},c}^{\text{rep}} - p_{\text{PN},c}^{\text{rep}}$
    \end{itemize}
    donde $w_{jc}$ representa la proporci\'on de votos del lema $j$ en el circuito $c$ respecto al total de primera vuelta.

    \item Se calcularon estad\'isticos de prueba tanto para los datos observados como para cada r\'eplica generada.
\end{enumerate}

\subsubsection{Estad\'isticos de Prueba}

Se evaluaron cuatro estad\'isticos de prueba fundamentales para cada variable de resultado (FA, PN, Blancos), operando sobre las proporciones de votos por circuito:

\begin{itemize}
    \item \textbf{Media}: Mide la tendencia central de la distribuci\'on de proporciones entre circuitos.
    \begin{equation}
    T_{\text{mean}}(y) = \frac{1}{n}\sum_{c=1}^{n} y_c
    \end{equation}

    \item \textbf{Desviaci\'on est\'andar}: Captura la dispersi\'on de las proporciones entre circuitos.
    \begin{equation}
    T_{\text{std}}(y) = \sqrt{\frac{1}{n-1}\sum_{c=1}^{n} (y_c - \bar{y})^2}
    \end{equation}

    \item \textbf{M\'inimo}: Identifica el comportamiento en circuitos con menor proporci\'on de votaci\'on.
    \begin{equation}
    T_{\text{min}}(y) = \min_{c} y_c
    \end{equation}

    \item \textbf{M\'aximo}: Identifica el comportamiento en circuitos con mayor proporci\'on de votaci\'on.
    \begin{equation}
    T_{\text{max}}(y) = \max_{c} y_c
    \end{equation}
\end{itemize}

\subsubsection{P-valores Bayesianos}

Para cada estad\'istico de prueba, se calcul\'o un p-valor bayesiano que cuantifica la proporci\'on de r\'eplicas posteriores cuyo estad\'istico iguala o supera al observado:

\begin{equation}
p_B = P(T(y^{\text{rep}}) \geq T(y_{\text{obs}}) | y_{\text{obs}}) = \frac{1}{S}\sum_{s=1}^{S} \mathbb{1}(T(y^{\text{rep}(s)}) \geq T(y_{\text{obs}}))
\end{equation}

donde $S = 500$ es el n\'umero de r\'eplicas posteriores generadas.

La interpretaci\'on de los p-valores bayesianos es la siguiente:
\begin{itemize}
    \item $p_B \approx 0.5$: Excelente ajuste --- el valor observado se sit\'ua cerca del centro de la distribuci\'on posterior predictiva.
    \item $0.05 < p_B < 0.95$: Ajuste aceptable --- el valor observado es compatible con las r\'eplicas generadas.
    \item $p_B < 0.05$: El modelo genera sistem\'aticamente valores \emph{menores} que el observado (subestimaci\'on sistem\'atica del estad\'istico).
    \item $p_B > 0.95$: El modelo genera sistem\'aticamente valores \emph{mayores} que el observado (sobreestimaci\'on sistem\'atica del estad\'istico).
    \item $p_B = 0.000$ o $p_B = 1.000$: Discrepancia extrema --- el valor observado se encuentra fuera del rango de todas las r\'eplicas posteriores, lo que indica desajuste sistem\'atico del modelo en ese aspecto espec\'ifico.
\end{itemize}

\subsection{Resultados}
\label{subsec:ppc_resultados}

La Tabla~\ref{tab:ppc_summary} presenta un resumen detallado de los posterior predictive checks realizados para el modelo nacional de 2024. Se reportan los estad\'isticos de prueba calculados sobre las proporciones observadas y predichas, junto con sus correspondientes intervalos de credibilidad al 95\% y p-valores bayesianos.

\begin{table}
\caption{Posterior Predictive Checks: Estadísticos de Prueba}
\label{tab:ppc_summary}
\begin{tabular}{llcccr}
\toprule
Partido & Estadístico & Observado & Predicho (Mediana) & IC 95\% & p-valor \\
\midrule
FA & Media & 0.4880 & 0.4866 & {[}0.4859, 0.4872{]} & 0.000 \\
FA & Desv. Est. & 0.1348 & 0.1244 & {[}0.1241, 0.1247{]} & 0.000 \\
FA & Mínimo & 0.0000 & 0.1033 & --- & 0.000 \\
FA & Máximo & 0.8179 & 0.7785 & --- & 0.000 \\
PN & Media & 0.4696 & 0.4698 & {[}0.4691, 0.4705{]} & 0.648 \\
PN & Desv. Est. & 0.1373 & 0.1246 & {[}0.1245, 0.1248{]} & 0.000 \\
PN & Mínimo & 0.1643 & 0.1773 & --- & 0.000 \\
PN & Máximo & 1.0000 & 0.8538 & --- & 0.000 \\
Blancos & Media & 0.0424 & 0.0437 & {[}0.0434, 0.0439{]} & 1.000 \\
Blancos & Desv. Est. & 0.0145 & 0.0002 & {[}0.0000, 0.0005{]} & 0.000 \\
Blancos & Mínimo & 0.0000 & 0.0429 & --- & 0.000 \\
Blancos & Máximo & 0.1667 & 0.0443 & --- & 0.000 \\
\bottomrule
\end{tabular}
\end{table}


Los resultados revelan un patr\'on diferencial de ajuste seg\'un la variable de resultado y el estad\'istico evaluado. A continuaci\'on se analiza cada componente en detalle.

\subsubsection{Ajuste para Frente Amplio}

\begin{itemize}
    \item \textbf{Media} ($p = 0.000$): La proporci\'on media observada de votos a FA por circuito es 0.4880, mientras que la mediana de las r\'eplicas posteriores es 0.4866, con un intervalo de credibilidad al 95\% de $[0.4859,\; 0.4872]$. El valor observado se encuentra \emph{por encima} del l\'imite superior de este intervalo, lo que indica que el modelo subestima sistem\'aticamente la proporci\'on media de FA. El p-valor de 0.000 confirma que ninguna de las 500 r\'eplicas alcanza el valor observado. No obstante, la magnitud del sesgo es reducida (0.0014 puntos, equivalente a 0.14 puntos porcentuales), lo que se hace detectable \'unicamente por el gran n\'umero de circuitos ($n = 7{,}271$).

    \item \textbf{Desviaci\'on est\'andar} ($p = 0.000$): La desviaci\'on est\'andar observada entre circuitos es 0.1348, frente a una mediana predicha de 0.1244 (IC~95\%: $[0.1241,\; 0.1247]$). El modelo subestima la variabilidad en 0.0104 puntos, lo que equivale a un 7.7\% de la dispersi\'on real. Este resultado indica que el modelo genera r\'eplicas m\'as homog\'eneas que los datos observados.

    \item \textbf{M\'inimo} ($p = 0.000$): El m\'inimo observado es 0.0000 (circuitos con proporci\'on nula de FA), mientras que el modelo predice un m\'inimo mediano de 0.1033. El modelo no es capaz de reproducir circuitos con proporci\'on extremadamente baja de FA, lo cual refleja la existencia de circuitos at\'ipicos (posiblemente rurales o de muy baja habilitaci\'on) con din\'amicas locales que el modelo nacional no captura.

    \item \textbf{M\'aximo} ($p = 0.000$): El m\'aximo observado es 0.8179, superior a la predicci\'on mediana de 0.7785. El modelo subestima el valor m\'aximo, lo que sugiere la presencia de circuitos con concentraciones excepcionales de voto a FA que exceden lo previsto por las probabilidades de transferencia nacionales.
\end{itemize}

\subsubsection{Ajuste para Partido Nacional}

\begin{itemize}
    \item \textbf{Media} ($p = 0.648$): La proporci\'on media observada de votos a PN es 0.4696, y la mediana predicha es 0.4698 (IC~95\%: $[0.4691,\; 0.4705]$). El valor observado se sit\'ua c\'omodamente dentro del intervalo de credibilidad, y el p-valor de 0.648 ---pr\'oximo al \'optimo de 0.5--- indica un \textbf{excelente ajuste} del modelo a la tendencia central de la proporci\'on de PN. Este es el mejor resultado de ajuste obtenido en todo el an\'alisis PPC.

    \item \textbf{Desviaci\'on est\'andar} ($p = 0.000$): La desviaci\'on est\'andar observada es 0.1373, mientras que la mediana predicha es 0.1246 (IC~95\%: $[0.1245,\; 0.1248]$). Al igual que para FA, el modelo subestima la variabilidad entre circuitos en 0.0127 puntos (9.2\% de la dispersi\'on observada). Todas las r\'eplicas exhiben menor dispersi\'on que los datos reales.

    \item \textbf{M\'inimo} ($p = 0.000$): El m\'inimo observado es 0.1643, frente a una predicci\'on mediana de 0.1773. El modelo no reproduce los circuitos con menor proporci\'on de PN, aunque la discrepancia (0.013) es menor que la observada para FA.

    \item \textbf{M\'aximo} ($p = 0.000$): El m\'aximo observado es 1.0000 (circuitos con la totalidad de votos a PN), frente a una predicci\'on mediana de 0.8538. El modelo no puede generar proporciones cercanas a la unidad, lo que refleja la existencia de circuitos muy peque\~nos donde el 100\% de los votos se dirigen a un solo candidato.
\end{itemize}

\subsubsection{Ajuste para Votos en Blanco}

Los votos en blanco presentan el patr\'on de ajuste m\'as problem\'atico:

\begin{itemize}
    \item \textbf{Media} ($p = 1.000$): La proporci\'on media observada de votos en blanco es 0.0424, mientras que la mediana predicha es 0.0437 (IC~95\%: $[0.0434,\; 0.0439]$). El modelo sobreestima sistem\'aticamente la proporci\'on de votos en blanco. El p-valor de 1.000 indica que \emph{todas} las r\'eplicas generan una proporci\'on media superior a la observada. Este resultado es consistente con la estructura del modelo: dado que los votos en blanco se calculan por diferencia ($1 - p_{\text{FA}} - p_{\text{PN}}$), la leve subestimaci\'on de las proporciones de FA y PN se acumula en una sobreestimaci\'on de los votos en blanco.

    \item \textbf{Desviaci\'on est\'andar} ($p = 0.000$): La desviaci\'on est\'andar observada es 0.0145, mientras que la mediana predicha es pr\'acticamente nula: 0.0002 (IC~95\%: $[0.0000,\; 0.0005]$). Esta es la discrepancia m\'as severa de todo el an\'alisis: el modelo predice una dispersi\'on de votos en blanco casi inexistente, 72 veces menor que la observada. El modelo trata los votos en blanco como un residuo virtualmente constante, sin capturar la variabilidad real entre circuitos. Esto sugiere que la dispersi\'on de votos en blanco responde a factores locales (accesibilidad, composici\'on del electorado, particularidades de la votaci\'on) que el modelo agregado ignora por completo.

    \item \textbf{M\'inimo} ($p = 0.000$): El m\'inimo observado es 0.0000 (circuitos sin ning\'un voto en blanco), frente a una predicci\'on mediana de 0.0429. El modelo predice un piso de votos en blanco pr\'acticamente igual a la media, confirmando su incapacidad para generar variabilidad en esta categor\'ia.

    \item \textbf{M\'aximo} ($p = 0.000$): El m\'aximo observado es 0.1667 (16.7\% de votos en blanco en alg\'un circuito), mientras que la predicci\'on mediana es apenas 0.0443. El modelo no puede reproducir circuitos con proporciones altas de voto en blanco, resultado directo de la pr\'acticamente nula dispersi\'on predicha.
\end{itemize}

\subsubsection{Visualizaci\'on de Diagn\'osticos}

La Figura~\ref{fig:ppc_diag} presenta una comparaci\'on visual entre las distribuciones observadas y las distribuciones posteriores predictivas para cada variable de resultado. Los histogramas muestran la densidad de los datos observados (en azul) superpuesta con la densidad de las r\'eplicas posteriores (en rojo transl\'ucido).

\begin{figure}[H]
\centering
\includegraphics[width=\textwidth]{../../outputs/figures/ppc_diagnostics.pdf}
\caption{Diagn\'osticos de posterior predictive checks: distribuciones observadas (azul) versus distribuciones posteriores predictivas (rojo transl\'ucido) para proporciones de votos en ballotage. Se visualizan 500 r\'eplicas posteriores para evaluar la capacidad del modelo de reproducir las caracter\'isticas distribucionales de los datos observados.}
\label{fig:ppc_diag}
\end{figure}

Se observa que:
\begin{itemize}
    \item Para FA y PN, las distribuciones posteriores predictivas envuelven razonablemente la distribuci\'on observada en t\'erminos de localizaci\'on, pero presentan menor amplitud que los datos reales, consistente con la subdispersi\'on detectada en los estad\'isticos de desviaci\'on est\'andar.
    \item Para votos en blanco, se aprecia una concentraci\'on extrema de las r\'eplicas posteriores en torno a un valor \'unico, en marcado contraste con la distribuci\'on observada que exhibe dispersi\'on apreciable.
\end{itemize}

La Figura~\ref{fig:ppc_qq} complementa el an\'alisis mediante Q-Q plots (quantile-quantile plots), que comparan los cuantiles de la distribuci\'on observada con los cuantiles de la distribuci\'on posterior predictiva. Una l\'inea diagonal perfecta indicar\'ia ajuste perfecto.

\begin{figure}[H]
\centering
\includegraphics[width=\textwidth]{../../outputs/figures/ppc_qqplot.pdf}
\caption{Q-Q plots comparando cuantiles observados versus cuantiles posteriores predictivos para cada variable de resultado. La l\'inea roja representa la identidad (ajuste perfecto). Las desviaciones de esta l\'inea indican aspectos en los que el modelo no captura completamente la distribuci\'on observada. Los intervalos de credibilidad al 95\% (banda gris) permiten evaluar la significancia de las desviaciones.}
\label{fig:ppc_qq}
\end{figure}

Los Q-Q plots revelan:
\begin{itemize}
    \item \textbf{FA}: Ajuste razonable en la porci\'on central de la distribuci\'on, con desviaciones en ambas colas. La cola inferior (circuitos con baja proporci\'on de FA) muestra que el modelo sobreestima estos valores, mientras que la cola superior revela subestimaci\'on en circuitos con alta proporci\'on de FA. Este patr\'on es consistente con la subdispersi\'on general detectada.

    \item \textbf{PN}: Patr\'on similar al de FA, con buen ajuste central y desviaciones en las colas. Las desviaciones son particularmente notables en la cola superior, donde el m\'aximo observado de 1.0000 no es reproducible por el modelo.

    \item \textbf{Blancos}: Desviaci\'on severa en toda la distribuci\'on. Los cuantiles predichos se concentran en un rango extremadamente estrecho, mientras que los observados se distribuyen en un rango mucho m\'as amplio. Este resultado confirma visualmente la pr\'acticamente nula variabilidad que el modelo asigna a los votos en blanco.
\end{itemize}

\subsection{Interpretaci\'on}
\label{subsec:ppc_interpretacion}

El an\'alisis de posterior predictive checks permite establecer las siguientes conclusiones sobre la validez y las limitaciones del modelo de inferencia ecol\'ogica implementado.

\subsubsection{Ajuste de las Medias: Diagn\'ostico Diferencial}

Los resultados para las medias revelan un patr\'on heterog\'eneo que merece an\'alisis cuidadoso:

\begin{itemize}
    \item \textbf{PN --- Excelente ajuste} ($p = 0.648$): El modelo captura con alta precisi\'on la proporci\'on media de votos a PN (observado: 0.4696; predicho: 0.4698). El p-valor pr\'oximo a 0.5 indica que el valor observado se sit\'ua cerca del centro de la distribuci\'on posterior predictiva, lo cual constituye el resultado ideal.

    \item \textbf{FA --- Desajuste sistem\'atico detectable} ($p = 0.000$): El modelo subestima levemente la proporci\'on media de FA (0.4866 vs.\ 0.4880). Aunque la discrepancia absoluta es peque\~na (0.14 puntos porcentuales), es estad\'isticamente significativa debido al gran n\'umero de circuitos. Esto indica que las probabilidades de transferencia estimadas dirigen, en promedio, una fracci\'on marginalmente inferior de votos hacia FA respecto a lo observado.

    \item \textbf{Blancos --- Sobreestimaci\'on sistem\'atica} ($p = 1.000$): El modelo sobreestima la proporci\'on media de votos en blanco (0.0437 vs.\ 0.0424). Esta sobreestimaci\'on es una consecuencia directa del c\'alculo residual: la leve subestimaci\'on de FA se traduce mec\'anicamente en una sobreestimaci\'on de los votos en blanco.
\end{itemize}

Es importante destacar que un p-valor de 0.000 no debe interpretarse como indicador de ``buen ajuste'', sino como evidencia de que el modelo no logra reproducir el estad\'istico observado en ninguna de sus r\'eplicas. Sin embargo, la relevancia pr\'actica de esta discrepancia depende de su magnitud: en el caso de FA, la diferencia de 0.14 puntos porcentuales es lo suficientemente peque\~na como para no alterar las conclusiones sustantivas del an\'alisis.

\subsubsection{Subdispersi\'on: Limitaci\'on Estructural del Modelo}

El hallazgo m\'as relevante de los PPC es la \textbf{subdispersi\'on sistem\'atica}: el modelo genera r\'eplicas con menor variabilidad que los datos observados. Todos los p-valores de desviaci\'on est\'andar son 0.000, lo que confirma que \emph{ninguna} r\'eplica alcanza la dispersi\'on observada en los datos reales:

\begin{itemize}
    \item FA: dispersi\'on observada 0.1348 vs.\ predicha 0.1244 (subestimaci\'on del 7.7\%).
    \item PN: dispersi\'on observada 0.1373 vs.\ predicha 0.1246 (subestimaci\'on del 9.2\%).
    \item Blancos: dispersi\'on observada 0.0145 vs.\ predicha 0.0002 (subestimaci\'on del 98.6\%).
\end{itemize}

Esta subdispersi\'on es una limitaci\'on conocida de los modelos de inferencia ecol\'ogica con par\'ametros globales, y puede atribuirse a varios factores:

\begin{enumerate}
    \item \textbf{Heterogeneidad no modelada}: El modelo nacional asume probabilidades de transferencia $\beta_{ij}$ constantes en todos los circuitos. En la realidad, estas probabilidades var\'ian seg\'un caracter\'isticas geogr\'aficas, socioecon\'omicas y demogr\'aficas, como se documenta en el an\'alisis departamental (Secci\'on~\ref{sec:resultados_departamento}).

    \item \textbf{Pooling completo}: Al estimar par\'ametros globales sin jerarqu\'ias subnacionales, se produce un efecto de \textit{shrinkage} hacia la media nacional que reduce artificialmente la variabilidad predicha.

    \item \textbf{Factores locales}: Din\'amicas pol\'iticas espec\'ificas de ciertos circuitos (campa\~nas locales, candidatos con arraigo territorial, eventos puntuales) introducen variabilidad adicional que el modelo agregado no captura.
\end{enumerate}

El caso de los votos en blanco es particularmente ilustrativo: la proporci\'on de votos en blanco por circuito depende de factores altamente locales (facilidad de acceso, caracter\'isticas del electorado, condiciones de votaci\'on), y el modelo, al tratarla como residuo de proporciones globales, genera una predicci\'on pr\'acticamente constante que no captura esta heterogeneidad.

\subsubsection{Ajuste en Extremos de la Distribuci\'on}

Los p-valores de m\'inimos y m\'aximos ($p = 0.000$ en todos los casos) revelan una limitaci\'on adicional: el modelo no puede reproducir los valores extremos observados en los circuitos at\'ipicos. En particular:

\begin{itemize}
    \item Circuitos con proporci\'on nula de FA o Blancos (m\'inimo = 0.0000) y circuitos donde el 100\% de los votos se dirigen a PN (m\'aximo = 1.0000) reflejan la existencia de circuitos muy peque\~nos con comportamientos degenerados que el modelo, basado en probabilidades de transferencia suaves, no puede generar.

    \item Estos circuitos extremos son escasos en n\'umero pero tienen impacto en los estad\'isticos de m\'inimo y m\'aximo. Su presencia no invalida las estimaciones de transferencia promedio, que constituyen el objetivo central del an\'alisis.
\end{itemize}

\subsubsection{Validaci\'on General del Modelo}

El an\'alisis de posterior predictive checks revela tanto fortalezas como limitaciones del modelo implementado. Es necesario distinguir entre ambas con rigor:

\textbf{Fortalezas:}
\begin{enumerate}
    \item La proporci\'on media de votos a PN se reproduce con precisi\'on pr\'acticamente perfecta ($p = 0.648$), lo que valida la consistencia interna de las estimaciones de transferencia hacia el candidato oficialista.

    \item La subestimaci\'on de la media de FA, aunque estad\'isticamente significativa, es de magnitud reducida (0.14 puntos porcentuales) y no altera las conclusiones sustantivas sobre las transferencias de votos.

    \item Las distribuciones centrales de FA y PN son razonablemente bien reproducidas, como evidencian los Q-Q plots y los histogramas superpuestos.
\end{enumerate}

\textbf{Limitaciones:}
\begin{enumerate}
    \item La subdispersi\'on sistem\'atica (p-valores de desviaci\'on est\'andar = 0.000 para las tres variables) indica que el modelo asume mayor homogeneidad entre circuitos que la existente en la realidad. Este es el desajuste m\'as importante y afecta la precisi\'on de las estimaciones a nivel de circuito individual, aunque no las estimaciones agregadas.

    \item La sobreestimaci\'on de votos en blanco y la virtual ausencia de variabilidad predicha para esta categor\'ia representan una limitaci\'on inherente al c\'alculo residual empleado.

    \item El modelo no reproduce los valores extremos observados en circuitos at\'ipicos, lo que sugiere que los intervalos de credibilidad a nivel de circuito pueden ser demasiado estrechos.
\end{enumerate}

Estas limitaciones son consistentes con lo reportado en la literatura sobre modelos de inferencia ecol\'ogica con par\'ametros globales \citep{king1997} y no invalidan las conclusiones centrales del estudio, que se refieren a las transferencias \emph{promedio} de votos entre la primera vuelta y el ballotage.

\subsubsection{Recomendaciones para Mejoras Futuras}

Los resultados de los PPC sugieren dos v\'ias principales de mejora para reducir la subdispersi\'on y mejorar el ajuste global:

\begin{enumerate}
    \item \textbf{Modelo jer\'arquico con variabilidad subnacional}: Implementar un modelo multinivel donde las probabilidades de transferencia $\beta_{ij}$ var\'ien seg\'un caracter\'isticas del circuito:
    \begin{equation}
    \beta_{jk,c} \sim \text{Dirichlet}(\alpha_{jk,g(c)})
    \end{equation}
    donde $g(c)$ indica el grupo (departamento, regi\'on, clasificaci\'on urbano/rural) al que pertenece el circuito $c$. Esto permitir\'ia capturar mayor variabilidad entre circuitos y reducir la subdispersi\'on observada.

    \item \textbf{Modelado expl\'icito de votos en blanco}: En lugar de calcular los votos en blanco por diferencia, estimar directamente las tres probabilidades de transferencia ($\beta_{j,\text{FA}}$, $\beta_{j,\text{PN}}$, $\beta_{j,\text{Blancos}}$) bajo una distribuci\'on de Dirichlet con restricci\'on de suma unitaria. Esto permitir\'ia asignar variabilidad propia a los votos en blanco y evitar la acumulaci\'on de sesgos residuales.
\end{enumerate}

\subsubsection{Conclusi\'on}

Los posterior predictive checks revelan que el modelo de inferencia ecol\'ogica de King implementado con PyMC~v5 presenta un ajuste adecuado para la estimaci\'on de transferencias promedio de votos, con un excelente desempe\~no en la media de PN ($p = 0.648$) y discrepancias menores pero detectables en la media de FA ($p = 0.000$, sesgo de 0.14 pp). La limitaci\'on principal es la subdispersi\'on sistem\'atica: el modelo subestima la variabilidad entre circuitos, particularmente para los votos en blanco. Esta limitaci\'on, inherente a los modelos de inferencia ecol\'ogica con par\'ametros globales, afecta la precisi\'on de las estimaciones a nivel de circuito individual pero no compromete las inferencias agregadas sobre transferencias de votos, que constituyen el objetivo central del an\'alisis. El an\'alisis departamental presentado en secciones anteriores proporciona una v\'ia complementaria para capturar la heterogeneidad geogr\'afica que el modelo nacional no incorpora.

\newpage

% Covariables Demográficas del Censo 2023
\section{Covariables Demográficas del Censo 2023 (Análisis Exploratorio con Datos Sintéticos)}
\label{sec:census}

\textbf{Advertencia fundamental}: Las covariables demográficas utilizadas en esta sección son \textbf{datos sintéticos generados para fines demostrativos}, no datos oficiales del Censo 2023. Al momento de realizar este estudio, los microdatos censales oficiales no se encontraban disponibles. Por tanto, \textbf{todas las correlaciones y patrones reportados en esta sección carecen de valor empírico} y no deben interpretarse como evidencia de relaciones reales entre variables demográficas y comportamiento electoral. El propósito de esta sección es exclusivamente metodológico: ilustrar cómo se integrarían covariables censales en el análisis de inferencia ecológica una vez que se disponga de datos oficiales.

El comportamiento electoral no ocurre en el vacío, sino que se enmarca en contextos demográficos, socioeconómicos y culturales específicos. En la literatura se ha documentado que factores como el nivel educativo y la composición etaria pueden estar asociados con patrones de transferencia de votos \citep{Dalton2002, Franklin2004}. Para ilustrar cómo se complementaría el análisis de inferencia ecológica con covariables demográficas, se presenta este ejercicio exploratorio a nivel departamental.

\subsection{Datos Censales}
\label{subsec:census_data}

Se reitera que las cifras presentadas son datos sintéticos que reflejan patrones realistas basados en tendencias históricas y estimaciones publicadas por el Instituto Nacional de Estadística (INE), pero \textbf{no representan valores oficiales del Censo 2023}.

Las variables demográficas consideradas son:

\begin{itemize}
    \item \textbf{Educación terciaria (\%)}: Porcentaje de población mayor de 25 años con educación terciaria completa (universitaria o terciaria no universitaria).
    \item \textbf{Edad mediana (años)}: Edad mediana de la población departamental, indicador de la estructura etaria.
\end{itemize}

El Cuadro~\ref{tab:census_data} presenta los valores sintéticos para los 19 departamentos del país.

\begin{table}[H]
\centering
\caption{Covariables demográficas por departamento (datos sintéticos --- no oficiales)}
\label{tab:census_data}
\begin{tabular}{lrr}
\toprule
\textbf{Departamento} & \textbf{Educación terciaria (\%)} & \textbf{Edad mediana (años)} \\
\midrule
Montevideo           & 35.2 & 36.8 \\
Canelones            & 21.4 & 34.5 \\
Maldonado            & 23.8 & 37.2 \\
Colonia              & 19.6 & 38.5 \\
Salto                & 17.8 & 35.9 \\
Paysandú             & 18.9 & 37.1 \\
Artigas              & 16.4 & 35.2 \\
Rivera               & 17.2 & 36.3 \\
Tacuarembó           & 16.8 & 36.7 \\
Soriano              & 18.3 & 38.9 \\
San José             & 20.1 & 35.6 \\
Rocha                & 17.5 & 37.8 \\
Río Negro            & 17.9 & 38.2 \\
Flores               & 16.7 & 39.8 \\
Florida              & 18.7 & 37.4 \\
Lavalleja            & 17.1 & 38.1 \\
Cerro Largo          & 15.3 & 36.9 \\
Treinta y Tres       & 16.1 & 37.5 \\
Durazno              & 15.7 & 38.6 \\
\bottomrule
\end{tabular}
\end{table}

Los valores sintéticos reproducen la estructura centro-periferia del Uruguay: mayor concentración de población con educación terciaria en la zona metropolitana (Montevideo, Canelones) y la costa (Maldonado), frente a valores menores en departamentos del interior rural.

\subsection{Correlaciones con Comportamiento Electoral}
\label{subsec:census_correlations}

Se calcularon coeficientes de correlación de Pearson entre las variables censales sintéticas y las estimaciones puntuales de las tasas de transferencia de votos obtenidas mediante inferencia ecológica. El Cuadro~\ref{tab:census_correlations} presenta la matriz de correlaciones. Se enfatiza que estas correlaciones se basan en datos sintéticos y, por tanto, \textbf{no constituyen hallazgos empíricos}.

\begin{table}
\caption{Correlaciones entre variables demográficas (Censo 2023) y comportamiento electoral}
\label{tab:census_correlations}
\begin{tabular}{llccc}
\toprule
Variable Demográfica & Variable Electoral & Correlación & Valor $p$ & Significativo \\
\midrule
edad_mediana & ca_defection_to_fa & -0.159 & 0.5162 & No \\
edad_mediana & ca_retention_by_pn & 0.097 & 0.6924 & No \\
edad_mediana & pc_retention_by_pn & 0.414 & 0.0780 & No \\
edad_mediana & pn_retention & -0.189 & 0.4374 & No \\
pct_educacion_terciaria & ca_defection_to_fa & 0.379 & 0.1090 & No \\
pct_educacion_terciaria & ca_retention_by_pn & -0.419 & 0.0740 & No \\
pct_educacion_terciaria & pc_retention_by_pn & 0.298 & 0.2145 & No \\
pct_educacion_terciaria & pn_retention & 0.404 & 0.0859 & No \\
poblacion & ca_defection_to_fa & 0.144 & 0.5564 & No \\
poblacion & ca_retention_by_pn & -0.178 & 0.4657 & No \\
poblacion & pc_retention_by_pn & 0.255 & 0.2919 & No \\
poblacion & pn_retention & 0.327 & 0.1718 & No \\
\bottomrule
\end{tabular}
\end{table}


A continuación se describen los patrones observados, con la salvedad de que su interpretación es puramente ilustrativa:

\paragraph{Nivel educativo y defección de Cabildo Abierto}
Se observa una correlación positiva moderada entre el porcentaje de educación terciaria (sintético) y la tasa de defección de CA hacia FA ($r = +0.380$, $p = 0.109$). La Figura~\ref{fig:census_edu} ilustra esta relación. Este patrón, de confirmarse con datos reales, sería consistente con la hipótesis de que departamentos con mayor nivel educativo presentan mayor volatilidad electoral. Sin embargo, dado el carácter sintético de los datos, esta correlación no puede considerarse evidencia de dicha relación.

\begin{figure}[H]
\centering
\includegraphics[width=0.85\textwidth]{../../outputs/figures/scatter_ca_defection_vs_education.pdf}
\caption{Defección CA vs nivel educativo (sintético) por departamento. Correlación basada en datos sintéticos ($r = 0.380$); no constituye evidencia empírica. Los puntos representan los 19 departamentos, con tamaño proporcional al electorado de CA en primera vuelta.}
\label{fig:census_edu}
\end{figure}

\paragraph{Nivel educativo y retención de CA por PN}
La correlación entre educación terciaria y retención de CA por el Partido Nacional es negativa ($r = -0.419$, $p = 0.074$). Este valor complementa la correlación anterior, pero su interpretación sustantiva queda condicionada a la disponibilidad de datos censales reales.

\paragraph{Edad mediana}
Las correlaciones con la edad mediana son débiles y no significativas ($r = -0.162$, $p = 0.509$ para defección de CA; $r = +0.193$, $p = 0.429$ para retención por PN).

\paragraph{Otras transferencias}
Las transferencias de PC a PN y la retención del PN no muestran correlaciones sustanciales con las variables demográficas sintéticas ($r = +0.043$ y $r = -0.098$ con educación, respectivamente).

La Figura~\ref{fig:census_heatmap} presenta visualmente la matriz de correlaciones mediante un mapa de calor.

\begin{figure}[H]
\centering
\includegraphics[width=0.8\textwidth]{../../outputs/figures/heatmap_census_correlations.pdf}
\caption{Matriz de correlaciones: variables demográficas (sintéticas) y electorales. Los tonos rojos indican correlaciones positivas, azules negativas. Se recuerda que estas correlaciones se basan en datos no oficiales.}
\label{fig:census_heatmap}
\end{figure}

\subsection{Limitaciones y Trabajo Futuro}
\label{subsec:census_interpretation}

Las limitaciones de este ejercicio exploratorio son sustanciales y deben considerarse antes de cualquier interpretación:

\paragraph{Datos sintéticos (limitación crítica)}
La limitación más importante es que todas las covariables demográficas son datos sintéticos. Las correlaciones reportadas en esta sección \textbf{no tienen valor empírico alguno}. Cualquier patrón observado podría ser un artefacto de la generación sintética de datos y no reflejar relaciones reales entre demografía y comportamiento electoral.

\paragraph{Tamaño muestral}
Con solo $N = 19$ observaciones (departamentos), el poder estadístico es limitado. Un análisis a nivel de municipios ($N \approx 125$) o circuitos ($N > 7000$) con covariables correspondientes proporcionaría mayor poder estadístico.

\paragraph{Falacia ecológica}
Aun con datos oficiales, las correlaciones a nivel agregado (departamental) no implican necesariamente relaciones a nivel individual. Esto es una manifestación de la falacia ecológica, el mismo problema que la inferencia ecológica busca mitigar. Para establecer relaciones individuales se requerirían encuestas post-electorales con datos de votación autorreportada.

\paragraph{Causalidad}
Las correlaciones no implican causalidad. Múltiples factores no observados (contexto político local, campañas departamentales, liderazgos locales) pueden confundir las relaciones entre variables demográficas y comportamiento electoral.

\paragraph{Trabajo futuro}
Una vez que el Instituto Nacional de Estadística (INE) publique los microdatos del Censo 2023, se recomienda replicar este análisis con las siguientes mejoras:

\begin{itemize}
    \item Variables adicionales: porcentaje de población urbana, tasa de desempleo, nivel socioeconómico (NSE), acceso a servicios.
    \item Análisis multinivel: modelar la variación electoral a nivel de circuito o sección con covariables censales correspondientes.
    \item Modelos jerárquicos: incorporar covariables en el modelo de inferencia ecológica como predictores de las tasas de transferencia (modelo de King con regresión sobre $\beta_{ij}$).
\end{itemize}

En resumen, esta sección ilustra la metodología para integrar covariables demográficas en el análisis de inferencia ecológica. Los resultados aquí presentados son exclusivamente demostrativos y deberán ser reemplazados por análisis con datos oficiales cuando estos se encuentren disponibles.


\end{document}

% ================================================================================
% FIN DEL DOCUMENTO
% ================================================================================
